\documentclass[11pt]{article}

% Change "review" to "final" to generate the final (sometimes called camera-ready) version.
% Change to "preprint" to generate a non-anonymous version with page numbers.
\usepackage[review]{acl}

% Standard package includes
\usepackage{times}
\usepackage{latexsym}
    
% For proper rendering and hyphenation of words containing Latin characters (including in bib files)
\usepackage[T1]{fontenc}
% For Vietnamese characters
% \usepackage[T5]{fontenc}
% See https://www.latex-project.org/help/documentation/encguide.pdf for other character sets

% This assumes your files are encoded as UTF8
\usepackage[utf8]{inputenc}

% This is not strictly necessary, and may be commented out,
% but it will improve the layout of the manuscript,
% and will typically save some space.
\usepackage{microtype}

% This is also not strictly necessary, and may be commented out.
% However, it will improve the aesthetics of text in
% the typewriter font.
\usepackage{inconsolata}

\usepackage[utf8]{inputenc} % allow utf-8 input
\usepackage[T1]{fontenc}    % use 8-bit T1 fonts
\usepackage{hyperref}       % hyperlinks
\usepackage{url}            % simple URL typesetting
\usepackage{booktabs}       % professional-quality tables
\usepackage{amsfonts}       % blackboard math symbols
\usepackage{nicefrac}       % compact symbols for 1/2, etc.
\usepackage{microtype}      % microtypography
\usepackage{xcolor}         % colors
\usepackage{amsmath}     % argmax, align, equation environments
\usepackage{amssymb}     % mathcal, mathbb
\usepackage[table]{xcolor}
\usepackage{tabularx}

\definecolor{echoshade}{HTML}{FFF2CC}
\definecolor{grposhade}{HTML}{E8F1FF}   % soft blue for GRPO
\definecolor{oracleshade}{HTML}{EAF7EA} % soft green for greedy oracle
\usepackage{bm}
\newcommand{\ms}[2]{#1{\scriptsize $_{\pm #2}$}}
% \newcommand{\ms}[2]{\ensuremath{#1_{\scriptscriptstyle \pm #2}}}
% \newcommand{\bms}[2]{\ensuremath{\bm{#1_{\scriptscriptstyle \pm #2}}}}
% expectation operator
\newcommand{\EE}{\mathbb{E}}
\usepackage{bm}          % bold math if needed
\usepackage{bbm}         % \mathbbm{1} for indicator function
\usepackage{cleveref}
 \usepackage{wrapfig}
\usepackage{graphicx}    % \resizebox
\usepackage{subcaption}
\usepackage{multirow}    % optional but good to have for tables
\usepackage{array}       % better column formatting
\definecolor{darkgold}{RGB}{184, 134, 11} % dark goldenrod
\newcommand{\an}[1]{\textcolor{black}{#1}}
\newcommand{\nk}[1]{\textcolor{teal}{#1}}
\newcommand{\todo}[1]{\textcolor{red}{TODO: #1}}
\DeclareMathOperator*{\argmax}{arg\,max}
\DeclareMathOperator*{\argmin}{arg\,min}
% in preamble
\newcommand{\std}[1]{{\scriptsize$\pm$#1}}

% Theorem/proof environments
\usepackage{amsthm}
\usepackage{arydshln}
\theoremstyle{plain}
\newtheorem{theorem}{Theorem}[section]
\newtheorem{proposition}[theorem]{Proposition}
\newtheorem{lemma}[theorem]{Lemma}
\newtheorem{corollary}[theorem]{Corollary}

\theoremstyle{definition}
\newtheorem{definition}[theorem]{Definition}
\newtheorem{assumption}[theorem]{Assumption}
\newtheorem{example}[theorem]{Example}

\theoremstyle{remark}
\newtheorem{remark}[theorem]{Remark}
\usepackage[most]{tcolorbox}
\usepackage{pdflscape}
\newtcolorbox{theorybox}{
    enhanced,
    breakable=false,
      colback=white,
    colframe=blue!50!black,
    boxrule=0pt,
    leftrule=3pt,
    arc=1.5mm,
    left=1.2mm,
    right=1.2mm,
    top=1mm,
    bottom=1mm,
    before skip=0.8em,
    after skip=0.8em
}
\usepackage[edges]{forest} \usepackage{xcolor} \definecolor{echoshade}{HTML}{FFF2CC} \definecolor{lightblue}{HTML}{EAF3FF} \definecolor{lightgreen}{HTML}{EAF7EA} \definecolor{lightgraybox}{HTML}{F3F4F6}

\usepackage{tcolorbox}
\usepackage{listings}
 
\tcbuselibrary{skins, breakable}
 
% ── Colour palette ───────────────────────────────────────────────────────────
\definecolor{promptgray}{RGB}{245,245,245}
\definecolor{promptborder}{RGB}{180,180,180}
\definecolor{headergray}{RGB}{210,210,210}
\definecolor{keywordblue}{RGB}{0,70,150}
 
% ── Reusable tcolorbox style ─────────────────────────────────────────────────
\tcbuselibrary{skins, breakable}
\newtcolorbox{promptbox}[2][]{%
  enhanced, breakable,
  colback=promptgray, colframe=promptborder,
  coltitle=black, fonttitle=\small\bfseries\sffamily,
  title={#2},
  attach boxed title to top left={yshift=-2mm, xshift=4mm},
  boxed title style={colback=headergray, colframe=promptborder, sharp corners},
  sharp corners,
  top=4mm, left=4mm, right=4mm, bottom=3mm,
  #1
}

%Including images in your LaTeX document requires adding
%additional package(s)
\usepackage{graphicx}
\usepackage{xcolor}


% If the title and author information does not fit in the area allocated, uncomment the following
%
%\setlength\titlebox{<dim>}
%
% and set <dim> to something 5cm or larger.

\title{\textsc{ECHO}: Learning Epistemically Adaptive Language Agents with Turn-Level Credit
%Teaching LLM Agents to Seek Information Without Reasoning Traces
}

% Author information can be set in various styles:
% For several authors from the same institution:
% \author{Author 1 \and ... \and Author n \\
%         Address line \\ ... \\ Address line}
% if the names do not fit well on one line use
%         Author 1 \\ {\bf Author 2} \\ ... \\ {\bf Author n} \\
% For authors from different institutions:
% \author{Author 1 \\ Address line \\  ... \\ Address line
%         \And  ... \And
%         Author n \\ Address line \\ ... \\ Address line}
% To start a separate ``row'' of authors use \AND, as in
% \author{Author 1 \\ Address line \\  ... \\ Address line
%         \AND
%         Author 2 \\ Address line \\ ... \\ Address line \And
%         Author 3 \\ Address line \\ ... \\ Address line}

\author{First Author \\
  Affiliation / Address line 1 \\
  Affiliation / Address line 2 \\
  Affiliation / Address line 3 \\
  \texttt{email@domain} \\\And
  Second Author \\
  Affiliation / Address line 1 \\
  Affiliation / Address line 2 \\
  Affiliation / Address line 3 \\
  \texttt{email@domain} \\}

%\author{
%  \textbf{First Author\textsuperscript{1}},
%  \textbf{Second Author\textsuperscript{1,2}},
%  \textbf{Third T. Author\textsuperscript{1}},
%  \textbf{Fourth Author\textsuperscript{1}},
%\\
%  \textbf{Fifth Author\textsuperscript{1,2}},
%  \textbf{Sixth Author\textsuperscript{1}},
%  \textbf{Seventh Author\textsuperscript{1}},
%  \textbf{Eighth Author \textsuperscript{1,2,3,4}},
%\\
%  \textbf{Ninth Author\textsuperscript{1}},
%  \textbf{Tenth Author\textsuperscript{1}},
%  \textbf{Eleventh E. Author\textsuperscript{1,2,3,4,5}},
%  \textbf{Twelfth Author\textsuperscript{1}},
%\\
%  \textbf{Thirteenth Author\textsuperscript{3}},
%  \textbf{Fourteenth F. Author\textsuperscript{2,4}},
%  \textbf{Fifteenth Author\textsuperscript{1}},
%  \textbf{Sixteenth Author\textsuperscript{1}},
%\\
%  \textbf{Seventeenth S. Author\textsuperscript{4,5}},
%  \textbf{Eighteenth Author\textsuperscript{3,4}},
%  \textbf{Nineteenth N. Author\textsuperscript{2,5}},
%  \textbf{Twentieth Author\textsuperscript{1}}
%\\
%\\
%  \textsuperscript{1}Affiliation 1,
%  \textsuperscript{2}Affiliation 2,
%  \textsuperscript{3}Affiliation 3,
%  \textsuperscript{4}Affiliation 4,
%  \textsuperscript{5}Affiliation 5
%\\
%  \small{
%    \textbf{Correspondence:} \href{mailto:email@domain}{email@domain}
%  }
%}

\begin{document}
\maketitle
\begin{abstract}
What does it mean for a language agent to be adaptive? Effective multi-turn agents must decide what information to seek, how to use new evidence, and when they are certain enough to act.
We introduce \textbf{Epistemic Decision Processes} (EDPs), a belief-state formulation of multi-turn information seeking in which actions produce external observations that update the agent's posterior over a latent task variable. EDPs make epistemic adaptivity explicit: good policies choose actions that are useful under the current belief, not merely those that correlate with eventual success. We prove that belief-agnostic policies can suffer errors that compound exponentially over the horizon, and that aggregate trajectory returns can fail to identify the per-turn Bayesian advantage needed for epistemic credit. We then introduce \textsc{ECHO} (\textbf{E}pistemic \textbf{C}redit for \textbf{H}istory-Conditioned \textbf{O}ptimization), a practical clipped policy-gradient objective that assigns turn-level credit using posterior-sensitive rewards. In the \textbf{Clue Selector Game}, a novel controlled evidence-seeking benchmark, \nk{we show that} \textsc{ECHO} substantially improves resolution, information gain, and efficiency over trajectory-level GRPO, and matches or exceeds frontier baselines on epistemic metrics such as grounding, recovery, and calibration while producing almost no visible reasoning text.
\todo{Add CRAFT mention}
\end{abstract}

\section{Introduction}
\vspace*{-2mm}

% \an{epistemic adaptivity is important, but is an unsolved problem from two angles: 1) most research focus on final task accuracy, which limits credit assignment in sparse-reward settings/long horizon and 2) this and additional cost of deploying multiturn environments for RL research is another problem.  }
% Information seeking is a sequential behavior. In diagnosis, troubleshooting, scientific investigation, and search, the central challenge is rarely to act once with perfect information; it is to decide what evidence to gather next. A useful question is not useful in isolation. Its value depends on what is already known, what uncertainty remains, and whether the answer will change the agent’s belief enough to guide future action. This makes “adaptivity” a fundamentally epistemic property: good agents must track what information has already been acquired and choose actions that reduce the uncertainty that still matters.

%Information seeking is inherently sequential. 
\nk{In multi-step problem solving, including medical diagnoses, technical troubleshooting, scientific inquiry, and search, 
%In diagnosis, troubleshooting, scientific inquiry, and search, 
effective agents must decide what evidence to gather next; they do not merely act upon incomplete information~\citep{howard1966information,pauker1980threshold,pirolli1999information,settles2009active}.} \an{This skill is central to human learning and question asking~\citep{chouinard2007children,geva2021didaristotleuselaptop,meder2019stepwise,rothe2018people}, and increasingly so to LLM agents that search the web~\citep{deng2023mind2web,yao2022webshop,Zhou2023WebArenaAR}, call tools~\citep{schick2023toolformer}, inspect documents~\citep{zhang-etal-2025-belle}, and revise actions as observations arrive~\citep{ferrag2025llm,yao2022react}.} In such settings, the value of an \nk{information-seeking} question depends on the current belief state: what is already known, what uncertainty remains, and whether the answer would change future action~\citep{kaelbling1998planning}. Thus, adaptivity is fundamentally \textit{epistemic}~\citep{meder2019stepwise,rothe2018people}---good agents choose actions that reduce \nk{not just global uncertainty, but the uncertainty} that \nk{matters most in the current context}.
% as a user's intent, a preference, a diagnosis, or a target entity

\begin{figure*}
    \centering
    \begin{subfigure}[t]{0.58\textwidth}
        \centering
        \includegraphics[width=\linewidth]{figures/intro_metrics.png}
    \end{subfigure}
    \hfill
    \begin{subfigure}[t]{0.35\textwidth}
        \centering
        \includegraphics[width=\linewidth]{figures/sample_efficiency.png}
    \end{subfigure}

    \vspace*{-3mm}
  \caption{
Results on multi-turn information seeking in the \textbf{Clue Selector Game (CSG)}. \textbf{A 1.5B \textsc{ECHO}-trained Qwen 2.5 policy reaches frontier-level task resolution, by asking questions that more reliably update the candidate belief state.} Compared with GRPO and prompting baselines, \textsc{ECHO} produces fewer zero-elimination turns and learns substantially \textit{faster}, illustrating the central idea of epistemic credit: useful actions should be rewarded when they reduce uncertainty under the current belief state, not only when the final trajectory succeeds.
}
    \label{fig:intro_teaser}
    \vspace*{-3mm}
\end{figure*}

\nk{In multi-turn tasks, final success may only be known at the end of the trajectory. Intermediate rewards are often unavailable, learned approximately through process reward models~\citep{choudhury2025processrewardmodelsllm,hu2026playing20questiongame,lee2025prints}, or disconnected from whether a query actually changed the agent's belief (e.g., when rewarding only task outcomes). The agent's uncertainty is often organized around a hidden variable of interest, such as a goal or sub-goal, that must be inferred through interaction rather than revealed directly~\citep{abdulhai2023lmrl,mazzaccara-etal-2024-learning,zhang-etal-2024-probing,zhang2024the}. While the agent may observe the full history of its own queries and the resulting feedback, the value of the next action still depends on how that history changes its belief. This creates a \textit{credit-assignment} problem that existing methods do not fully address~\citep{bertolazzi2023chatgpt,laban2026llms, patil2025berkeley}.}


\an{We argue that information-seeking agents need more than temporal credit assignment \nk{for optimal behavior}: they need \textit{epistemic} credit assignment, which \nk{reflects the usefulness of evidence elicited by an action given the agent's current belief, not merely that it appears in a trajectory that eventually succeeds.}
%\an{Our core insight is that in information-seeking tasks, an action should be credited for being informative under the agent's current belief, not merely for appearing in a trajectory that eventually succeeds.} 
We formalize this through \emph{Epistemic Decision Processes} (EDPs), a belief-state view\footnote{
EDPs are a task-level specialization of belief-state decision processes: unlike BAMDPs~\citep{duff2002optimal,martin1967bayesian}, which typically model uncertainty over transition/reward parameters or latent MDP identities, EDPs focus on uncertainty over a hidden task variable under a fixed interaction interface; see Sec.~\ref{sec:edp_theory} and Appendix~\ref{app:proofs} for details.
} of multi-turn information seeking in which actions produce external observations that update a posterior over a latent task variable. Our theory shows why this distinction matters: there exist long-horizon information-seeking tasks where belief-sensitive policies are necessary for strong performance, while trajectory-level returns can be too coarse to identify which turns actually reduced uncertainty.
}
% This connects to Bayes-adaptive reflective reasoning methods such as BARL~\citep{zhang2025markovianreflectiveexplorationbayesadaptive}, but our focus is external observation-driven interaction, where the policy must learn from its own queries rather than only revise an internal reasoning trace.

% This motivates \textsc{ECHO}, a clipped policy-gradient objective that assigns turn-level credit using posterior-sensitive rewards.}

 
\nk{We then} introduce \textsc{ECHO}, a clipped policy-gradient objective that assigns turn-level credit using posterior-sensitive rewards. \nk{\textsc{ECHO} trains} agents to ask questions that are useful under the current belief state, not merely questions that correlate with eventual success. We instantiate this setting in the \emph{Clue Selector Game} (CSG), \nk{a deduction task} where 
%the task variable remains hidden from the policy, but 
the hypothesis space, candidate posterior, and belief updates are exactly measurable. Inspired by prior work in step-level epistemic competence~\citep{shao2025llm}, CSG evaluates both final resolution and process-level epistemic behavior, using diagnostics 
%inspired by epistemic-competence evaluations 
such as groundedness, zero-information turns, recovery after low-information actions, and late-stage calibration. \textbf{Empirically, \textsc{ECHO} trains a 1.5B policy to reach frontier-level task resolution, reduce zero-elimination turns, and learn substantially faster than standard RL methods, while producing almost no visible reasoning text.} This reveals a form of ``silent exploration'': adaptive information seeking expressed through belief-conditioned actions rather than long chains of thought.

Our novel contributions are as follows:
\begin{enumerate}
    \item We introduce \emph{Epistemic Decision Processes} (EDPs), a belief-state formulation of multi-turn information-seeking language agents, and define epistemic adaptivity as posterior-sensitive action selection.
    \item We prove that belief-sensitive policies are necessary for strong performance in EDPs, with an exponential gap over belief-agnostic policies, and that aggregate trajectory return is generally insufficient for recovering per-turn Bayesian advantage.

    \item We introduce \textsc{ECHO}, a clipped policy-gradient objective for epistemic credit assignment over posterior-sensitive turn rewards, and evaluate it in the Clue Selector Game, a controlled EDP benchmark with explicit beliefs and measurable posterior contraction. \textsc{ECHO} outperforms trajectory-level GRPO and prompting baselines in resolution, information gain, and sample efficiency, while showing strong grounding, recovery, and calibration with little visible reasoning.

    % \item We propose \textsc{ECHO}, a practical clipped policy-gradient objective that applies group-normalized credit to posterior-sensitive turn rewards rather than full episode returns.
    % \item We introduce and evaluate on the Clue Selector Game, a controlled EDP benchmark where the belief state is explicit and reward measures posterior contraction.
    % \item We show that \textsc{ECHO} substantially outperforms trajectory-level GRPO and prompting baselines on resolution, information gain, and efficiency, while exhibiting strong process-level epistemic quality---grounding, recovery, and calibration---with little visible reasoning.
\end{enumerate}

\iffalse
We argue that information-seeking agents require a further distinction: temporal credit assignment is not enough; agents also need \textit{epistemic} credit assignment. We formalize this view through Epistemic Decision Processes\footnote{EDPs can be viewed as a belief-state specialization of POMDPs~\citep{kaelbling1998planning} for multi-turn information-seeking agents. We use the term to emphasize the epistemic structure of such tasks: actions are valuable not only because they lead to terminal success, but because they change the agent's posterior over a latent task variable.}, a belief-state formulation of multi-turn information seeking in which actions produce external observations that update a posterior over a latent task variable. This connects to recent Bayes-adaptive reflective reasoning methods such as BARL~\citep{zhang2025markovianreflectiveexplorationbayesadaptive}, which uses RL to encourage reflective exploration in single-turn LLM reasoning, but our focus is on external observation-driven environments where the agent must learn from its own queries. We then introduce ECHO, a practical clipped policy-gradient objective that assigns turn-level credit using posterior-sensitive rewards. \nk{The goal is to train agents that ask questions that are truly useful under the current belief state, not ones that merely correlate with eventual success.}
\fi


% \an{Importantly, this form of adaptivity that is crucial, but often limited, in multiturn interactions involving LLMs~\citep{bertolazzi2023chatgpt, patil2025berkeley, laban2026llms} is the ability to consistently display a robust balance between exploring the environment, while also exploiting the policy that results in strong task performance. But this quality is often challenging for current LLMs due to the problem of credit assignment, where the long trajectories and sparse-reward settings make it difficult for the agent to attribute which of its actions led to the observed outcome, even though multiple Reinforcement Learning (RL) frameworks have been proposed to tackle this challenging problem. While such credit assignment is crucial for robust learning, current multi-turn turn RL methods designed for LLMs majorly focus on “temporal” credit assignment—which is the standard method~\citep{sutton1984temporal, barto2003recent, schulman2017proximal, shao2024deepseekmath} that fails to explicitly account for the epistemic nature of the problem. Specifically, information-seeking acts in such multi-step settings require tracking uncertainty, not just about the environment where the agent operates, but also its own belief-states as new evidence comes in. }


\iffalse
\todo{Temporal credit assignment angle might fit here: temporal credit assignment addresses how to credit intermediate turns given sparse observables but does not directly address how actions affect agent beliefs}



% Large language models (LLMs) are increasingly deployed in exactly these interactive settings~\an{citations}. Modern language agents search the web, call tools, query databases, inspect documents, and interact with stochastic environments over multiple turns. 

\nk{Large language model (LLM)} agents now routinely operate in interactive loops: they search the web~\citep{deng2023mind2web,yao2022webshop,Zhou2023WebArenaAR}, call tools~\citep{schick2023toolformer}, inspect documents~\citep{zhang-etal-2025-belle}, schedule meetings~\citep{zou2026calbenchevaluatingcoordinationprivacytradeoffs}, and revise their actions as new observations arrive~\citep{yao2022react,ferrag2025llm}. In these settings, the agent’s central problem is not only to produce a correct final answer~\citep{fan2025ssrl, jinsearch,shao2025llm}, but to decide which observation would be most useful next and act accordingly~\citep{bae2025alignsearchbeliefguidedexploratory, hua2026learningexplorescalingagentic}. \todo{$\leftarrow$ this is fundamentally an {\it epistemic} task, not simply one of action selection}



This shift has motivated recent work on multi-turn reinforcement learning for LLM agents, much of which focuses on temporal credit assignment~\citep{sutton1984temporal}: how to assign credit to intermediate turns when only delayed trajectory-level outcomes are observed~\citep{shani2024multiturnreinforcementlearningpreference, wei2025reinforcingmultiturnreasoningllm}. This is an important problem, and recent methods explicitly target fine-grained turn-level advantage estimation for long-horizon LLM interactions.

\an{\textbf{core point: EDP must be highlighted in intro, but the CSG environment contribution must also be part of intro problem setting}: long-horizon task environments are typically costly for RL training, which limits research in novel ways of training information seeking agents~\citep{abdulhai2023lmrl}; CSG provides this contribution many works do not truly propose environments where language-based RL agents can be trained and evaluated in reasonably multiturn settings~\citep{singh2026multilevel, chen2026learning}}
\an{\citep{nideployable}: another work in adaptive policies}

We argue that information-seeking agents require a further distinction: temporal credit assignment is not enough; agents also need \textit{epistemic} credit assignment. In an information-seeking task, the value of an action is determined by how it changes the agent’s \textit{belief} at the time it is chosen~\citep{kausik2026context}. A question that reduces the remaining uncertainty early in a trajectory may be valuable, while the same question later may be redundant or completely uninformative. Thus, optimizing only aggregate trajectory return can obscure the per-turn belief updates that make an agent adaptive.
\fi

 

% \paragraph{Epistemic Decision Process (EDP).}
% We study a finite-horizon interaction in which an agent must act under uncertainty about
% a latent task variable $z^\star \sim \rho_0$. At turn $t$, the agent observes a history
% \[
% h_t = (a_1,o_1,\ldots,a_{t-1},o_{t-1}),
% \]
% selects an action $a_t \in \mathcal{A}$, and receives evidence
% $o_t \sim O(\cdot \mid z^\star,h_t,a_t)$. The history induces a belief
% \[
% b_t(z) = p(z^\star=z \mid h_t),
% \]
% which summarizes the agent's epistemic state. The goal is to select actions that both
% make progress toward the task objective and reduce uncertainty over $z^\star$.



% Previous work~\citep{choudhury2025better} has explored in the form of “privileged” information that allows the models to learn from an expert teacher using imitation learning at training time, but such privileged information is not realistic to assume in real-world situations where such a teacher or expert may not be available; and one may not have access to “expert-trajectories” when frontier models themselves may not perform well in the benchmark (our empirical results from the Clue Game benchmark). 
% \an{Adaptive agency~\citep{shao2025llm, zhang2025beyond} in LLM agents for long-horizon tasks is not necessarily a property of what the model says in its reasoning trace; it can be a property of how the policy updates its actions under uncertainty.}
 


\vspace*{-3mm}
\section{Related Work}

\vspace*{-3mm}
\paragraph{Temporal credit assignment and in-context learning}
Standard temporal credit assignment~\citep{barto2003recent,schulman2017proximalpolicyoptimizationalgorithms,shao2024deepseekmath,sutton1984temporal} can propagate delayed rewards to earlier turns in a task, but unless the state or reward explicitly represents belief change, it does not identify whether an action reduced uncertainty under the belief state at that moment.
%One might hope that this issue disappears 
LLM agents' full query--observation history can be placed \textit{in-context}~\citep{brown2020language,olsson2022context}; the model could, in principle, update its belief in-context and ask the next informative question. However, recent work suggests that in-context updating is sensitive to attention biases, context-length, and implicit priors shaped by pretraining~\citep{kossen2024incontext,liu-etal-2024-lost, falck2024context}. Bayesian methods such as BED-LLM compute informative queries at inference time using expected information gain over the variable of interest~\citep{choudhury2025bed,echarghaoui2026balar,hu2026playing20questiongame,mazzaccara-etal-2024-learning,zhang-etal-2024-probing}. However, these methods also highlight why belief-sensitive learning is difficult in realistic language tasks: the hypothesis space's pretraining and implicit priors make the agent's posterior difficult to inspect or credit at each step. Moreover, inference-time strategies~\citep{jeong2025reflectthenplan,yao2022react} must recompute query utility at every turn; \nk{they do not bake in adaptive querying to the policy.}

 

\vspace*{-3mm}
\paragraph{Bayesian learning and belief-conditioned policies}
\an{Epistemic adaptivity builds on belief-conditioned decision making: agents represent uncertainty, update it from evidence, and act to make future decisions better informed. Bayes-Adaptive MDPs formalize this by augmenting state with a posterior over environment parameters or latent MDP identities~\citep{bellman1956problem,duff2002optimal,Ghavamzadeh_2015}, with deep RL extensions using approximate inference~\citep{Zintgraf2020VariBAD}. Language-based information seeking \textit{also} requires crediting actions for changing belief; otherwise, outcome-based RL can lock agents into low-information interaction patterns~\citep{zou2026informationselflockingreinforcementlearning}. 
%One might expect LLM agents to update beliefs “in-context”~\citep{brown2020language, olsson2022context} from the full query--observation history, but such updating is unreliable under attention biases, context-usage effects, and pretraining-shaped priors~\citep{kossen2024incontext,liu-etal-2024-lost}. 
%Inference-time methods instead select queries by expected information gain or mutual information over a variable of interest~\citep{choudhury2025bed,echarghaoui2026balar,hu2026playing20questiongame,mazzaccara-etal-2024-learning,zhang-etal-2024-probing}, but they must recompute query utility at each turn and often rely on hypothesis spaces whose posterior is difficult to inspect. 
Training-time Bayes-adaptive methods such as BARL encourage reflective exploration within generated reasoning traces~\citep{zhang2025markovianreflectiveexplorationbayesadaptive}, but recent work cautions that visible reasoning traces are not necessarily faithful evidence of reasoning competence~\citep{palod2025performative,samineni2025localcoherenceglobalvalidity}. \textsc{ECHO} via EDPs instead studies ``silent exploration'' across external query--observation turns in a multiturn setup, where adaptivity is expressed through belief-informative actions rather than verbalized rationales. }

% EDPs differ in scope and object of uncertainty: they model external multi-turn interactions in which the hidden quantity is a task variable under a fixed interface, making belief change over that variable the direct target of credit assignment.}

% \textsc{ECHO} shifts the locus of adaptivity from verbal reflection to external query--observation turns: the policy explores by choosing belief-informative actions, not by producing longer rationales.
 
 
\vspace*{-3mm}
\paragraph{Controlled environments for information-seeking agents}
 
%\an{Controllable environments are needed to evaluate belief-conditioned behavior in multi-turn language agents~\citep{abdulhai2023lmrl,shojaee2026illusion,yue2026does}. 
\an{Common question-asking settings---including entity deduction, 20-questions-style tasks, active task disambiguation, and preference elicitation---study how agents infer hidden targets through interaction~\citep{andukuri2024stargateteachinglanguagemodels,kobalczyk2025activetaskdisambiguationllms,lialfa,zhang-etal-2024-probing,zhou2024archer}. However, non-uniform LLM priors affect question selection and evaluation~\citep{mazzaccara-etal-2024-learning}; intermediate progress is usually tracked via hand-crafted knowledge-bases~\citep{zhao2016towards} or approximated through learned reward networks, process reward models, or inference-time information-gain estimates~\citep{choudhury2025bed,choudhury2025processrewardmodelsllm,hu2026playing20questiongame,lee2025prints}. Final accuracy can miss process-level failures such as information self-locking, poor grounding, or failure to recognize knowledge gaps~\citep{shao2025llm,zou2026informationselflockingreinforcementlearning}. Web-search, tool-use, and embodied environments expose richer interactions~\citep{bae2025alignsearchbeliefguidedexploratory,hua2026learningexplorescalingagentic,jinsearch,rezende2014stochastic,song2026large}, but their complexity can obscure which action improved the agent's knowledge state. Our \textbf{Clue Selector Game} (CSG) is a controlled complement: it preserves a language-mediated multi-turn loop---action, response, belief update, and subsequent action---while making the hypothesis space explicit, the posterior exactly measurable, and process-level metrics such as grounding, redundancy, zero-information behavior, recovery, and calibration directly computable.}
 

 





% \paragraph{Controlled environments for information-seeking agents.}
% Progress in multi-turn language-agent RL requires controllable environments that expose long-horizon interaction and reasoning skills such as asking useful follow-up questions, gathering evidence, and acting under uncertainty~\citep{abdulhai2023lmrl,shojaee2026illusion, yue2026does}. Recent work shows why final accuracy alone is insufficient~\citep{zhang-etal-2024-probing,lialfa, choudhury2025bed}: outcome-based RL can produce “information self-locking”, where agents remain trapped in low-information interaction patterns~\citep{zou2026informationselflockingreinforcementlearning}, and agents may score well while still exhibiting poor epistemic behavior, such as hallucinating unsupported claims or failing to recognize knowledge gaps~\citep{shao2025llm}. Related belief-guided and exploration-aware methods address uncertainty in partially observable or embodied settings~\citep{rezende2014stochastic,bae2025alignsearchbeliefguidedexploratory,hua2026learningexplorescalingagentic}, while BARL motivates belief-conditioned adaptation for reflective reasoning~\citep{zhang2025beyond}. CSG follows this broader trend but deliberately isolates \textit{query-based} information seeking~\citep{andukuri2024stargateteachinglanguagemodels, mazzaccara-etal-2024-learning,zhang-etal-2024-probing}: the hidden variable is explicit, the candidate set is the belief state, and each arm--question action produces a measurable belief transition. This lets us directly evaluate whether an action reduced uncertainty, became redundant, or enabled recovery after a low-information turn, complementing more realistic but more complex setups like web-search, tool-use, and embodied-agent environments.
 



% \paragraph{Bayesian Learning in RL}
% BARL~\citep{zhang2025markovianreflectiveexplorationbayesadaptive} frames reasoning as a finite-horizon MDP over CoT prefixes, where exploration occurs within a single generated response. We study a complementary setting: multi-turn evidence-seeking agents whose actions elicit observations from an environment with inherent stochasticity (where the true transition function in the EDP is not fully deterministic). This distinction lets us directly evaluate whether RL induces adaptive exploration \textit{across} interaction turns~\citep{abdulhai2023lmrl}, rather than only within latent reasoning traces.

% \paragraph{Multiturn RL with GRPO and Environments}


% \paragraph{Reasoning-length vs accuracy in RL}
% Recent work has questioned the assumption that longer explicit reasoning traces are
% necessary for stronger LLM reasoning. Several RL-based methods reduce overthinking by
% penalizing or controlling reasoning length, showing that models can preserve accuracy
% while producing substantially shorter chains of thought. Other work finds that
% ``no-thinking'' or short-thinking inference can be competitive under fixed token budgets,
% and position papers have cautioned against treating intermediate tokens as faithful
% reasoning traces. These findings suggest that visible reasoning length is an unreliable
% proxy for reasoning competence~\citep{palod2025performative,samineni2025localcoherenceglobalvalidity}. However, existing concise-reasoning work primarily
% measures final-answer accuracy and token efficiency. In contrast, we ask whether a model
% that emits little or no reasoning text can still act adaptively under uncertainty. We
% call this phenomenon \emph{silent exploration}: belief-sensitive information gathering
% without verbalized rationales. 


\begin{wrapfigure}{r}{0.45\linewidth}
    \vspace*{-2mm}
    \centering
    \includegraphics[
        width=\linewidth,
        trim=0.3in 0.25in 0.3in 0.25in, clip,
        clip
    ]{figures/edp_highlevel.pdf}
    \caption{\small
EDP intuition. Two histories can share the \an{same belief-independent task context} \(s(h_t)=s(h'_t)\) while inducing different beliefs \(b_t\neq b'_t\) over the latent task variable. An \textit{epistemically adaptive} policy conditions on the belief state at each turn $t$ to choose different useful actions, so \(a^\star(b_t)\neq a^\star(b'_t)\). A belief-agnostic policy that depends only on \(s(h_t)\) and \(t\), but not \(b_t\), cannot make this distinction.
}
  
     % In our clue-selection task, the belief corresponds to the current candidate set.
    \label{fig:edp}
    \vspace*{-6mm}
\end{wrapfigure}

\vspace*{-3mm}
\section{Epistemic Decision Processes}
\label{sec:edp_theory}
\vspace*{-3mm}

% \[
% \boxed{
% \underbrace{s(h_t)=s(h'_t)}_{\text{same visible state}},
% \quad
% \underbrace{b_t \neq b'_t}_{\text{different beliefs}}
% \;\Longrightarrow\;
% \underbrace{a^\star(h_t)\neq a^\star(h'_t)}_{\text{different optimal actions}}
% }
% \]

% \[
% \pi^\star(a_t\mid h_t)=\mu(a_t\mid s(h_t),b_t,t).

% s(h_t)=s(h'_t),\; b_t\neq b'_t
% \;\Rightarrow\;
% a^\star(b_t)\neq a^\star(b'_t)
% \]
 

We now formalize multi-turn information seeking as a decision problem over beliefs. An agent acts under uncertainty about a latent task variable \(z^\star\), chooses actions that elicit external observations, and updates its belief \nk{from those observations} before choosing the next action. The key EDP property is that the main dense reward is \textit{posterior-sensitive}: the same question can be useful under one candidate set and uninformative under another.

% The central feature of this setting is that action value is belief-dependent: two histories may look similar at the surface level while inducing different posteriors, and therefore requiring different useful actions.


% Many real-world LLM agents operate as information-seeking systems: they search, ask questions, call tools, inspect intermediate results, and revise their next actions based on new evidence. Epistemic Decision Processes capture this interaction pattern by treating uncertainty itself as part of the decision state. In this work, we study “multi-turn” information-seeking agents that must act under uncertainty about a latent task variable.

 

This differs from single-response reasoning settings, where adaptation occurs primarily within a generated trace before feedback is received~\citep{zhang2025markovianreflectiveexplorationbayesadaptive}. In an EDP, the policy interacts with an external environment over multiple turns: each action produces an observation, the observation changes the posterior, and the next action should be chosen under the updated belief~\citep{abdulhai2023lmrl}. 
\nk{Intuitively, optimal policies in this setting are}
%Our goal is to characterize and train 
\emph{epistemically adaptive}: they choose actions because they are useful under the current belief, not merely because they are plausible \nk{given the observable} surface history.


% Unlike single-response reasoning settings~\citep{zhang2025markovianreflectiveexplorationbayesadaptive}, which resemble contextual bandit-style language RL~\citep{samineni2025rlonlyanalyzingstructural} in that an agent generates a complete chain of thought before receiving feedback, the agent in our setting interacts with an external environment over multiple turns~\citep{abdulhai2023lmrl}: each action elicits an observation, and this observation changes the agent's uncertainty before the next action is chosen. \textbf{Our goal is to learn an epistemically adaptive policy, which we hypothesize will perform well in long-horizon information seeking tasks with stochastic environment feedback that reflects real world agentic tasks. }

%\subsection{Epistemic Decision Process.}

We formalize this as an \emph{Epistemic Decision Process} (EDP), a finite-horizon interaction defined by:
\[
    \mathcal{E}=\left(\mathcal{Z},\rho_0,\mathcal{A},\Omega,O,u,T\right).
\]

Here, \(z^\star \in \mathcal{Z}\) denotes the latent task variable sampled from prior \(\rho_0\)---for example, a target hypothesis, user intent, diagnosis, or answer state---while \(\mathcal{A}\), \(\Omega\), \(O\), \(u\), and \(T\) denote the action space, observation space, observation kernel, evidential utility, and horizon. At turn \(t\), the interaction history \(h_t=(a_0,o_0,\ldots,a_{t-1},o_{t-1})\) induces a posterior belief \(b_t(z):=\mathbb{P}(z^\star=z \mid h_t)\).\footnote{In LLM implementations, the serialized prompt may include a textual state summary, such as remaining hypotheses. We treat this as an encoding of the epistemic component \(b_t\).} The agent samples an action \(a_t \sim \pi_\theta(\cdot \mid h_t)\), receives an external observation \(o_t \sim O(\cdot \mid z^\star,h_t,a_t)\), and \an{updates its posterior via a Bayesian belief-update operator \(b_{t+1}=\tau(b_t,h_t,a_t,o_t)\), where \(\tau\) is induced by \(O\). }Fig.~\ref{fig:edp} shows a conceptualization of the EDP.
\iffalse
Here, \(z^\star \in \mathcal{Z}\) is a latent task variable sampled from a prior \(\rho_0\), \(\mathcal{A}\) is the action space, \(\Omega\) is the observation space, \(O(o_t \mid z^\star,h_t,a_t)\) is the observation kernel, \(u\) is an evidential utility function, and \(T\) is the horizon. Intuitively, \(z^\star\) is the hidden task state the agent is trying to discover through interaction, i.e., the goal state. At turn \(t\), the agent has observed an interaction history \(h_t=(a_0,o_0,\ldots,a_{t-1},o_{t-1})\), which induces a posterior belief over the latent task variable\footnote{In LLM implementations, the serialized prompt may include a textual state summary, such as remaining hypotheses. We treat this as an encoding of the epistemic component \(b_t\).}, \(b_t(z):=\mathbb{P}(z^\star=z \mid h_t)\). The agent samples an action \(a_t \sim \pi_\theta(\cdot \mid h_t)\), receives an external observation \(o_t \sim O(\cdot \mid z^\star,h_t,a_t)\), \an{and updates its posterior belief via a belief-update operator \(b_{t+1}=\tau(b_t,h_t,a_t,o_t)\).  Here, \(\tau\) denotes the belief-update operator induced by the observation kernel \(O\).} Fig.~\ref{fig:edp} shows a conceptualization of the EDP.
\fi

\vspace*{-3mm}
\paragraph{LLM Instantiation}
The policy \(\pi_\theta\) is an autoregressive LLM that gets a serialized history \(h_t\), including the task instruction, current state summary, and previous action--observation pairs. The model generates a structured output that is parsed as the action \(a_t\), such as a query or information-seeking move. The observation \(o_t\) from the environment %produced in our experiments by an oracle/verifier model, % distracting in this context - introduce later
\nk{is used to update} the environment state by filtering the hypotheses still consistent with the evidence. Thus, the transition is not merely deterministic text concatenation~\citep{samineni2026rlonlyanalyzingstructural}: although the next prompt appends \((a_t,o_t)\), the observation itself is produced by an external kernel \(o_t\sim O(\cdot\mid z^\star,h_t,a_t)\). Each turn is therefore a genuine environment interaction, and the decision-relevant state is the posterior uncertainty induced by the history, i.e., the augmented epistemic state \((s(h_t),b_t,t)\), rather than only the visible text or turn index.

 
% \paragraph{LLM instantiation.}
% In our language-agent setting, the policy \(\pi_\theta\) is an autoregressive LLM. At turn \(t\), the input to the model is the serialized interaction history \(h_t\), formed by concatenating the task instruction, the current state summary, and all previous action--observation pairs \((a_i,o_i)_{i<t}\). The model samples tokens autoregressively until a structured action is completed; this parsed output constitutes the action \(a_t\), such as a query or information-seeking move. After the action is issued, the environment returns an external observation \(o_t\). In our experiments, this observation is produced by an oracle/verifier model and then used to update the environment state, for example by reducing the set of hypotheses still consistent with the evidence. Thus the true transition function is \textit{not} simply deterministic text concatenation~\citep{samineni2025rlonlyanalyzingstructural}: the next history is \(h_{t+1}=h_t \oplus (a_t,o_t)\), but the observation is generated according to an external observation kernel \(o_t \sim O(\cdot \mid z^\star,h_t,a_t)\), which captures the stochasticity of the oracle response and the evidence it provides that are part of the environment. This makes each turn a genuine environment interaction rather than merely another segment of a single generated response. Because each external observation changes what the agent knows about the latent task variable, the relevant state for decision-making is not just the text history itself, but the posterior uncertainty induced by that history. \textbf{The central object in an EDP is therefore the evolving belief state \(b_t\), or more generally the augmented epistemic state \((s(h_t),b_t,t)\), rather than merely the visible surface state or turn index.}
\vspace{-3mm}

\paragraph{Bayesian Value and Advantage}
\label{ssec:bayesian_value_advantage}
To train an epistemically adaptive policy (Fig.~\ref{fig:edp}), we need a credit signal that evaluates actions under the agent's current posterior belief, rather than only under the final episode outcome. For a fixed latent task \(z\), let
\(Q_z^\pi(h_t,a_t)=\EE_{\pi,O}[\sum_{s=t}^{T-1}u(b_s,a_s,o_s)+R_T(h_T,z)\mid h_t,a_t,z^\star=z]\)
denote the value of taking action \(a_t\) after history \(h_t\), where \(R_T\) is any terminal task reward. Since the agent does not observe \(z^\star\), the relevant value is posterior-weighted: \(Q_{\mathrm{B}}^\pi(h_t,a_t)=\EE_{z\sim b_t}[Q_z^\pi(h_t,a_t)]\). The corresponding Bayesian advantage is \(A_{\mathrm{B}}^\pi(h_t,a_t)=Q_{\mathrm{B}}^\pi(h_t,a_t)-\EE_{a'\sim\pi(\cdot\mid h_t)}[Q_{\mathrm{B}}^\pi(h_t,a')]\). \an{
Since \(b_t=b(h_t)\), the dependence on the current belief is implicit in the history argument; equivalently, \(A_{\mathrm{B}}^\pi(h_t,a_t)\) evaluates the action under the posterior induced by \(h_t\).} Intuitively, this asks whether \(a_t\) is better than alternative actions available under the same posterior belief. At first glance, estimating this quantity exactly would require branching multiple actions from the same history \(h_t\), which is expensive in multi-turn interactive environments for LLMs~\citep{abdulhai2023lmrl,guo2025segment,kazemnejad2025vinepporefiningcreditassignment,zhang2025markovianreflectiveexplorationbayesadaptive}. This raises two crucial questions: what policy class is this Bayesian advantage meant to favor, and how can we approximate its credit signal without exhaustive same-history branching? We first answer the policy question by distinguishing belief-agnostic policies from epistemically adaptive ones; we then introduce a tractable Monte Carlo policy-gradient estimator that uses turn-level epistemic rewards as a practical surrogate for the belief-conditioned advantage.
\vspace{-4mm}

 \subsection{Distinguishing Epistemically Adaptive from Belief-Agnostic Policies}
\label{ssec:epistemically_adaptive_policies}
 \vspace{-3mm}
 
\an{\nk{An {\it epistemically adaptive policy} uses the {\it belief} induced by the interaction history, not just the history itself.} Let \(s(h_t)\) denote the belief-independent task context, as shown in Fig.~\ref{fig:edp}. These variables constrain what the agent can do at turn \(t\), but do not themselves encode the posterior over \(z^\star\). A policy is \emph{belief-agnostic} if its dependence on history factors only through this interface state, i.e., there exists a decision rule \(g\) such that \(\pi(a_t\mid h_t)=g(a_t\mid s(h_t))\). In contrast, a policy is \emph{epistemically adaptive} if its action distribution can also depend on the posterior belief induced by the history, i.e., there exists a decision rule \(f\) such that \(\pi(a_t\mid h_t)=f(a_t\mid s(h_t),b_t)\). We say the policy is \emph{strictly} epistemically adaptive if there exist histories \(h_t,h'_t\) with the same interface state \(s(h_t)=s(h'_t)\) but different beliefs \(b_t\neq b'_t\), such that \(\pi(\cdot\mid h_t)\neq\pi(\cdot\mid h'_t)\).}
 
 

% We now define what we mean by epistemically adaptive policies. Specifically, we distinguish policies by whether they use the belief induced by interaction history. A policy is \emph{belief-agnostic} if \(\pi(a_t\mid h_t)=\mu(a_t\mid s(h_t),t)\), where \(s(h_t)\) denotes the non-epistemic interface state of the task, such as the action schema, tool interface, turn index, and remaining budget. In contrast, a policy is \emph{epistemically adaptive} if its action distribution can also depend on the posterior belief, \(\pi(a_t\mid h_t)=\mu(a_t\mid s(h_t),b_t,t)\). We say it is \emph{strictly} epistemically adaptive if there exist histories \(h_t,h'_t\) with the same interface state \(s(h_t)=s(h'_t)\) but different beliefs \(b_t\neq b'_t\), such that \(\pi(\cdot\mid h_t)\neq\pi(\cdot\mid h'_t)\). 

\begin{theorybox}
    

\paragraph{Silent exploration.}
When the policy adapts to the posterior without emitting an explicit natural-language rationale, \nk{we call this behavior \emph{silent exploration}}. A silently exploratory policy changes what it does in response to evidence, even if its output is only a compact structured action. This matters because visible reasoning is neither necessary nor sufficient for epistemic adaptivity: as our experimental results suggest (Sec.~\ref{sec:experiments}), a model may produce long explanations without acting on the current belief, while another may output no rationale and still act adaptively.

\end{theorybox}


\begin{theorybox}
\begin{proposition}[Belief-state sufficiency]
Assume the observation kernel and utility function depend on \(h_t\) only through the surface state \(s(h_t)\) and belief \(b_t\). Then there exists an optimal policy $\pi^\star$ for the EDP that is Markovian in the augmented epistemic state \((s(h_t),b_t,t)\). Full proof is in Appendix~\ref{app:proofs}.
\label{proposition:belief_sufficiency_main}
\end{proposition}
\end{theorybox}
\vspace{-3mm}

\paragraph{Remark on Proposition~\ref{proposition:belief_sufficiency_main}} Once the belief \(b_t\) is part of the state, optimal behavior can be written as a Markov policy over \((s(h_t),b_t,t)\). The important question is therefore \textbf{not} whether a Markov policy exists, but whether the training signal can teach the policy to use the belief as uncertainty is resolved~\citep{zhang2025markovianreflectiveexplorationbayesadaptive}. In long-horizon information-seeking tasks, this is a depth-wise credit-assignment problem: at each turn, the useful action depends on the current posterior, and failing to adapt even once can derail the remaining trajectory. The first part of Theorem~\ref{thm:necessity} makes this sharp: there are EDPs where an adaptive policy succeeds with probability \(1\), while any belief-agnostic policy succeeds with probability at most \(2^{1-d}\), decaying exponentially with horizon. The second part of Theorem~\ref{thm:necessity}  then shows why sparse trajectory-level returns are a poor learning signal for this behavior: they reveal final success, but not the belief-conditioned credit of each intermediate action. Together, these results motivate turn-level epistemic credit rather than trajectory-level credit.

% \clearpage
\begin{theorybox}
\begin{theorem}[Adaptive policies require epistemic credit]
\label{thm:necessity}
Let $\mathcal{E}$ be an Epistemic Decision Process.
\begin{enumerate} 
    \item \textbf{(Performance gap.)}
    For every depth $d \geq 1$, there exists an EDP $\mathcal{E}_d$ with
    horizon $d$ in which an epistemically adaptive policy achieves
    $\mathbb{P}(\mathrm{success}) = 1$, while every belief-agnostic policy
    satisfies $\mathbb{P}(\mathrm{success}) \leq 2^{1-d}$. Thus, the cost of ignoring belief compounds exponentially with the decision horizon.

    \item \textbf{(Estimator insufficiency.)}
    Let $G = \sum_{s=0}^{T-1} r_s$ be the total episode return and fix
    any turn $t$. For any trajectory-level \nk{advantage} estimator
    $\widehat{A}_t^{\mathrm{traj}} = \phi(G)$,
    \begin{align}
        &\notag\inf_{\phi}\,
        \mathbb{E}\!\left[
            \bigl(\phi(G) - A_{\mathrm{B}}^\pi(h_t,a_t)\bigr)^2
        \right]\\
        = &\notag\mathbb{E}\!\left[
            \operatorname{Var}\!\left(
                A_{\mathrm{B}}^\pi(h_t,a_t) \mid G
            \right)
        \right] > 0
    \end{align}
 whenever the same total return can arise from belief-distinct
decisions with different Bayesian advantages.
\end{enumerate}

Together, (i) and (ii) show that, for a family of EDPs where a useful next action depends on the current posterior, epistemic adaptivity is necessary for strong long-horizon performance, and aggregate trajectory return is too coarse to recover the per-turn Bayesian advantage. It is not that full-return policy gradients are invalid for optimizing expected return; rather, total return alone aliases belief-distinct decisions that require different local credit. Full proofs are in Appendix~\ref{app:proofs}. 
% \an{This distinction is also visible in Fig.~\ref{fig:clue_overview}: \textsc{ECHO} receives useful local credit because its questions steadily increase information gain, whereas episodic-return GRPO can assign similar trajectory-level credit to turns that are fluent but epistemically redundant. 
 
\end{theorem}
\end{theorybox}

\vspace*{-3mm}
\section{Depth-wise Epistemic Credit Optimization (\textsc{ECHO})}
\label{sec:echo}
\vspace*{-3mm}
Theorem~\ref{thm:necessity} shows that aggregate trajectory return can alias belief-distinct decisions. More generally, two trajectories can have the same final reward while containing very different intermediate evidence-gathering behavior, so total return alone cannot identify which actions reduced uncertainty. This leads to our core insight—\textit{actions should receive credit based on whether they were useful under the belief state in which they were chosen, not only based on whether the entire episode eventually succeeded}. We therefore introduce \textsc{ECHO} (\textbf{E}pistemic \textbf{C}redit for \textbf{H}istory-Conditioned \textbf{O}ptimization), a practical turn-level surrogate for Bayesian advantage in multi-turn EDPs. Like GRPO~\citep{shao2024deepseekmath}, \textsc{ECHO} uses group-normalized advantages and a clipped policy-gradient objective~\citep{schulman2017proximalpolicyoptimizationalgorithms}; unlike trajectory-level GRPO, it normalizes posterior-sensitive turn rewards rather than full episode returns. As such, \textsc{ECHO} keeps the practical group-relative clipped optimization form of GRPO, but shifts the credit signal from trajectory-level success to turn-level “epistemic” progress. 

 % Unlike trajectory-level GRPO, however, the normalization is applied to posterior-sensitive turn rewards rather than full episode returns.

Specifically, for each training instance, we sample \(G\) parallel rollouts from the old policy. At turn \(t\), rollout \(g\) has history \(h_t^{(g)}\), belief \(b_t^{(g)}=\mathbb{P}(z^\star\mid h_t^{(g)})\), sampled action \(a_t^{(g)}\sim\pi_{\theta_{\mathrm{old}}}(\cdot\mid h_t^{(g)})\), and external observation \(o_t^{(g)}\). We assign a posterior-sensitive turn reward:
\[
    r_t^{(g)}
    =
    u\!\left(b_t^{(g)},a_t^{(g)},o_t^{(g)}\right),
\]
which measures local epistemic progress, such as information gain, reduction in posterior uncertainty, or elimination of inconsistent hypotheses. Sec.~\ref{ssec:clue_task} shows how this reward design fits into an EDP environment. At each decision depth \(t\), we normalize these rewards across the active rollouts: 

% \begin{theorybox}
    \begin{equation}
    \widehat{A}_{t}^{(g),\mathrm{ECHO}}
    =
    \frac{
        r_t^{(g)}
        -
        \frac{1}{|\mathcal{G}_t|}
        \sum_{j\in\mathcal{G}_t} r_t^{(j)}
    }{
        \operatorname{std}_{j\in\mathcal{G}_t}(r_t^{(j)})+\epsilon
    }
    \label{eq:echo_advantage}
\end{equation}
% \end{theorybox}
Unlike trajectory-level GRPO, which broadcasts a single full-episode advantage across all actions in a rollout, \textsc{ECHO}—as shown in Eq.~\ref{eq:echo_advantage}—computes group-normalized advantages at each decision depth. This gives each structured turn action credit according to the epistemic progress it produced under the current posterior~\citep{jin2026dgpo, nath2026owen, samineni2026rlonlyanalyzingstructural}. We note that our approach is not an exact belief-local advantage after rollouts diverge, since different rollouts may reach different histories and beliefs at the same turn—given the exponential state-action spaces in language based RL using LLMs~\citep{dong2024rlhf, srivastava2024policy, zhou2024archer}. The key point that each reward is evaluated relative to its rollout's own posterior, so the update preserves belief-dependent variation in action quality while avoiding exhaustive same-history branching.
%
Let
\[
    \rho_t^{(g)}(\theta)
    =
    \frac{
        \pi_\theta(a_t^{(g)}\mid h_t^{(g)})
    }{
        \pi_{\theta_{\mathrm{old}}}(a_t^{(g)}\mid h_t^{(g)})
    },
\]
where for LLM agents \(\pi_\theta(a_t^{(g)}\mid h_t^{(g)})\) is the product of token probabilities for the structured action. \textsc{ECHO} optimizes the clipped objective:
\begin{theorybox}
\begin{align}
    \mathcal{J}_{\mathrm{ECHO}}(\theta)
    =
    &\notag\mathbb{E}
    \Bigg[
    \frac{1}{\sum_t |\mathcal{G}_t|}
    \sum_{t=0}^{T-1}
    \sum_{g\in\mathcal{G}_t}
    \min\!\Big(
        \rho_t^{(g)}\widehat{A}_{t}^{(g)},\\
        &\operatorname{clip}\!\left(
            \rho_t^{(g)},
            1-\eta,
            1+\eta
        \right)
        \widehat{A}_{t}^{(g)}
    \Big)
    \Bigg]
    \label{eq:echo-objective}
\end{align}
\textbf{Remark.}~\textsc{ECHO} retains the group-relative clipped 
trust-region form of standard policy gradient methods~\citep{schulman2015trust,sutton2018reinforcement,williams92}, but replaces 
trajectory-level success with turn-level epistemic progress as the 
credit signal.
\end{theorybox}
 


\vspace*{-3mm}
\section{Experiments}
\label{sec:experiments}
\vspace*{-3mm}
\subsection{Clue Selector Game Environment as an EDP}
\label{ssec:clue_task}

\vspace*{-3mm}
\begin{figure*}[t]
    \centering
    \includegraphics[width=\textwidth]{figures/csg_final.png}
\vspace*{-3mm}
    \caption{
    Clue Selector Game overview and example trajectory from real evaluation data: \textsc{ECHO} contracts the belief state, while episodic-return GRPO stalls with zero-elimination queries.
    }
    \label{fig:clue_overview}
\vspace*{-3mm}
\end{figure*}

We instantiate EDPs through the \emph{Clue Selector Game} (CSG), a controlled multi-turn information-seeking task inspired by prior work in multiturn RL~\citep{abdulhai2023lmrl}. In the CSG, the agent must identify a hidden secret number by sequentially querying a clue-holder oracle that knows the secret but is not allowed to reveal it directly. Instead, the oracle can only answer property-based questions about the secret. To make the action space structured\footnote{The action space is still diverse due to the query component, which admits open-ended question generation.}, the agent acts through an \emph{arm--question interface}: at each turn it routes a natural-language property question through one of several clue-holder arms, where each arm ``owns'' certain domains of information. This setup lets us test whether a policy can do more than ask plausible questions in isolation: under the current candidate set, a successful agent identifies {\it what information is most useful} to seek (the arm), avoids repeatedly asking for information already obtained, and adapts its query strategy as its belief over the secret changes. The arms therefore provide a lightweight yet explicit scaffolding to examine the adaptivity of the query strategy. Figs.~\ref{fig:clue_overview} and~\ref{fig:csg-prompts} show a high-level overview of the CSG environment and the prompts used, respectively. 

%\vspace*{-3mm}
\paragraph{Game dynamics}  A secret integer \(n^\star \in \{1,\ldots,N\}\) is sampled uniformly at random; we use \(N=100\) for training, while evaluation can include secrets up to \(N=200\) to test generalization. \an{At each turn, the agent poses a natural-language property
question to a fixed oracle that answers truthfully without
revealing the secret. The action space pairs each question
with a categorical arm label that groups questions by
property family (e.g., parity, divisibility, range);
this provides lightweight structure over the action space (see above)
but does not affect the oracle's response. The dominant
reward is posterior-sensitive candidate elimination; a
small shaping bonus for arm--question consistency is
included but is not the focus of our analysis.}


% At each turn, the LLM observes the current candidate set \(\mathcal{C}_t\), the previous arm--question--answer history, and a fixed set of clue-holder arms. The agent's action is an arm--question pair \(a_t=(k_t,q_t)\), where \(k_t\) selects an arm and \(q_t\) is a natural-language property question. The arms provide a lightweight “routing scaffold” over coarse property families, such as parity/range, primality/squares, divisibility, digit properties, and interval queries. The selected question is then sent to a fixed clue-holder oracle\footnote{We use GPT-4o mini as the oracle.}, which answers truthfully about the queried property without revealing the secret directly. The environment then removes candidates inconsistent with the returned observation. Thus, the arm does not by itself determine the oracle answer; rather, it structures the policy's action space and lets us test whether models route questions through semantically appropriate arms while still adapting those questions to the current belief state. The episode ends when \(|\mathcal{C}_t|=1\), meaning the secret has been resolved, or when the turn limit is reached. 

% At each turn, the LLM observes the current candidate set \(\mathcal{C}_t\), the previous question--answer history, and a fixed set of clue-holder arms. Each arm corresponds to a family of yes/no-style questions, such as parity, primality, divisibility, digit properties, or range. The agent must select an arm and issue a natural-language question. A fixed oracle/verifier\footnote{We use the GPT-4o mini model as the oracle.} responds truthfully to the queried property, and the environment removes candidates inconsistent with the observation. The episode ends when \(|\mathcal{C}_t|=1\), meaning the secret has been resolved, or when the turn limit is reached. 

\vspace*{-3mm}
\paragraph{CSG's explicit EDP interpretation} The latent task variable is the secret number \(z^\star=n^\star\). The history \(h_t\) consists of the previous arm selections, questions, and oracle observations. The belief state \(b_t\) is the posterior over the secret number induced by this history. Because the secret is sampled uniformly and oracle observations only eliminate inconsistent candidates, the posterior is uniform over the remaining candidate set:
\[
    b_t(n)=
    \begin{cases}
    1/|\mathcal{C}_t|, & n\in\mathcal{C}_t,\\
    0, & n\notin\mathcal{C}_t.
    \end{cases}
\]
For example, if the initial candidates are \(\{1,\ldots,10\}\) and the agent asks whether the number is even, a truthful ``yes'' response updates the candidate set to \(\{2,4,6,8,10\}\), making the posterior uniform over those five numbers. Thus, in CSG, the candidate set is not merely a textual state variable; it is a direct representation of the agent's epistemic state.

Turn-level utility measures epistemic progress. The dominant reward term is fractional candidate elimination, \(r_{\mathrm{info}}=(|\mathcal{C}_t|-|\mathcal{C}_{t+1}|)/|\mathcal{C}_t|\), which measures how much the action contracts the current posterior. Training also includes small shaping terms for valid arm--question matching and query diversity, a redundancy penalty for repetition, and sparse bonuses for resolving the task efficiently. \an{The action value is therefore belief-dependent: two histories may share the same belief-independent task context---for example, the same turn index, action schema, and remaining budget---while inducing different posteriors over \(z^\star\), and therefore requiring different useful actions. }
%The key EDP property is that the main dense reward is \textit{posterior-sensitive}: the same question can be useful under one candidate set and uninformative under another. 
The posterior-sensitive reward signal is used by \textsc{ECHO} and other RL baselines for training; it assigns credit based on the epistemic progress made at the turn where the action was chosen, rather than broadcasting only the final episode outcome. CSG is therefore a clean testbed for our theory. 

\vspace*{-3mm}
\paragraph{Evaluation Strategy}
Our metrics test whether a policy uses accumulated evidence to choose useful actions throughout the trajectory, rather than merely succeeding at the end or producing more reasoning text.
%Our goal is to evaluate agents on both task success and epistemic adaptability in a multi-turn setup. 
In CSG, optimal information seeking is not determined by turn number alone, but by the structure of the current posterior over candidates, \an{i.e., the specific composition of the surviving candidate set $\mathcal{C}_t$}. An \textit{epistemically adaptive} agent should therefore ask questions that reduce uncertainty under the current \(\mathcal{C}_t\), avoid redundant queries, recover after low-information observations, and not ask broad questions if the belief state is narrow. Accordingly, we evaluate not only \textbf{final task resolution}, but also \textbf{process-level epistemic quality}~\citep{shao2025llm} as a proxy for epistemic adaptability (see Sec.~\ref{sec:results}). We report \textbf{Zero}, the rate of non-redundant turns that eliminate no candidates; \textbf{Qual.}, the candidate-elimination quality normalized by an ideal split; and \textbf{Ground}, the rate of turns that are both non-redundant and informative under the current candidate set. To measure robustness after local epistemic failures, we also report \textbf{GRecover}, the rate at which the next turn after a zero-elimination action is \textit{grounded} as defined above, and \textbf{ResAfterZero}, the fraction of episodes that still resolve after at least one zero-elimination event. Finally, \textbf{Reason\%} measures explicit reasoning text outside the required action format. \an{Further details on evaluation metrics are provided in Appendix~\ref{ssec:main_task_epistemic_metrics_full}.}
 
Empirically, these metrics are strongly tied to task success: across evaluation episodes, resolution correlates positively with elimination quality and groundedness, and negatively with zero-elimination and low-information behavior, supporting their use as indicators of epistemically adaptive behavior.

% \paragraph{Evaluation Strategy} Our main goal is to evaluate agents on \textit{both} task success as well as epistemic adaptability. Notably, the optimal information-seeking behavior is not determined by turn number alone, but by the structure of the current posterior over candidates. Therefore, an \textit{epistemically adaptive} agent should ask questions that reduce uncertainty under the current \(\mathcal{C}_t\), avoid redundant queries, recover after low-information observations, and stop asking broad questions once the belief is already narrow. Accordingly, we evaluate not only \textbf{final task resolution}, but also \textbf{process-level epistemic quality~\citep{shao2025llm}} as a proxy for epistemic adaptability. For the latter, we report groundedness, recovery after low-information or redundant actions, and late-stage calibration. These metrics are intended to capture epistemic adaptivity: whether the policy uses evidence to choose useful actions throughout the trajectory, rather than merely succeeding at the end. \an{one sentence here connecting to results showing high correlation with task success and epistemic adaptability metrics}
\vspace*{-3mm}
\paragraph{Baselines}
%We compare against baselines spanning five primary categories. 
\textbf{Frontier or proprietary models} include Claude Sonnet 4.6, OpenAI's o3-mini, GPT-4o, GPT-4.1-mini, GPT-4.1-nano, and GPT-4o-mini, evaluated zero-shot with the standard task prompt. These models represent strong off-the-shelf performance without task-specific training. \textbf{Prompting variants} evaluate Qwen2.5-1.5B Instruct with explicit reasoning prompts: chain-of-thought (CoT), full ReAct~\citep{yao2022react} using a complete Thought/Action/Observation loop with growing history context, and a simple version of ReAct using a Thought/Action format in the system prompt. These baselines test whether eliciting explicit reasoning can substitute for RL training on the same base model. We also include \textbf{GPT-4o w/ full ReAct}, which tests whether a structured reasoning loop improves a capable frontier model. \textbf{RL baselines} are trained on Qwen2.5-1.5B Instruct using the same environment and reward function as \textsc{ECHO}, differing only in the advantage estimator. The main RL baseline is \textbf{episode-return RL}, which broadcasts the total episode reward uniformly to all turns (and thereby, to tokens) consistent with the standard GRPO formulation when extended to multi-turn settings; we also include RLOO~\citep{ahmadian-etal-2024-back} per-turn to isolate the effect of per-turn relative credit. \textbf{Offline SFT+RL} is a single-turn offline baseline trained from pre-collected trajectories, representing training on fixed demonstrations rather than online interaction. \textbf{Base Qwen 2.5} at sizes from 0.5B--14B provide a scaling reference without task-specific prompting or RL.  Finally, a \textbf{greedy oracle} reference policy with access to the secret during action selection, that chooses the arm-question pair with maximal candidate elimination, serves as an “upper-bound”~\citep{analyzingnas} diagnostic.% rather than a directly comparable learned baseline.

 

\vspace*{-3mm}
\section{Results}
\label{sec:results}

\input{final_clue_results_provenance}
\vspace*{-3mm}
\paragraph{ECHO improves both resolution and epistemic behavior.}
Table~\ref{tab:csg_main_results} shows that \textsc{ECHO} achieves the strongest learned-policy performance in CSG. Its resolve rate of \(45.3\%\) exceeds the best frontier baseline, Claude Sonnet 4.6 (\(43.1\%\)), and substantially outperforms episodic-return GRPO (\(14.7\%\)), RLOO per-turn (\(13.8\%\)), and all base/prompting baselines. More importantly, the improvement is reflected in process-level epistemic behavior: \textsc{ECHO} has the lowest zero-elimination rate (\(18.6\%\)), indicating that its questions most often update the agent belief state. While Claude attains slightly higher per-turn quality and groundedness, \textsc{ECHO} combines high elimination quality (\(0.670\)) and groundedness (\(0.718\)) with much stronger trajectory-level robustness: after zero-elimination events, it has the highest grounded recovery (\(0.406\)) and resolution-after-zero rate (\(0.318\)). This suggests that \textsc{ECHO}'s advantage is not merely asking strong individual questions, but that \textit{the \textsc{ECHO}-trained policy learns more robust information-gathering behavior across the trajectory}: it asks well-grounded questions, avoids zero-elimination actions, and recovers when a local epistemic failure occurs. \textsc{ECHO} also achieves these gains with almost no explicit reasoning text (\(0.9\%\)) exposing its silently exploratory nature, whereas ReAct variants reason on every turn but do not improve resolution, indicating that explicit reasoning traces alone do not substitute for epistemically adaptive training. \an{Arm-level diagnostics further show that 
%\textsc{ECHO}'s gains are not driven by broad arm switching or static arm-label matching. As shown in \Cref{tab:arm_routing_belief_diagnostics} in Appendix~\ref{app:binary_search_echo}, 
\textsc{ECHO}'s arm--question choices remain informative even under same-arm reuse, suggesting belief-conditioned routing rather than superficial arm switching (full results in \Cref{tab:arm_routing_belief_diagnostics} in Appendix~\ref{app:binary_search_echo}).}
Additional results are reported in Appendix~\ref{app:training_dynamics} and~\ref{app:binary_search_echo}. 

\vspace*{-3mm}
\paragraph{Behavioral diagnostics reveal why \textsc{ECHO} succeeds.}
To understand whether \textsc{ECHO}'s gains come from genuine belief-state adaptation,
%rather than superficial arm exploration or explicit reasoning, 
we analyze process-level diagnostics across policies in Fig.~\ref{fig:csg_behavior_diagnostics}. Results show that resolution is tightly linked to posterior contraction: models with fewer zero-elimination turns and higher elimination quality achieve higher solve rates, with \textsc{ECHO} occupying the favorable region of high resolution and low wasted updates. Fig.~\ref{fig:csg_behavior_diagnostics}(c) further shows that \textsc{ECHO} continues shrinking the candidate set across turns, whereas episodic-return GRPO and prompting baselines flatten earlier, indicating stalled information acquisition. Fig.~\ref{fig:csg_behavior_diagnostics}(d) shows that \textsc{ECHO} is more robust after local epistemic failures: following zero-elimination events, it produces more grounded recoveries and is substantially more likely to still resolve the episode. Together, these diagnostics support the central mechanism behind \textsc{ECHO}: it learns to ask questions that remain informative under the evolving belief state, rather than being merely plausible.% or producing more explicit reasoning text.
\vspace*{-3mm}
\begin{figure*}[t]
    \centering
    \includegraphics[width=\textwidth]{figures/main_behavior_four_panel_clean.pdf}
    \vspace*{-3mm}
    \caption{
    \textbf{Behavioral diagnostics in CSG}. \textsc{ECHO} succeeds by reducing zero-elimination turns, maintaining high-quality belief contraction, shrinking the candidate set faster, and recovering more robustly after epistemic failures. Results on question type are shown in \Cref{tab:csg_hint_distribution_all} in Appendix~\ref{app:binary_search_echo}.}
    \vspace*{-3mm} 
    \label{fig:csg_behavior_diagnostics}
\end{figure*}

\begin{figure}[t]
    \centering

    \begin{subfigure}[t]{0.235\linewidth}
        \centering
        \includegraphics[width=\linewidth]{figures/rl_per_step/reward.pdf}
      
        \label{fig:csg_reward}
    \end{subfigure}
    \hfill
    \begin{subfigure}[t]{0.235\linewidth}
        \centering
        \includegraphics[width=\linewidth]{figures/rl_per_step/clue_episode_resolve_rate.pdf}
 
        \label{fig:csg_episode_resolve_rate}
    \end{subfigure}
    \hfill
    \begin{subfigure}[t]{0.235\linewidth}
        \centering
        \includegraphics[width=\linewidth]{figures/rl_per_step/clue_eliminated_mean.pdf}
 
        \label{fig:csg_redundancy_rate}
    \end{subfigure}
    \hfill
    \begin{subfigure}[t]{0.235\linewidth}
        \centering
        \includegraphics[width=\linewidth]{figures/rl_per_step/clue_redundancy_rate.pdf}
 
        \label{fig:csg_arm_match_rate}
    \end{subfigure}

    \vspace*{-3mm} 
    \caption{Performance and behavioral trends during \textbf{RL training} in the Clue Selector Game (CSG) environment. Plots show metrics SEM across 3 independent training runs. Additional training dynamics are shown in Fig.~\ref{fig:clue_steps_additional_metrics} in Appendix~\ref{app:training_dynamics}.}
    \vspace*{-3mm} 
    \label{fig:clue_steps_four_panel_training_main}
\end{figure}


% \paragraph{What does epistemic adaptability look like behaviorally?}

% \paragraph{ECHO achieves frontier-level epistemic quality without explicit reasoning.}

% Without leveraging a significant amount of reasoning tokens, our approach is able to outperform models 100X larger in the frontier range, as well as base models larger in the parameter scale. At just 1.5B scale, ECHO-trained models outperform strong frontier models including reasoning models like O3-mini as well as strong baselines like CoT~\citep{wei2023chainofthought} and React~\citep{yao2022react} with GPT-4o. This shows the adaptability in multiturn information-seeking tasks within an EDP setting requires “adaptability” as well as “recovery” from bad decisions—a behavior we see emerge in models trained with our approach. Notably, this boost is overall performance is seen even without requiring the costly amount of reasoning tokens before the final answer. This suggests that ECHO trains models to be adaptable using belief-sensitive reward structures and can provide strong performance with higher token cost-efficiency—a trend consistent with prior works~\citep{ma2025reasoningmodelseffectivethinking}.



% \paragraph{Training dynamics support turn-level epistemic credit assignment.}
% Finally, we examine whether these behavioral gains emerge during optimization. \Cref{fig:clue_steps_four_panel_training_main} shows that \textsc{ECHO}'s improvement is not merely a generic increase in reward, but a turn-level epistemic learning effect. Across training, \textsc{ECHO} increases episode resolution while also improving mean candidate elimination and reducing redundancy, whereas episodic-return GRPO remains comparatively flat and continues to produce more low-information behavior. By the end of training, \textsc{ECHO} achieves substantially higher resolution than episodic-return GRPO, eliminates more candidates per turn, and nearly removes redundant querying. This pattern matches the EDP perspective: when credit is assigned at the decision depth where an action changes the posterior, the policy learns to choose questions that contract the current belief state rather than relying on sparse trajectory-level outcomes.


\paragraph{What epistemically adaptive behavior looks like.}
The qualitative trajectory in Fig.~\ref{fig:clue_overview} illustrates the same mechanism at the episode level. For the same secret, \textsc{ECHO} asks questions that remain useful under the changing candidate set, progressively contracting the belief state until only the target remains, while episodic-return GRPO asks questions that are initially plausible, but then non-informative, producing several zero-elimination turns. Thus, the failure mode is \emph{stale information seeking}: otherwise reasonable questions that are no longer informative after earlier observations.

\vspace*{-3mm}
\paragraph{Explicit reasoning does not substitute for epistemic training.}
Explicit reasoning traces alone do not produce epistemic adaptability in long-horizon EDPs. Full ReAct forces models to reason on nearly every turn, yet GPT-4o + Full ReAct resolves only \(25.3\%\) of episodes. Qwen2.5-1.5B + Full ReAct resolves only \(4.4\%\). By contrast, \textsc{ECHO} resolves \(45.3\%\) of episodes and produces explicit reasoning text on only \(0.9\%\) of turns. This suggests that the central bottleneck for epistemic adaptivity is \textit{not} verbalizing intermediate reasoning, but learning a policy conditioned on the evolving belief state. This is “silent exploration” (Sec.~\ref{sec:edp_theory}) manifest: adaptive information seeking expressed through query choices and posterior-sensitive feedback rather than long test-time reasoning traces. \textsc{ECHO} trains adaptable models using belief-sensitive reward structures and provides strong performance with higher token cost-efficiency—a trend consistent with prior works~\citep{ma2025reasoningmodelseffectivethinking} in LLM reasoning.



\vspace*{-3mm}
\paragraph{Training dynamics support turn-level epistemic credit assignment.}
Fig.~\ref{fig:clue_steps_four_panel_training_main} 
%with additional training dynamics (\Cref{fig:clue_steps_additional_metrics}; Appendix~\Cref{app:training_dynamics}) 
shows that \textsc{ECHO}'s gains emerge as a turn-level epistemic learning effect rather than as a generic reward increase. By the end of training, \textsc{ECHO} reaches much higher episode resolution than trajectory-level GRPO (\(70.1\%\) vs.\ \(15.9\%\)), eliminates more candidates per turn (\(13.04\) vs.\ \(10.33\)), and reduces redundancy almost completely (\(0.2\%\) vs.\ \(2.5\%\)). It also resolves more often at the turn level (\(9.3\%\) vs.\ \(1.8\%\)) and shortens average episodes (\(7.67\) vs.\ \(9.29\) turns). These trends match our theoretical analysis of EDPs (Sec.~\ref{sec:edp_theory}): when credit is assigned at the decision depth where an action changes the posterior, the policy learns to ask questions that contract the current candidate set rather than merely optimize a sparse final outcome in a multiturn task. The optimization diagnostics suggest specialization rather than mode collapse~\citep{janus2022mysteries,zhang2025verbalized}: \textsc{ECHO}'s entropy decreases while query diversity remains near saturation, indicating that
the policy becomes more selective without repeating the same questions. In contrast, GRPO maintains higher entropy but achieves lower information gain and far lower resolution, consistent with trajectory-level credit failing to identify which intermediate actions were useful under the current belief. Overall, these results indicate that \textbf{aggregate trajectory returns can
fail to identify the per-turn Bayesian advantage needed for epistemic credit.}

\vspace*{-3mm}

\section{Conclusion}
\label{sec:conclusion}
\vspace*{-3mm}
We introduced Epistemic Decision Processes (EDPs), a belief-state formulation of multi-turn information seeking, and showed theoretically why trajectory-level success is insufficient for epistemic credit assignment: policies must receive credit at the decision points where actions change the posterior. Building on this view, we proposed \textsc{ECHO}, a turn-level optimization framework that rewards posterior-sensitive information gain and trains policies to be \textit{epistemically adaptive}. In the Clue Selector Game, \textsc{ECHO} trains a 1.5B policy that exceeds strong proprietary frontier baselines such as Claude Sonnet 4.6 and reasoning-observation focused baselines like ReAct on final task resolution, while also reducing zero-elimination turns, maintaining high-quality belief contraction, and recovering more robustly after local epistemic failures. These gains do not come from longer visible reasoning traces: \textsc{ECHO} exhibits ``silent exploration,'' learning adaptive information seeking through belief-conditioned actions rather than explicit rationales at inference time.

Overall, our results suggest that adaptive language agents in information-seeking long-horizon tasks should be trained and evaluated not only by whether they eventually succeed, but by how they gather evidence, update beliefs, and recover from uncertainty over time. Extending EDP-style credit assignment to richer tool-use, search, and multi-agent settings is a natural next step for future work. 

\vspace*{-3mm}
\section*{Limitations}
 
CSG provides exact Bayesian posteriors by construction
(Lemma~\ref{lemma:exact_posterior}), but real information-seeking
tasks---medical diagnosis, web search, preference
elicitation---involve implicit, high-dimensional hypothesis spaces
where the belief state must be
approximated~\citep{choudhury2025bed,choudhury2025processrewardmodelsllm,lee2025prints}, and where information
sources may be noisy or
adversarial. In such
settings, the posterior-sensitive turn reward that drives
\textsc{ECHO} could be challenging to compute analytically. %, reintroducing the
%need for learned process reward models that CSG was designed to
%avoid. 
Our experiments demonstrate \textsc{ECHO} at the 1.5B
scale; whether epistemic credit remains necessary at larger scales
where models may have stronger in-context belief
tracking~\citep{kossen2024incontext}, or provides complementary
gains, is an open question. Similarly, CSG operates over a
relatively short horizon (e.g., 10 total turns) with a structured arm--question action
space; scaling epistemic credit assignment to longer horizons with relatively
open-ended actions including general-purpose tool calls and richer observation
spaces~\citep{abdulhai2023lmrl,jinsearch} like the web is unexplored.
Extending EDPs to settings with approximate beliefs, unreliable
information sources, and compositional latent structure is a
natural direction for future work.

\begin{table*}[t]
\centering
\scriptsize
\setlength{\tabcolsep}{4.2pt}
\renewcommand{\arraystretch}{1.08}
\resizebox{\linewidth}{!}{
\begin{tabular}{lccccccc}
\toprule
\textbf{Condition} &
\textbf{Final Prog. $\uparrow$} &
\textbf{Completion $\uparrow$} &
\textbf{Oracle Match $\uparrow$} &
\textbf{Invalid $\downarrow$} &
\textbf{Parse Fail. $\downarrow$} &
\textbf{Prog. Gained $\uparrow$} &
\textbf{Ep. Length $\downarrow$} \\
\midrule

\multicolumn{8}{l}{\textit{Frontier baselines}} \\
\hdashline

Claude Sonnet 5
& 80.9 $\pm$ 0.5
& 31.7 $\pm$ 3.3
& \textbf{86.0 $\pm$ 2.7}
& 4.5 $\pm$ 1.5
& 0.3 $\pm$ 0.3
& 46.4 $\pm$ 1.4
& 16.9 $\pm$ 0.2 \\

Claude Sonnet 4.6
& \textbf{82.1 $\pm$ 0.5}
& 33.3 $\pm$ 6.0
& 83.6 $\pm$ 1.4
& 7.5 $\pm$ 0.6
& \textbf{0.0 $\pm$ 0.0}
& \textbf{47.6 $\pm$ 1.8}
& 16.4 $\pm$ 0.6 \\

GPT-5.4 mini
& 72.8 $\pm$ 0.8
& 23.3 $\pm$ 1.7
& 60.8 $\pm$ 1.8
& 27.4 $\pm$ 3.7
& \textbf{0.0 $\pm$ 0.0}
& 38.3 $\pm$ 0.5
& 17.9 $\pm$ 0.3 \\

GPT-4.1 mini
& 78.3 $\pm$ 1.6
& 28.3 $\pm$ 6.0
& 64.6 $\pm$ 3.4
& 25.8 $\pm$ 3.1
& 0.8 $\pm$ 0.2
& 43.8 $\pm$ 2.8
& 17.5 $\pm$ 0.4 \\

\midrule
\multicolumn{8}{l}{\textit{Base Qwen2.5 baselines}} \\
\hdashline

Qwen2.5-72B
& 63.1 $\pm$ 0.5
& 8.3 $\pm$ 1.7
& 66.3 $\pm$ 2.1
& 13.2 $\pm$ 2.0
& 28.8 $\pm$ 1.4
& 28.7 $\pm$ 1.8
& 19.4 $\pm$ 0.2 \\

Qwen2.5-32B
& 67.5 $\pm$ 0.6
& 11.7 $\pm$ 4.4
& 47.0 $\pm$ 3.4
& 35.6 $\pm$ 4.8
& 2.4 $\pm$ 0.2
& 33.0 $\pm$ 1.9
& 19.1 $\pm$ 0.3 \\

Qwen2.5-14B
& 67.6 $\pm$ 1.0
& 10.0 $\pm$ 2.9
& 43.0 $\pm$ 2.2
& 41.7 $\pm$ 2.0
& 4.1 $\pm$ 0.7
& 33.1 $\pm$ 1.5
& 19.1 $\pm$ 0.3 \\

Qwen2.5-7B
& 51.4 $\pm$ 0.6
& 1.7 $\pm$ 1.7
& 49.5 $\pm$ 1.7
& 14.7 $\pm$ 2.0
& 1.0 $\pm$ 0.6
& 17.0 $\pm$ 1.4
& 20.0 $\pm$ 0.0 \\

\midrule
\multicolumn{8}{l}{\textit{RL-trained Qwen2.5-7B models}} \\
\hdashline

Qwen2.5-7B + GRPO
& 74.8 $\pm$ 0.3
& 25.0 $\pm$ 5.0
& 44.9 $\pm$ 1.5
& 35.6 $\pm$ 1.2
& 2.0 $\pm$ 0.6
& 40.3 $\pm$ 1.5
& 17.8 $\pm$ 0.3 \\

\rowcolor{echoshade}
\textbf{Qwen2.5-7B + \textsc{ECHO} (Ours)}
& 80.2 $\pm$ 0.3
& \textbf{35.0 $\pm$ 5.0}
& 71.9 $\pm$ 3.0
& \textbf{3.2 $\pm$ 1.0}
& 0.4 $\pm$ 0.3
& 45.7 $\pm$ 1.6
& \textbf{16.2 $\pm$ 0.6} \\

\bottomrule
\end{tabular}
}
\caption{
\textbf{Main CRAFT results.}
Values are mean with SEM shown in subscript across 3 independent runs on the
evaluation set.
}
\label{tab:craft_results}
\end{table*}

 
% \section*{Acknowledgments}

% This document has been adapted
% by Steven Bethard, Ryan Cotterell and Rui Yan
% from the instructions for earlier ACL and NAACL proceedings, including those for
% ACL 2019 by Douwe Kiela and Ivan Vuli\'{c},
% NAACL 2019 by Stephanie Lukin and Alla Roskovskaya,
% ACL 2018 by Shay Cohen, Kevin Gimpel, and Wei Lu,
% NAACL 2018 by Margaret Mitchell and Stephanie Lukin,
% Bib\TeX{} suggestions for (NA)ACL 2017/2018 from Jason Eisner,
% ACL 2017 by Dan Gildea and Min-Yen Kan,
% NAACL 2017 by Margaret Mitchell,
% ACL 2012 by Maggie Li and Michael White,
% ACL 2010 by Jing-Shin Chang and Philipp Koehn,
% ACL 2008 by Johanna D. Moore, Simone Teufel, James Allan, and Sadaoki Furui,
% ACL 2005 by Hwee Tou Ng and Kemal Oflazer,
% ACL 2002 by Eugene Charniak and Dekang Lin,
% and earlier ACL and EACL formats written by several people, including
% John Chen, Henry S. Thompson and Donald Walker.
% Additional elements were taken from the formatting instructions of the \emph{International Joint Conference on Artificial Intelligence} and the \emph{Conference on Computer Vision and Pattern Recognition}.

% Bibliography entries for the entire Anthology, followed by custom entries
%\bibliography{custom,anthology-overleaf-1,anthology-overleaf-2}

% Custom bibliography entries only
\bibliography{abhijnan_lit_survey,colm2026_conference, rlhf, rlhf_2, nips_deception, nk_cites, nk_hannah_cites}

\appendix

\section{Theory and Proofs}
\label{app:proofs}
In this section, we provide our main theoretical results and proofs. 
\begin{proposition}[Belief-state sufficiency]
\label{prop:belief-state-sufficiency}
Let \(\mathcal{E}=(\mathcal{Z},\rho_0,\mathcal{A},\Omega,O,u,T)\) be an Epistemic Decision Process. Let \(h_t=(a_0,o_0,\ldots,a_{t-1},o_{t-1})\) denote the interaction history, and let \(b_t(z)=\mathbb{P}(z^\star=z\mid h_t)\) be the posterior belief over the latent task variable. Suppose there exists a surface-state map \(s:\mathcal{H}\rightarrow\mathcal{S}\) such that the following hold:
\begin{enumerate}
    \item \textbf{Observation sufficiency:} the observation kernel depends on the history only through the latent variable, surface state, and action:
    \[
        O(o_t\mid z^\star,h_t,a_t)=O(o_t\mid z^\star,s(h_t),a_t).
    \]
    \item \textbf{Surface-state update sufficiency:} there exists a transition map \(F:\mathcal{S}\times\mathcal{A}\times\Omega\rightarrow\mathcal{S}\) such that
    \[
        s(h_{t+1})=F(s(h_t),a_t,o_t).
    \]
    \item \textbf{Utility sufficiency:} the per-turn utility can be written as a measurable function of the surface state, belief, action, and observation:
    \[
        u_t=u(s(h_t),b_t,a_t,o_t).
    \]
    The terminal reward can likewise be written as \(R_T=R_T(s(h_T),b_T)\), or equivalently as the posterior expectation of a latent terminal reward.
\end{enumerate}
Then there exists an optimal policy $\pi^\star$ for the EDP that depends on the history \(h_t\) only through the augmented epistemic state \((s(h_t),b_t,t)\).
\end{proposition}

\begin{proof}
We prove the result by reducing the EDP to a belief-state MDP~\citep{aastrom1965optimal, smallwood1973optimal}.

\paragraph{Step 1: The augmented epistemic state determines future observation distributions.}
Fix any history \(h_t\), action \(a_t\), surface state \(s_t=s(h_t)\), and belief \(b_t\). Since \(b_t(z)=\mathbb{P}(z^\star=z\mid h_t)\), the predictive distribution over the next observation is
\[
    \mathbb{P}(o_t\mid h_t,a_t)
    =
    \int_{\mathcal{Z}}
    O(o_t\mid z,s_t,a_t)b_t(z)\,dz.
\]
Thus, by the observation sufficiency assumption, the distribution of \(o_t\) depends on the full history \(h_t\) only through \((s_t,b_t)\) and the chosen action \(a_t\).

\paragraph{Step 2: The belief update is determined by the augmented epistemic state.}
After observing \(o_t\), the posterior is updated by Bayes' rule:
\begin{align}
    b_{t+1}(z)
    &\notag=
    \tau(b_t,a_t,o_t)(z)\\
    &\notag=
    \frac{O(o_t\mid z,s_t,a_t)b_t(z)}
    {\int_{\mathcal{Z}}O(o_t\mid z',s_t,a_t)b_t(z')\,dz'}.
\end{align}
Therefore, \(b_{t+1}\) is fully determined by \((s_t,b_t,a_t,o_t)\). By the surface-state update assumption, \(s_{t+1}=s(h_{t+1})=F(s_t,a_t,o_t)\). Hence the next augmented epistemic state \((s_{t+1},b_{t+1},t+1)\) depends on the history only through \((s_t,b_t,t)\), the action \(a_t\), and the observation \(o_t\).

\paragraph{Step 3: The reward is determined by the augmented epistemic state.}
By the utility sufficiency assumption, the per-turn utility is \(u(s_t,b_t,a_t,o_t)\). Thus the distribution of both the next augmented state and the immediate utility is determined by \((s_t,b_t,t)\) and \(a_t\). No additional information from the raw history \(h_t\) is needed for predicting future rewards.

\paragraph{Step 4: Bellman recursion on the belief-state MDP.}
Define the optimal value function on the augmented epistemic state by backward induction~\citep{sutton2018reinforcement}. At the terminal step,
\[
    V_T^\star(s,b)=R_T(s,b).
\]
For \(t<T\), define
\begin{align}
    V_t^\star(s,b)
    =&\notag
    \max_{a\in\mathcal{A}}
    \mathbb{E}_{z\sim b}
    \mathbb{E}_{o\sim O(\cdot\mid z,s,a)}
    \Big[\\
        &\notag u(s,b,a,o)
        +
        V_{t+1}^\star(F(s,a,o),\\
        &\tau(b,a,o))
    \Big].
\end{align}
This Bellman recursion is well-defined because the predictive observation distribution, the belief update, the surface-state update, and the utility all depend only on \((s,b,t)\) and \(a\).

\paragraph{Step 5: Constructing an optimal policy.}
For any history \(h_t\), let \(s_t=s(h_t)\) and \(b_t=\mathbb{P}(z^\star\mid h_t)\). Define a greedy policy
\begin{align}
    \pi^\star(h_t)
    \in&\notag
    \arg\max_{a\in\mathcal{A}}
    \mathbb{E}_{z\sim b_t}
    \mathbb{E}_{o\sim O(\cdot\mid z,s_t,a)}
    \Big[\\
        &\notag u(s_t,b_t,a,o)
        +
        V_{t+1}^\star(F(s_t,a,o),\\
        &\tau(b_t,a,o))
    \Big].
\end{align}
By construction, \(\pi^\star(h_t)\) depends on \(h_t\) only through \((s(h_t),b_t,t)\). Moreover, by the Bellman recursion~\citep{sutton2018reinforcement}, this policy attains \(V_t^\star(s(h_t),b_t)\) for every history \(h_t\). Therefore, \(\pi^\star\) is optimal, and an optimal policy exists that is Markovian~\citep{zhang2025markovianreflectiveexplorationbayesadaptive} in the augmented epistemic state \((s(h_t),b_t,t)\).
\end{proof}

\paragraph{Significance.}
Proposition~\ref{prop:belief-state-sufficiency} justifies the EDP formulation. Although the agent observes a growing text history \(h_t\), the full history is not the fundamental decision state. What matters for optimal action selection is the visible surface state together with the posterior belief induced by the history. This separates ordinary text-state conditioning from epistemic adaptivity: a policy that ignores \(b_t\) may condition on the prompt and turn index, but it cannot in general choose actions based on what the agent has learned from external observations. The result is the EDP analogue of the standard belief-MDP reduction in Bayesian RL and POMDPs, but applied to multi-turn LLM agents whose observations come from an external environment rather than from internal chain-of-thought continuation.

\paragraph{Connection to Bayesian RL for LLM reasoning.}
This proposition parallels the belief-sufficiency argument used in Bayesian RL accounts of LLM reasoning. In BARL~\citep{zhang2025markovianreflectiveexplorationbayesadaptive}, the hidden uncertainty is over candidate MDPs or reasoning hypotheses, and the belief is updated over those hypotheses inside a generated reasoning trace. In our EDP setting, the hidden uncertainty is over an external latent task variable \(z^\star\), and the belief \(b_t(z)=\mathbb{P}(z^\star=z\mid h_t)\) is updated through external observations. Thus, the proof technique is the same belief-MDP reduction, but our setting is explicitly multi-turn: the agent acts, receives evidence from the environment, updates its posterior, and then chooses the next action.



\begin{theorem}[Adaptive policies can strictly outperform belief-agnostic policies]
\label{thm:adaptive-gap}
For every depth $d \geq 1$, there exists an Epistemic Decision Process $\mathcal{E}_d$ 
with horizon $d$ such that:
\begin{enumerate}
    \item there exists an epistemically adaptive policy $\pi^\star$ with 
          $\mathbb{P}(\mathrm{success} \mid \pi^\star) = 1$;
    \item every belief-agnostic policy $\mu$ satisfies 
          $\mathbb{P}(\mathrm{success} \mid \mu) \leq 2^{1-d}$.
\end{enumerate}
\end{theorem}

\begin{proof}
We construct a family of EDPs in which success requires using evidence 
revealed in previous turns. The construction separates policies that 
condition on the epistemic state from policies that do not.

\paragraph{Step 1: Construction of $\mathcal{E}_d$.}
Let the latent task variable be a binary string 
$z^\star = (z^\star_1, \ldots, z^\star_d) \in \mathcal{Z} = \{0,1\}^d$, 
sampled uniformly from $\rho_0$. The action space is $\mathcal{A} = \{0,1\}$, 
and the observation space is $\Omega = \{0,1,\bot\}$. The surface state 
is only the turn index, $s(h_t) = t$, so a belief-agnostic policy may 
know the depth but not the posterior induced by observed evidence.

At turn $t = 0$, action $a_0 = 0$ always reveals the first latent bit: 
$o_0 = z^\star_1$. Action $a_0 = 1$ returns $o_0 = \bot$.
For turns $t \geq 1$, the environment reveals the \emph{next} latent bit 
if and only if the agent chooses the \emph{previously revealed} bit:
\[
    o_t =
    \begin{cases}
        z^\star_{t+1}, & \text{if } a_t = z^\star_t, \\
        \bot,          & \text{otherwise.}
    \end{cases}
\]
The terminal reward is $R_T = 1$ if the agent has observed all $d$ latent 
bits by the end of the episode, and $R_T = 0$ otherwise. Since at most one 
bit is revealed per turn, any $\bot$ observation prevents success.

The key design principle is that the correct action at turn $t$ is 
determined by the observation $o_{t-1} = z^\star_t$ received at the 
previous turn. This observation is part of the history $h_t$ and therefore 
encoded in the belief $b_t$, but it is not a function of the surface state 
$s(h_t) = t$ alone.

\paragraph{Step 2: An epistemically adaptive policy succeeds with 
probability one.}
Define $\pi^\star$ as follows. At $t = 0$, choose $a_0 = 0$, which reveals 
$z^\star_1$. At each later turn $t \geq 1$, choose $a_t = z^\star_t$, 
the bit revealed as observation $o_{t-1}$ at the previous turn.

This action is accessible to $\pi^\star$: the value $z^\star_t = o_{t-1}$ 
is contained in the history $h_t = (a_0, o_0, \ldots, a_{t-1}, o_{t-1})$ 
and hence in the posterior belief $b_t$. An epistemically adaptive policy 
can therefore compute $a_t = z^\star_t$ from $b_t$.

By induction, after turn $t$, the agent has observed 
$z^\star_1, \ldots, z^\star_{t+1}$. After $d - 1$ turns of this form, 
the agent has observed the full latent string $z^\star$, so 
$\mathbb{P}(\mathrm{success} \mid \pi^\star) = 1$.

\paragraph{Step 3: Belief-agnostic policies cannot use the revealed bit.}
Let $\mu$ be any belief-agnostic policy. Since $s(h_t) = t$, the policy 
conditions only on the turn index: 
$\mu(a_t \mid h_t) = \mu_t(a_t)$ for some distribution 
$\mu_t \in \Delta(\{0,1\})$, independent of the observations received.

At $t = 0$, the best belief-agnostic policy chooses $a_0 = 0$, so we 
upper bound its probability of passing the first turn by $1$.

For any later turn $t \geq 1$, success requires $a_t = z^\star_t$. 
The value $z^\star_t$ was revealed as observation $o_{t-1}$, but $\mu$ 
does not condition on $o_{t-1}$. Conditional on success at all previous 
turns---which is required for $z^\star_t$ to have been revealed---the 
revealed bit $z^\star_t$ is uniformly distributed on $\{0,1\}$ by symmetry 
of the prior. Therefore, for any fixed action distribution $\mu_t$:
\begin{align}
    &\notag\mathbb{P}(a_t = z^\star_t \mid \mathrm{success\ up\ to\ } t-1,\, \mu)
    \\=&
    \sum_{a \in \{0,1\}} \mu_t(a)\,\mathbb{P}(z^\star_t = a)
    =
    \frac{\mu_t(0) + \mu_t(1)}{2}
    = \frac{1}{2}.
\end{align}
The equality holds regardless of the mixing weights $\mu_t(0), \mu_t(1)$ 
because the bit $z^\star_t$ is symmetric.

\paragraph{Step 4: Chaining the per-turn bounds.}
For the terminal reward $R_T = 1$, the agent must choose the correct 
action at every turn $t = 1, \ldots, d-1$ (turn $t = 0$ succeeds 
with probability $1$ under the best belief-agnostic policy). The 
events $\{a_t = z^\star_t\}$ are conditionally independent across turns 
given the latent string $z^\star$, because $\mu_t$ is a fixed distribution 
and $z^\star_1, \ldots, z^\star_d$ are independent bits. Therefore:
\[
    \mathbb{P}(\mathrm{success} \mid \mu)
    \leq
    1 \cdot \prod_{t=1}^{d-1} \frac{1}{2}
    =
    2^{1-d}.
\]
Combined with Part~1, this gives $\mathbb{P}(\mathrm{success} \mid \pi^\star) 
= 1$ versus $\mathbb{P}(\mathrm{success} \mid \mu) \leq 2^{1-d}$ for any 
belief-agnostic $\mu$, establishing the claimed separation.

\paragraph{Remark.}
This construction is related to the binary tree construction in 
\citep{zhang2025markovianreflectiveexplorationbayesadaptive} (Theorem~4.3), 
where an epistemically adaptive policy succeeds with probability~$1$ in a 
depth-$d$ tree while any Markovian policy succeeds with probability 
$1/2^{d-1}$. The key difference is the proof technique: their argument 
uses a visit-frequency induction on tree nodes under Markovian policies, 
while our argument uses the direct symmetry of the per-turn bit under 
belief-agnostic action selection. Both constructions achieve an 
exponential gap of $2^{\Omega(d)}$, but ours is more direct for 
the EDP setting because it does not require the full tree structure—only a sequential bit-reveal chain.
\end{proof}

% \begin{theorem}[Trajectory return is not sufficient for epistemic credit]
% \label{thm:trajectory-insufficiency}
% Let \(\mathcal{E}\) be a non-trivial EDP and fix a policy \(\pi\). Let \(G=\sum_{s=0}^{T-1}r_s\) denote the total episode return, and let \(A_{\mathrm{B}}^\pi(h_t,a_t)\) denote the Bayesian advantage at turn \(t\). Consider any trajectory-level estimator of the form \(\widehat{A}^{\mathrm{traj}}_t=\phi(G)\), where \(\phi\) is any measurable function of the total return.

% Then the best possible trajectory-level estimator satisfies
% \[
%     \inf_{\phi}
%     \mathbb{E}
%     \left[
%         \left(
%         \phi(G)-A_{\mathrm{B}}^\pi(h_t,a_t)
%         \right)^2
%     \right]
%     =
%     \mathbb{E}
%     \left[
%         \operatorname{Var}
%         \left(
%         A_{\mathrm{B}}^\pi(h_t,a_t)\mid G
%         \right)
%     \right].
% \]
% Consequently, if \(A_{\mathrm{B}}^\pi(h_t,a_t)\) is not measurable with respect to \(G\), then every trajectory-level estimator has strictly positive irreducible error:
% \[
%     \inf_{\phi}
%     \mathbb{E}
%     \left[
%         \left(
%         \phi(G)-A_{\mathrm{B}}^\pi(h_t,a_t)
%         \right)^2
%     \right]
%     > 0.
% \]
% \end{theorem}

\begin{theorem}[Spectral aliasing of trajectory-level credit]
\label{thm:spectral-aliasing}
Fix a policy \(\pi\) in an EDP and a turn \(t\). Let \(\omega\) denote a full sampled trajectory, let \(G(\omega)=\sum_{s=0}^{T-1}r_s(\omega)\) be its total return, and let
\[
    A(\omega)
    =
    A_{\mathrm{B}}^\pi(h_t(\omega),a_t(\omega))
\]
be the Bayesian advantage of the action taken at turn \(t\). Consider the class of trajectory-level estimators \(\widehat A_t^{\mathrm{traj}}=\phi(G)\), where \(\phi\) is any measurable function of the total return.

Let \(\mathcal{H}=L^2(\Omega,\mathbb{P}_\pi)\) and let \(\mathcal{H}_G=\{\phi(G):\phi\in L^2\}\) be the closed subspace of return-measurable functions. Let \(\Pi_G:\mathcal{H}\rightarrow \mathcal{H}_G\) be the orthogonal projection onto this subspace, i.e. \(\Pi_G A=\mathbb{E}[A\mid G]\). Then
\begin{align}
    \inf_{\phi}
    \mathbb{E}\left[
        \left(
        \phi(G)-A
        \right)^2
    \right]
    &\notag=
    \left\|
        (I-\Pi_G)A
    \right\|_2^2\\
    &\notag=
    \mathbb{E}\left[
        \operatorname{Var}(A\mid G)
    \right].
\end{align}
Consequently, if \(A_{\mathrm{B}}^\pi(h_t,a_t)\) is not measurable with respect to the total return \(G\), then any trajectory-level estimator has strictly positive irreducible statewise error.
\end{theorem}

\begin{proof}
We view advantage estimation as an approximation problem in the Hilbert space \(\mathcal{H}=L^2(\Omega,\mathbb{P}_\pi)\), equipped with inner product \(\langle f,g\rangle=\mathbb{E}[fg]\). The class of trajectory-level estimators \(\phi(G)\) is exactly the subspace \(\mathcal{H}_G\) of functions measurable with respect to the sigma-algebra generated by \(G\).

The conditional expectation \(\Pi_G A=\mathbb{E}[A\mid G]\) is the orthogonal projection of \(A\) onto \(\mathcal{H}_G\). Therefore, for any \(\phi(G)\in \mathcal{H}_G\), we have the orthogonal decomposition
\[
    \phi(G)-A
    =
    \underbrace{\phi(G)-\Pi_G A}_{\in \mathcal{H}_G}
    -
    \underbrace{(I-\Pi_G)A}_{\perp \mathcal{H}_G}.
\]
By Pythagoras,
\[
    \|\phi(G)-A\|_2^2
    =
    \|\phi(G)-\Pi_G A\|_2^2
    +
    \|(I-\Pi_G)A\|_2^2.
\]
The first term is minimized by choosing \(\phi^\star(G)=\Pi_G A=\mathbb{E}[A\mid G]\), leaving the irreducible residual \(\|(I-\Pi_G)A\|_2^2\). Since conditional expectation gives the best \(L^2\) predictor, this residual is
\begin{align}
    \|(I-\Pi_G)A\|_2^2
    &\notag=
    \mathbb{E}\left[
        \left(A-\mathbb{E}[A\mid G]\right)^2
    \right]\\
    &\notag=
    \mathbb{E}\left[
        \operatorname{Var}(A\mid G)
    \right].
\end{align}
If \(A\) is not measurable with respect to \(G\), then \(A\neq \Pi_GA\) on a set of positive probability, so \(\|(I-\Pi_G)A\|_2^2>0\). This proves the claim.
\end{proof}



% \begin{proof}
% Let \(A=A_{\mathrm{B}}^\pi(h_t,a_t)\). Since \(\widehat{A}^{\mathrm{traj}}_t=\phi(G)\) is measurable with respect to \(G\), we can decompose its squared error as
% \[
%     \mathbb{E}\left[(\phi(G)-A)^2\right]
%     =
%     \mathbb{E}
%     \left[
%         \left(
%         \phi(G)-\mathbb{E}[A\mid G]
%         \right)^2
%     \right]
%     +
%     \mathbb{E}
%     \left[
%         \operatorname{Var}(A\mid G)
%     \right].
% \]
% This is the standard orthogonal decomposition of squared error with respect to the sigma-algebra generated by \(G\). The first term is minimized by choosing \(\phi^\star(G)=\mathbb{E}[A\mid G]\), in which case it becomes zero. Therefore,
% \[
%     \inf_{\phi}
%     \mathbb{E}
%     \left[
%         \left(
%         \phi(G)-A
%         \right)^2
%     \right]
%     =
%     \mathbb{E}
%     \left[
%         \operatorname{Var}(A\mid G)
%     \right].
% \]
% If \(A\) is not measurable with respect to \(G\), then \(\operatorname{Var}(A\mid G)>0\) on a set of positive probability. Hence the expectation is strictly positive. This proves the result.
% \end{proof}
\begin{corollary}[Finite-dimensional view of trajectory-level aliasing]
\label{cor:finite-spectral-aliasing}
Assume the trajectory space is finite, \(\Omega=\{\omega_1,\ldots,\omega_n\}\), with probabilities \(p_i=\mathbb{P}_\pi(\omega_i)>0\). Let \(A_i=A(\omega_i)\), where \(A(\omega)=A_{\mathrm{B}}^\pi(h_t(\omega),a_t(\omega))\), and let \(G_i=G(\omega_i)\) be the total return of trajectory \(\omega_i\). Let \(\Pi_G\) denote the orthogonal projection onto the subspace of functions measurable with respect to \(G\), under the weighted inner product \(\langle f,g\rangle=\sum_{i=1}^n p_i f_i g_i\). Then
\[
    (\Pi_G A)_i
    =
    \mathbb{E}[A\mid G=G_i],
\]
i.e., \(\Pi_G\) replaces each trajectory-level advantage by the average Bayesian advantage among trajectories with the same total return. Moreover, \(\Pi_G\) is self-adjoint and idempotent, so its spectrum is contained in \(\{0,1\}\). The eigenspace with eigenvalue \(1\) consists of functions that are constant on each return level set, while the eigenspace with eigenvalue \(0\) consists of within-level-set contrasts whose conditional mean given \(G\) is zero.

Consequently, the irreducible error of any trajectory-level estimator \(\phi(G)\) is exactly the squared norm of the component of \(A\) lying outside the return-measurable subspace:
\begin{align}
    \inf_{\phi}
    \|\phi(G)-A\|_2^2
    &\notag=
    \|(I-\Pi_G)A\|_2^2\\
    &\notag=
    \sum_g
    \mathbb{P}(G=g)\operatorname{Var}(A\mid G=g).
\end{align}
Thus, if two positive-probability trajectories have the same total return but require different Bayesian credit at turn \(t\), then trajectory-level estimators cannot recover that within-return difference.
\end{corollary}

\begin{proof}
Since \(\Pi_G\) is the conditional expectation operator \(\mathbb{E}[\cdot\mid G]\), it maps each value \(A_i\) to the average of \(A\) over all trajectories with the same total return \(G_i\). Thus \((\Pi_G A)_i=\mathbb{E}[A\mid G=G_i]\).

The operator \(\Pi_G\) is an orthogonal projection in \(L^2(\Omega,\mathbb{P}_\pi)\), so \(\Pi_G^\ast=\Pi_G\) and \(\Pi_G^2=\Pi_G\). Therefore every eigenvalue \(\lambda\) satisfies \(\lambda^2=\lambda\), implying \(\lambda\in\{0,1\}\). The \(1\)-eigenspace is the range of \(\Pi_G\), namely the set of functions measurable with respect to \(G\), which are exactly the functions that are constant on each return level set. The \(0\)-eigenspace is the nullspace of \(\Pi_G\), namely functions \(f\) satisfying \(\mathbb{E}[f\mid G]=0\), which are precisely within-level-set contrasts.

By Theorem~\ref{thm:spectral-aliasing}, the best trajectory-level estimator is \(\Pi_G A=\mathbb{E}[A\mid G]\), and its residual is \((I-\Pi_G)A\). Therefore
\[
    \inf_{\phi}\|\phi(G)-A\|_2^2
    =
    \|(I-\Pi_G)A\|_2^2
    =
    \mathbb{E}[\operatorname{Var}(A\mid G)].
\]
In the finite case, this conditional variance expands as
\[
    \mathbb{E}[\operatorname{Var}(A\mid G)]
    =
    \sum_g
    \mathbb{P}(G=g)\operatorname{Var}(A\mid G=g).
\]
If \(A\) varies among trajectories in the same return level set, then at least one conditional variance term is positive, so the residual is strictly positive.
\end{proof}

\paragraph{Why this matters for epistemic credit}
Corollary~\ref{cor:finite-spectral-aliasing} makes the limitation of trajectory-level credit explicit. Total return partitions trajectories into return level sets, and any estimator \(\phi(G)\) can only assign credit that is constant within each level set. However, in an EDP, two trajectories can have the same final return while requiring different credit for an intermediate action because the action was taken under a different belief state. This within-return variation is exactly the component \((I-\Pi_G)A\), which lies outside the return-measurable subspace and is therefore invisible to any trajectory-level estimator. The failure is not finite-sample noise or poor optimization; it is an aliasing\footnote{We use the term ``aliasing'' by analogy to signal processing, where distinct high-frequency components can become indistinguishable under a coarse sampling operator~\citep{oppenheim1999discrete,shannon2006communication}. In our setting, the coarse observation is the aggregate return \(G\): trajectories with different belief-conditioned advantages can collapse to the same return level set.} problem caused by projecting belief-dependent credit onto a coarser signal. 
 

\paragraph{Connection to policy-gradient baselines}
This result is complementary to classical policy-gradient variance-reduction analyses such as Greensmith et al.~\citep{greensmith2004variance}. Those analyses study how baselines reduce variance when they condition on useful state information. Our result identifies the opposite obstruction in EDPs: if the estimator conditions only on aggregate return, then the projection subspace is too small to recover belief-dependent advantage components. In other words, the issue is not merely choosing a better scalar baseline; the estimator must expose information about the epistemic state.
\end{remark}

\begin{lemma}[Exact posterior in CSG]
\label{lemma:exact_posterior}
Let $z^\star$ be drawn uniformly from $\{1,\ldots,N\}$ and let
$\mathcal{C}_0 = \{1,\ldots,N\}$. Suppose the oracle responds
truthfully at every turn, i.e.,
$p(o_t \mid z^\star\!=\!z,\, a_t) = \mathbf{1}[z \in \mathcal{C}_{t+1}]$,
where $\mathcal{C}_{t+1} \subseteq \mathcal{C}_t$ is the set of
candidates consistent with observation $o_t$. Then for all
$t \geq 0$, the posterior is exactly uniform over the remaining
candidate set:
\[
  b_t(z) \;=\; \mathbb{P}(z^\star = z \mid h_t)
  \;=\;
  \frac{\mathbf{1}[z \in \mathcal{C}_t]}{|\mathcal{C}_t|}.
\]
\end{lemma}

\begin{proof} We prove this by induction. Notice that the base case holds by construction:
$b_0(z) = 1/N = 1/|\mathcal{C}_0|$ for all $z$.

For the inductive step, assume
$b_t(z) = \mathbf{1}[z \in \mathcal{C}_t] / |\mathcal{C}_t|$.
After action $a_t$ and observation $o_t$, Bayes' rule gives
\[
  b_{t+1}(z)
  = \frac{b_t(z)\; p(o_t \mid z,\, a_t)}
         {\sum_{z'} b_t(z')\; p(o_t \mid z',\, a_t)}.
\]
Substituting the inductive hypothesis and the binary likelihood:
\begin{align}
  b_{t+1}(z)
  &\notag= \frac{
      \frac{\mathbf{1}[z \in \mathcal{C}_t]}{|\mathcal{C}_t|}
      \cdot \mathbf{1}[z \in \mathcal{C}_{t+1}]
    }{
      \sum_{z' \in \mathcal{C}_t}
      \frac{1}{|\mathcal{C}_t|}
      \cdot \mathbf{1}[z' \in \mathcal{C}_{t+1}]
    }\\
  &\notag= \frac{\mathbf{1}[z \in \mathcal{C}_{t+1}]}
         {|\mathcal{C}_{t+1}|},
\end{align}
where the numerator uses
$\mathcal{C}_{t+1} \subseteq \mathcal{C}_t$, so
$\mathbf{1}[z \in \mathcal{C}_t] \cdot
 \mathbf{1}[z \in \mathcal{C}_{t+1}]
 = \mathbf{1}[z \in \mathcal{C}_{t+1}]$,
and the denominator sums to
$|\mathcal{C}_{t+1}|/|\mathcal{C}_t|$, which cancels
$|\mathcal{C}_t|$.
\end{proof}

\paragraph{Remark.}
The oracle in CGS is implemented via GPT-4o mini, so the observation
kernel $O(o_t \mid z^\star, h_t, a_t)$ is stochastic\footnote{\href{https://community.openai.com/t/why-does-openai-api-behave-randomly-with-temperature-0-and-top-p-0/934104}{https://community.openai.com/t/why-does-openai-api-behave-randomly-with-temperature-0-and-top-p-0/934104}} rather than
deterministic—even at temperature set to zero. However, the candidate-set update removes
inconsistent candidates given the \emph{realized} observation, so
$b_t$ is exact conditional on $o_t$; the stochasticity enters
through the EDP environment dynamics
(\Cref{sec:edp_theory}), not through belief-approximation error.
This makes the turn reward
$r_t = (|\mathcal{C}_t| - |\mathcal{C}_{t+1}|)/|\mathcal{C}_t|$
a direct measure of posterior contraction, computable in closed
form without learned reward models or Monte Carlo EIG
estimation~\citep{choudhury2025bed}.
 
\paragraph{Connection to BAMDPs and epistemic planning.}
    EDPs are best viewed as a task-level specialization of belief-state decision processes, rather than as a representationally separate alternative to POMDPs or BAMDPs. The distinction is one of modeling focus, but one that leverages the expressivity of LLMs that naturally  operate in exponentially complex language-based state-action spaces. BAMDPs typically maintain beliefs over unknown transition or reward parameters, or over latent MDP identities that are computationally expensive to track with LLMs, and use interaction to improve control under environment uncertainty. Epistemic planning methods such as Reflect-then-Plan similarly reason over model uncertainty to improve offline-to-online planning~\citep{jeong2025reflectthenplan}. In contrast, EDPs target information-seeking language agents operating through a fixed interaction interface: the environment need not be unknown in its dynamics, but the current task variable \(z^\star\) is hidden (See \Cref{ssec:edp_formulation}). The relevant belief is therefore a posterior over task hypotheses, \(p(z^\star \mid h_t)\), and the central question is whether an action acquires evidence that changes this posterior. Prior work~\citep{zhang2025markovianreflectiveexplorationbayesadaptive} is closest in spirit, using Bayes-adaptive RL to induce uncertainty-adaptive reflection over reasoning chains or MDP hypotheses~\citep{zhang2025markovianreflectiveexplorationbayesadaptive}; EDPs instead focus on external\footnote{In contrast to \citep{zhang2025markovianreflectiveexplorationbayesadaptive}, where adaptivity emerges “in-context” aka in reasoning chains within a turn, our focus in on adaptivity in an external sense i.e., across interaction turns.} query actions and observable belief updates \textit{across} interaction turns. Thus, the contribution of EDPs is not a new belief-state formalism in the abstract, but a lens that makes the credit signal for information-seeking actions explicit.
 



% \paragraph{Relation to BAMDPs and epistemic POMDPs.}
% EDPs share the belief-state perspective of POMDPs, BAMDPs, and recent epistemic planning formulations, but differ in the object of uncertainty and the role of the language policy. In a BAMDP, the agent maintains a posterior over unknown transition or reward parameters and explores to learn which MDP it is acting in; similarly, epistemic POMDP-style planning methods such as Reflect-then-Plan reason over model uncertainty to improve offline planning under uncertain dynamics~\citep{jeong2025reflectthenplan}. BARL is closest in spirit to our motivation, using Bayes-adaptive RL to induce uncertainty-adaptive reflection over reasoning or MDP hypotheses~\citep{zhang2025markovianreflectiveexplorationbayesadaptive}. In contrast, EDPs focus on interactive information seeking under a fixed query interface: the environment need not be unknown in its dynamics, but the task variable \(z^\star\) is hidden, and the policy must choose actions that elicit observations reducing uncertainty about that variable. Thus, the EDP belief is not primarily a posterior over environment dynamics, but a posterior over the current task hypothesis, \(p(z^\star \mid h_t)\). This distinction lets us assign credit directly to external query actions according to how they update the task belief, rather than treating exploration as internal reflection, model identification, or planning under uncertain dynamics.

% Theorem~\ref{thm:trajectory-insufficiency} shows that aggregate episode return is not, in general, a sufficient statistic for per-turn epistemic credit: whenever the Bayesian advantage varies across belief states that can share the same return, any trajectory-level estimator has irreducible statewise error.



\section{ECHO Training Dynamics}
\label{app:training_dynamics}


\paragraph{Agent Prompts} \Cref{fig:csg-prompts} (A) shows prompts used for training of LLM agents in CSG environment as well as evaluation.  \Cref{fig:csg-prompts} (B) shows the prompt used to instantiate the clue (or hint) provider model with GPT-4o-mini with near-deterministic inference with temperature ($T=0$).  The CSG uses two prompts with distinct roles.
The selector prompt presents the agent with its current belief state ---
the remaining candidate list $C_t$ and the full arm query history ---
and requires structured JSON output specifying an arm index and a property question.
No chain-of-thought or reasoning scaffold is provided; the agent must select
informative actions directly from the belief state encoding.
The oracle prompt is given to a fixed \texttt{gpt-4o-mini} clue holder
at temperature zero, which receives the secret number and the agent's question
and returns a single deterministic sentence answering only the queried property.
The oracle is explicitly instructed not to volunteer information beyond what was asked
and to deflect vague questions, enforcing strict information separation across
the $K{=}5$ arms.
Repeat questions to the same arm incur a penalty of $r_t = -0.1$
rather than an oracle call, discouraging redundant queries.
This design ensures that all information gain flows exclusively through
the structured arm--question interface, making the per-turn reward signal
a clean measure of epistemic progress.


\input{csg_agent_prompt_clean}
\paragraph{Training dynamics and mechanism analysis.}
Figure~\ref{fig:clue_steps_additional_metrics} along with main metrics in \Cref{fig:clue_steps_four_panel_training_main} (\Cref{sec:results}) decomposes training into \textbf{outcome, mechanism, and optimization-health metrics}. The main pattern is that \textsc{ECHO} improves not only final task success, but also the intermediate epistemic behavior predicted by the EDP analysis. Episode resolve rate rises sharply, while mean eliminated candidates per turn increases, indicating that the policy learns to ask questions that contract the current candidate set more effectively. At the same time, redundancy falls to near zero, showing that the policy is not simply repeating high-probability questions, but is using the evolving history to avoid already-exhausted information. This supports the central motivation for \textbf{turn-level epistemic credit}: when the reward signal is assigned at the decision depth where an action changes the posterior, the policy receives direct supervision about which questions are useful under the current belief state.

\paragraph{Specialization without collapse.}
The diagnostic metrics further suggest that \textsc{ECHO}'s gains reflect healthy specialization rather than mode collapse. Entropy decreases during training, showing that the policy becomes more decisive, but diversity remains near saturation, indicating that it does not collapse onto a small set of repeated questions. Mechanistically, the policy appears to learn strong \textbf{within-step information gathering} before fully mastering higher-level arm selection: candidate elimination and efficiency improve substantially, while arm-match remains comparatively weaker. This is consistent with the task structure, where the policy must learn a joint action consisting of both an arm choice and a natural-language question. In EDP terms, \textsc{ECHO} first improves posterior-sensitive local action quality---asking useful questions for the current candidate set---while leaving additional headroom for learning better long-horizon coordination over which information source to query.

\begin{figure*}[t]
    \centering

    % ---------- Row 1 ----------
    \includegraphics[width=0.24\linewidth]{figures/rl_per_step/clue_normalized_reward.pdf}
    \hfill
    \includegraphics[width=0.24\linewidth]{figures/rl_per_step/reward_std.pdf}
    \hfill
    \includegraphics[width=0.24\linewidth]{figures/rl_per_step/clue_diversity_rate.pdf}
    \hfill
    \includegraphics[width=0.24\linewidth]{figures/rl_per_step/clue_resolve_rate.pdf}

    \vspace{0.8em}

    % ---------- Row 2 ----------
    \includegraphics[width=0.24\linewidth]{figures/rl_per_step/clue_mean_episode_length.pdf}
    \hfill
    \includegraphics[width=0.24\linewidth]{figures/rl_per_step/clue_eliminated_mean.pdf}
    \hfill
    \includegraphics[width=0.24\linewidth]{figures/rl_per_step/clue_efficiency_bonus_mean.pdf}
    \hfill
    \includegraphics[width=0.24\linewidth]{figures/rl_per_step/clue_turns_saved_mean.pdf}

    \vspace{0.8em}

    % ---------- Row 3 ----------
    \includegraphics[width=0.31\linewidth]{figures/rl_per_step/loss.pdf}
    \hfill
    \includegraphics[width=0.31\linewidth]{figures/rl_per_step/grad_norm.pdf}
    \hfill
    \includegraphics[width=0.31\linewidth]{figures/rl_per_step/entropy.pdf}

    \caption{Additional learning dynamics of RL baselines and diagnostic metrics across training steps in the Clue Selector Game (CSG) environment.}
    \label{fig:clue_steps_additional_metrics}
\end{figure*}

\begin{table*}[t]
\centering
\scriptsize
\setlength{\tabcolsep}{6pt}
\renewcommand{\arraystretch}{1.12}
\begin{tabular}{lrrrr}
\toprule
\textbf{Metric} & \textbf{\textsc{ECHO}} & \textbf{GRPO} & \textbf{$\Delta$} & \textbf{Paired $t$-test $p$} \\
\midrule
Reward mean & 0.598 & 0.342 & +0.257 & $7.55{\times}10^{-4}$ \\
Normalized reward & 0.419 & 0.264 & +0.155 & $8.43{\times}10^{-4}$ \\
Episode resolve rate & 66.1\% & 17.1\% & +48.9 pts & $2.42{\times}10^{-4}$ \\
Mean eliminated candidates & 12.924 & 10.539 & +2.385 & $1.29{\times}10^{-3}$ \\
Redundancy rate $\downarrow$ & 0.26\% & 2.31\% & $-$2.04 pts & $1.55{\times}10^{-2}$ \\
Mean episode length $\downarrow$ & 7.724 & 9.136 & $-$1.412 & $9.88{\times}10^{-4}$ \\
Turn resolve rate & 8.66\% & 1.94\% & +6.72 pts & $3.69{\times}10^{-4}$ \\
Turns saved & 0.304 & 0.066 & +0.238 & $4.08{\times}10^{-4}$ \\
Entropy $\downarrow$ & 0.171 & 0.663 & $-$0.492 & $2.44{\times}10^{-3}$ \\
\bottomrule
\end{tabular}
\caption{
Final-window comparison between \textsc{ECHO} and episodic return baseline (GRPO). Values are seed-level means over the last 100 training steps, averaged across three seeds. \(p\)-values are from paired seed-level \(t\)-tests. Lower is better for redundancy rate, mean episode length, and entropy.
}
\label{tab:echo_grpo_final_window_stats}
\end{table*}

\paragraph{Seed-level significance.}
We compute final-window means over the last 100 training steps for each seed and compare \textsc{ECHO} against episodic GRPO using paired seed-level tests. \Cref{tab:echo_grpo_final_window_stats} shows the results. Across all three seeds, \textsc{ECHO} improves the primary outcome and mechanism metrics: episode resolve rate increases from \(17.1\%\) to \(66.1\%\), mean eliminated candidates rises from \(10.54\) to \(12.92\), turn resolve rate increases from \(1.94\%\) to \(8.66\%\), and mean episode length decreases from \(9.14\) to \(7.72\). Paired \(t\)-tests are significant for these metrics despite the small number of seeds, reflecting consistent improvements across runs. Since \(n=3\), we interpret these results primarily through effect size and cross-seed consistency rather than relying solely on significance tests.
\section{Additional Results}
\label{app:binary_search_echo}


\subsection{Main Evaluation Metric Definitions}
\label{ssec:main_task_epistemic_metrics_full}
We compute task-level and process-level metrics from CSG episode logs over 10 maximum turns. Unless otherwise stated, turn-level metrics are averaged over turns within a run, and model-level values report mean and standard error across runs. Additional metrics that we consider and those that CSG RL training environment provides are noted in \Cref{tab:metric_definitions} (for training-related metrics) and \Cref{tab:arm_metric_definitions} (for arm-routing specific metrics). 

\begin{itemize}
    \item \textbf{Resolve} measures final task success: the fraction of episodes in which the agent identifies the secret within the turn budget.

    \item \textbf{Zero} measures zero-elimination behavior: the fraction of non-redundant turns that eliminate no candidates from the current candidate set. Redundant turns are excluded so that Zero captures new but uninformative questions.

    \item \textbf{Qual.} measures candidate-elimination quality relative to the current belief state. It normalizes the number of eliminated candidates by an ideal binary split of the current candidate set and caps the value at \(1\). Thus, high Qual. indicates questions that efficiently partition the remaining candidate set.

    \item \textbf{Ground} measures grounded information seeking: the fraction of turns that are both non-redundant and sufficiently informative under the current candidate set. In our implementation, a turn is grounded when its elimination quality is at least \(0.5\).

    \item \textbf{LateCalib} measures late-stage redundancy in narrow belief states. It is computed on turns with small remaining candidate sets late in the episode and records whether the action was redundant. Lower values indicate fewer wasted turns after the belief state has already narrowed.

    \item \textbf{Reason\%} measures visible reasoning text outside the required structured action format. It is used only as a diagnostic for whether a policy relies on explicit natural-language rationales rather than direct action selection.

    \item \textbf{ZeroEvents} counts zero-elimination non-redundant turns that have a following turn, and serves as the event denominator for recovery metrics.
 
\item \textbf{Rec1} measures weak immediate recovery after a Zero event: the fraction of Zero events after which the next turn is non-redundant and eliminates at least one candidate. It captures whether the agent resumes gaining any information at all.

\item \textbf{GRecover} measures grounded immediate recovery after a Zero event: the fraction of Zero events after which the next turn is grounded, meaning non-redundant and sufficiently informative under the current candidate set. This is stricter than Rec1 because the next turn must not merely eliminate something; it must pass the groundedness threshold.
 


    \item \textbf{RecQ} measures the elimination quality of the immediate next turn after a Zero event.

    \item \textbf{NextZero} measures whether a Zero event is followed by another zero-elimination non-redundant turn. Lower is better.

    \item \textbf{NextBad} measures whether a Zero event is followed by either another Zero event or a redundant turn. Lower is better.

    \item \textbf{TTR} measures the number of turns until the next grounded turn after a Zero event, averaged over cases where grounded recovery occurs. Lower is better.

    \item \textbf{TTRsucc} measures whether a grounded recovery occurs at any later turn after a Zero event before the episode ends.

    \item \textbf{ResAfterZero} measures episode-level robustness after local failure: among episodes containing at least one Zero event, it reports the fraction that still resolve.
\end{itemize}
 

\subsection{From Credit Signal to Strategy: The Emergence of Binary Search}

\paragraph{Oracle response distributions reveal emergent strategy.}
To characterize policy behavior beyond resolve rate, we analyze the distribution of oracle hint types received across evaluation episodes in CSG. This is shown in Table~\ref{tab:csg_hint_distribution_all}. For each turn, we classify the oracle response into coarse categories such as range, parity, divisibility, primality, square, digit properties, between, uninformative, and redundant, using keyword matching against the hint string. We also report the \emph{zero-elimination rate}: the fraction of non-redundant turns where the selected question eliminated no candidates. This captures cases where the agent asked a valid question, but the question was unhelpful under the current belief state. Table~\ref{tab:csg_hint_distribution_all} shows that \textsc{ECHO} has the lowest zero-elimination rate among evaluated policies (\(18.6\%\)), substantially below trajectory-level GRPO (Episodic Return) (\(39.8\%\)), RLOO per-turn (\(40.2\%\)), and prompting baselines such as Qwen2.5-1.5B + CoT (\(38.4\%\)) and full ReAct (\(49.2\%\)). Its hint distribution is also highly concentrated on \textit{range-based} clues (\(60.7\%\)), with low redundancy (\(0.6\%\)) and little reliance on digit-property or miscellaneous hints. This suggests that \textsc{ECHO} does not merely improve final resolution; it changes the mechanism of interaction toward consistently informative, belief-sensitive queries.

\paragraph{Bisection as a special case of the EDP objective.}
The emergence of binary search in \textsc{ECHO} is a consequence of the CSG environment, not a universal property of epistemically adaptive policies. In CSG, the initial prior is uniform, \(\rho_0=\mathrm{Unif}(\{1,\ldots,N\})\), and deterministic oracle filtering preserves a uniform posterior over the remaining candidates, \(b_t=\mathrm{Unif}(\mathcal{C}_t)\). For a binary question that partitions \(\mathcal{C}_t\) into fractions \(p\) and \(1-p\), the expected fractional elimination is
\[
    \mathbb{E}_{z^\star\sim b_t}[u(b_t,a_t,o_t)] = 2p(1-p),
\]
which is maximized at \(p=1/2\). Thus, under uniform beliefs, deterministic binary observations, and equal action costs, the EDP objective recovers bisection as the information-theoretically optimal strategy.

In general EDPs, these conditions need not hold. If the belief is non-uniform, as in diagnosis or retrieval settings, the optimal query should account for posterior mass rather than split the hypothesis set evenly. If observations are graded, noisy, or multi-valued, the optimal action depends on the task-specific observation model and belief geometry. The Bayesian value \(Q^\pi_{\mathrm{B}}(h_t,a_t)=\mathbb{E}_{z^\star\sim b_t}[Q^\pi(h_t,a_t\mid z^\star)]\) remains the relevant object, but its maximizer is not necessarily a bisection question. Thus, bisection is a special solution for the uniform-prior, deterministic-observation regime of CSGG, while the EDP framework and \textsc{ECHO}'s posterior-sensitive turn-level credit are more general.

\paragraph{Binary-search-like strategy in CSG} In our experiments in CSG, the hint distribution clearly explains why \textsc{ECHO} converges toward a binary-search-like strategy rather than the digit-property-heavy strategy used by the greedy oracle. In the CSG instantiation of the EDP (\Cref{sec:experiments}), the dominant evidential utility is fractional candidate elimination,
\(u(b_t,a_t,o_t)=|\mathcal{C}_t\setminus\mathcal{C}_{t+1}|/|\mathcal{C}_t|\), which measures posterior contraction relative to the current belief size. Since the secret is sampled uniformly and oracle observations only eliminate inconsistent candidates, \(b_t=\mathrm{Unif}(\mathcal{C}_t)\). For a binary question whose ``yes'' set occupies fraction \(p\) of the current candidate set, the expected fractional elimination under this uniform belief is \(2p(1-p)\), maximized at \(p=1/2\). Range-bisection questions can therefore achieve near-maximal expected utility by splitting \(\mathcal{C}_t\) roughly in half. By contrast, digit-property questions such as last digit or digit sum can be highly discriminative for some secrets, but often correspond to imbalanced predicates under the uniform posterior, giving lower expected utility on average. Thus, \textbf{\textsc{ECHO}} learns the policy that is optimal for an agent uncertain over \(z^\star\): ask questions that maximize expected posterior collapse. The \textbf{greedy oracle} baseline, in contrast, effectively conditions on the true latent variable and can exploit digit properties with high realized elimination for that specific secret, explaining its \(50.7\%\) combined use of last-digit and digit-sum hints versus \textsc{ECHO}'s \(60.7\%\) use of range clues.

\input{hint_classification_table_baselines}

\paragraph{Arm routing versus belief-conditioned information seeking.}

\begin{table}[t]
\centering
\scriptsize
\setlength{\tabcolsep}{4pt}
\renewcommand{\arraystretch}{1.08}
\begin{tabular}{p{0.18\linewidth}p{0.74\linewidth}}
\toprule
\textbf{Metric} & \textbf{Meaning} \\
\midrule
% \textbf{Res.} & Episode-level resolve rate: the fraction of games in which the policy uniquely identifies the secret number. \\

% \textbf{Zero} & Fraction of non-redundant turns that eliminate no candidates, measuring uninformative questions that fail to update the belief state. \\

% \textbf{Qual.} & Mean elimination quality: the fraction of an ideal binary split achieved by the selected question. \\

% \textbf{Ground} & Grounded-turn rate: the fraction of turns that are non-redundant and sufficiently informative relative to the current candidate set. \\

\textbf{ArmEnt} & Normalized entropy of arm choices, measuring how evenly the policy distributes actions across the available arms. \\

\textbf{DomArm} & Dominant-arm share: the fraction of turns assigned to the policy's most frequently selected arm. \\

\textbf{StaticComp} & Static arm compatibility: whether the question type is semantically compatible with the fixed property family assigned to the selected arm. \\

\textbf{ArmMatch} & Logged arm-match rate: agreement between the selected arm and the environment's target arm label. \\

\textbf{SameArm} & Fraction of turns in which the policy selects the same arm as in the previous turn. \\

\textbf{ZeroSame} & Zero-elimination rate conditioned on selecting the same arm as in the previous turn. \\

\textbf{ZeroSwitch} & Zero-elimination rate conditioned on switching to a different arm from the previous turn. \\

\textbf{QualSame} & Mean elimination quality conditioned on selecting the same arm as in the previous turn. \\

\textbf{QualSwitch} & Mean elimination quality conditioned on switching to a different arm from the previous turn. \\
\bottomrule
\end{tabular}
\caption{Definitions of arm-routing and belief-state diagnostic metrics.}
\label{tab:arm_metric_definitions}
\end{table}
Table~\ref{tab:arm_routing_belief_diagnostics} tests whether \textsc{ECHO}'s gains are explained by raw arm exploration, static arm matching, or belief-conditioned information gain. The results rule out a simple exploration account: \textsc{ECHO} has lower arm entropy than several weaker baselines such as GRPO, RLOO, and Qwen2.5-1.5B base model, yet achieves substantially higher resolve rate. Instead, \textsc{ECHO} exhibits the strongest learned-policy combination of low zero-elimination rate, high elimination quality, high groundedness, and high semantic compatibility between selected arms and question types. Raw ArmMatch alone is not sufficient to explain performance: \textsc{ECHO}'s ArmMatch is modest, but its selected arm-question pairs eliminate candidates much more reliably than GRPO or base-model baselines. The same-arm diagnostics further show that arm repetition is not inherently harmful. While frontier models often suffer when reusing an arm, \textsc{ECHO}'s same-arm turns remain highly informative, with lower zero-elimination and higher elimination quality than its switched-arm turns. This suggests that \textsc{ECHO} learns to reuse and route arms in a belief-conditioned manner rather than merely maximizing static arm labels.

\begin{table*}[t]
\centering
\scriptsize
\setlength{\tabcolsep}{4.2pt}
\renewcommand{\arraystretch}{1.08}
\resizebox{\linewidth}{!}{
\begin{tabular}{lccccccccccccc}
\toprule
\textbf{Model} &
\textbf{Res. $\uparrow$} &
\textbf{Zero $\downarrow$} &
\textbf{Qual. $\uparrow$} &
\textbf{Ground $\uparrow$} &
\textbf{ArmEnt} &
\textbf{DomArm} &
\textbf{StaticComp $\uparrow$} &
\textbf{ArmMatch $\uparrow$} &
\textbf{SameArm} &
\textbf{ZeroSame $\downarrow$} &
\textbf{ZeroSwitch $\downarrow$} &
\textbf{QualSame $\uparrow$} &
\textbf{QualSwitch $\uparrow$} \\
\midrule

% \multicolumn{14}{l}{\textit{Reference policy}} \\
% \hdashline
\rowcolor{oracleshade}
Greedy oracle (reference)
& 0.862 & 0.296 & 0.700 & 0.704 & 0.839 & 0.507 & 1.000 & 1.000 & 0.468 & 0.393 & 0.423 & 0.603 & 0.570 \\

\midrule
% \multicolumn{14}{l}{\textit{Ours}} \\
% \hdashline
\rowcolor{echoshade}
\textbf{\textsc{ECHO} (Ours)}
& \textbf{0.453} & \textbf{0.186} & \textbf{0.670} & \textbf{0.718} & 0.770 & 0.536 & \textbf{0.581} & 0.183 & 0.565 & \textbf{0.172} & \textbf{0.257} & \textbf{0.670} & \textbf{0.582} \\

\midrule
\multicolumn{14}{l}{\textit{Frontier baselines}} \\
\hdashline
Claude Sonnet 4.6
& 0.431 & 0.254 & 0.701 & 0.738 & 0.976 & 0.309 & 0.345 & 0.146 & 0.112 & 0.645 & 0.239 & 0.313 & 0.721 \\
GPT-4o
& 0.258 & 0.298 & 0.625 & 0.662 & 0.996 & 0.230 & 0.391 & 0.172 & 0.215 & 0.596 & 0.261 & 0.318 & 0.654 \\

\midrule
\multicolumn{14}{l}{\textit{RL-trained Qwen2.5-1.5B baselines}} \\
\hdashline
RLOO per-turn
& 0.138 & 0.402 & 0.446 & 0.445 & 0.960 & 0.276 & 0.311 & 0.135 & 0.319 & 0.406 & 0.458 & 0.421 & 0.390 \\
Episodic Return (GRPO)
& 0.118 & 0.398 & 0.423 & 0.419 & 0.944 & 0.292 & 0.288 & 0.141 & 0.378 & 0.393 & 0.469 & 0.409 & 0.354 \\

\midrule
\multicolumn{14}{l}{\textit{Base and prompting baselines}} \\
\hdashline
Qwen2.5-1.5B
& 0.080 & 0.386 & 0.352 & 0.340 & 0.944 & 0.319 & 0.238 & 0.138 & 0.391 & 0.401 & 0.446 & 0.300 & 0.287 \\
Qwen2.5-1.5B + ReAct
& 0.044 & 0.492 & 0.312 & 0.316 & 0.826 & 0.537 & 0.313 & 0.189 & 0.552 & 0.525 & 0.564 & 0.237 & 0.265 \\

\bottomrule
\end{tabular}
}
\caption{
\textbf{Arm-routing and belief-state diagnostics in the Clue Selector Game}. See \Cref{tab:arm_metric_definitions} for of arm-routing and belief-state diagnostic metrics in CSG.}
\label{tab:arm_routing_belief_diagnostics}
\end{table*}


\input{results_recovery_zero_info_events}
\input{results_recovery_bad_events}
\input{results_recovery_low_info_events}
\input{results_recovery_late_bad_events}
\iffalse

\section{Theoretical Analysis}
\label{sec:theory}






We develop a general theory of online policy gradient optimization 
in environments where reward accrues through sequential information 
acquisition from multiple independent sources. We formalize such 
environments as \emph{Information Acquisition MDPs}, derive conditions 
under which the policy gradient objective recovers Gittins-optimal 
arm selection, and bound the suboptimality gap in terms of advantage 
estimation quality. The Clue Selector Game is one instantiation of 
this framework; the results apply broadly to any environment 
satisfying the conditions below.

%─────────────────────────────────────────────────────────────────
\subsection{Information Acquisition MDPs}
%─────────────────────────────────────────────────────────────────

\begin{definition}[Information Acquisition MDP]
\label{def:ia-mdp}
An \emph{Information Acquisition MDP} (IA-MDP) is a tuple
$\mathcal{M} = (\mathcal{H}, \mathcal{A}, f, r, \rho_0, \gamma)$, 
where:
\begin{itemize}
    \item $\mathcal{H}$ is the space of \emph{interaction histories}
          $h_t = (a_0, \omega_0, \ldots, a_{t-1}, \omega_{t-1})$,
          where $a_s \in \mathcal{A}$ is an action taken at step $s$
          and $\omega_s \in \Omega$ is an \emph{observation} 
          received from the environment.
    \item $\mathcal{A} = \bigcup_{k=1}^{K} \mathcal{A}_k$ is the 
          action space, partitioned into $K$ \emph{arm-specific} 
          action sets $\{\mathcal{A}_k\}_{k=1}^K$.
    \item $f: \mathcal{H} \times \mathcal{A} \times \Omega 
          \rightarrow \mathcal{H}$ is the deterministic history 
          transition: $f(h_t, a_t, \omega_t) = h_t \oplus 
          (a_t, \omega_t)$.
    \item $r: \mathcal{H} \times \mathcal{A} \rightarrow \mathbb{R}$ 
          is the per-step reward measuring information gained.
    \item $\rho_0 \in \Delta(\mathcal{X})$ is the prior distribution 
          over a latent variable $x^* \in \mathcal{X}$ that the 
          agent seeks to identify.
    \item $\gamma \in (0,1)$ is a discount factor.
\end{itemize}
The \emph{belief state} at step $t$ is the posterior 
$b_t = P(x^* | h_t) \in \Delta(\mathcal{X})$, which is a 
sufficient statistic for the history $h_t$ with respect to $x^*$.
\end{definition}

\section{Notes on BARL~\citep{zhang2025markovianreflectiveexplorationbayesadaptive}}

Let us work through their method, proofs, and gaps carefully. I think the right framing for the Clue selector RL environment is to study the problem as a Bayesian RL problem, within the context of decision making under uncertainty and specifically, the uncertainty—adaptive nature of intelligent policies—which we want our LLMs to possess. So, instead of starting with the Gittins Index angle of explicitly learning a Gittins-optimal policy e which can cause challenges in the assumptions leading to the RL formulation (independence of arms and the entire Bandit machinery the connection with RL for text language does not exist at the moment), we could simply use the Bayesian RL formulation from BARL and their proof techniques. Our current empirical results are already good, so we just need to write up our multiturn RL (GRPO variant) baseline as a novel Bayesian RL policy. So, the correct setup would be either an epistemic POMDP or the information acquisition MDPs. 


\subsection{Problem Formulation}

\paragraph{LLM Information Acquisition.}
We formulate our setting as a finite-horizon MDP
\[
\mathcal{M} = (\mathcal{S}, \mathcal{A}, T, r_{\mathcal{M}}),
\]
where $\mathcal{S}$ is the state space, $\mathcal{A}$ is the action space,
$T$ is the horizon, and $r_{\mathcal{M}}(s,a)$ is the reward function.
The initial state $s_0$ is the task prompt, which specifies the information
acquisition problem. At each step $t$, the action $a_t$ corresponds to an
information-seeking decision, such as selecting a clue source, querying an
arm, or asking a question. The environment then returns an observation
$\omega_t$, and the state transition is deterministic in the interaction
history:
\[
s_{t+1} = s_t \oplus (a_t,\omega_t).
\]
Equivalently, we may write the history at time $t$ as
\[
h_t = (s_0,a_0,\omega_0,r_0,\ldots,a_{t-1},\omega_{t-1},r_{t-1},s_t),
\]
where $s(h_t)=s_t$ denotes the current state induced by the history.

We are interested in a training--deployment procedure. During training, the
agent only receives partial information about the true environment
\[
\mathcal{M}^* := (\mathcal{S},\mathcal{A},r,T)
\]
through the training data $\mathcal{D}$. During deployment, the learned policy
is frozen and deployed in $\mathcal{M}^*$. Importantly, $\mathcal{D}$ does not
uniquely identify $\mathcal{M}^*$, inducing epistemic uncertainty about the
underlying environment. Formally, a prior distribution over MDPs $p(\mathcal{M})$
and data $\mathcal{D}$ define a posterior
\[
p(\mathcal{M}\mid \mathcal{D})
\propto
p(\mathcal{M})p(\mathcal{D}\mid \mathcal{M}).
\]
The objective for a policy $\pi$ is therefore to maximize return in expectation
over MDPs sampled from this posterior:
\[
J(\pi)
=
\mathbb{E}_{\mathcal{M}\sim p(\mathcal{M}\mid \mathcal{D})}
\left[
J_{\mathcal{M}}(\pi)
\right],
\]
where
\[
J_{\mathcal{M}}(\pi)
=
\mathbb{E}_{s_0,\pi}
\left[
\sum_{t=0}^{T-1}
r_{\mathcal{M}}(s_t,a_t)
\right].
\]
This objective trains a policy that maximizes its Bayes-expected return under
the posterior uncertainty induced by the data.

\paragraph{Progress Reward.}
Prior work on LLM reasoning often uses an outcome-level reward verifier that
checks whether the final state $s_T$ contains the correct answer. In our setting,
the final outcome corresponds to identifying the hidden target, denoted
$x^*$, or reducing the candidate set sufficiently to recover it. A sparse
terminal verifier can therefore be written as
\[
\mathrm{verifier}(s_T,x^*)
\in \{0,1\}.
\]
However, this sparse reward does not expose which information-acquisition
actions were useful. We therefore use a progress reward that measures how much
the agent reduces uncertainty about the hidden target after taking action
$a_t$ and observing $\omega_t$.

Let
\[
b_t(x) := \mathbb{P}(x^*=x\mid h_t,\mathcal{D})
\]
denote the belief over hidden targets after history $h_t$. Let $U(b_t)$ be an
uncertainty functional, such as entropy, support size, or log candidate-set
size. We define the progress reward as
\[
r(s_t,a_t,\omega_t)
=
U(b_t)-U(b_{t+1}).
\]
In the Clue Selector Game with a uniform belief over the surviving candidate
set $\mathcal{C}_t$, a natural choice is
\[
U(b_t)=\log |\mathcal{C}_t|,
\]
which gives the reward
\[
r(s_t,a_t,\omega_t)
=
\log |\mathcal{C}_t|
-
\log |\mathcal{C}_{t+1}|.
\]
This reward is dense and directly measures the information gained by the
current action. The corresponding Q-value is
\[
Q^{\pi}_{\mathcal{M}}(h_t,a_t)
=
\mathbb{E}_{\pi}
\left[
\sum_{t'=t}^{T-1}
r_{\mathcal{M}}(s_{t'},a_{t'},\omega_{t'})
+
\mathrm{verifier}(s_T,x^*)
\right].
\]

\paragraph{Adaptive Information Acquisition.}
In our setting, the agent's decision should depend not only on the current
surface state, but also on the belief induced by the full interaction history.
For example, two histories may have the same turn number or the same apparent
state representation, but correspond to different surviving candidate sets
$\mathcal{C}_t$. The optimal next query may therefore differ across these two
histories.

This motivates the following definition.

\begin{definition}[Uncertainty-Adaptive Policy]
Define the belief as the posterior over MDPs
\[
b(h_t)(\mathcal{M})
:=
p(\mathcal{M}\mid h_t,\mathcal{D}).
\]
A policy $\pi$ is uncertainty-adaptive if there exists a measurable mapping
\[
\mu : \mathcal{S}\times \mathcal{P}(\mathfrak{M}) \rightarrow \Delta(\mathcal{A})
\]
such that for all histories $h_t$ and actions $a_t$,
\[
\pi(a_t\mid h_t)
=
\mu(a_t\mid s(h_t),b(h_t)).
\]
If, in addition, there exist histories $h_t$ and $h'_t$ such that
\[
s(h_t)=s(h'_t)
\]
but
\[
\mu(a_t\mid s(h_t),b(h_t))
\neq
\mu(a_t\mid s(h'_t),b(h'_t)),
\]
then $\pi$ is strictly uncertainty-adaptive, or equivalently, not Markovian
with respect to the surface state alone.
\end{definition}

From this perspective, effective clue selection is the behavioral signature of
an uncertainty-adaptive policy which by definition is non-Markovian in nature. In the general case, this is what is expected for an optimal reasoning agent in BARL~\citep{zhang2025markovianreflectiveexplorationbayesadaptive}, where the agent can be incentivized to “seek-out” certain states or “revisit” them in order to make different decisions in the future relative to what it did in the past. The adaptability is crucial since it allows the agent to \textit{backtrack} on its past decisions, but importantly by incentivizing it to arrive at identical states, where it can revisit its prior actions. 


As the agent receives observations, its belief
over the hidden target changes, and the value of each information source changes
accordingly. A purely state-only policy—like the conventional RL situation—that ignores this evolving belief would
choose the same action distribution whenever it encounters the same surface
state. In contrast, an uncertainty-adaptive policy conditions on the posterior
belief and can rationally switch to a different clue source when the current
belief makes that source more informative.



\subsection{Their Method (BARL) --- What It Actually Does}

The policy gradient is:
\[
\nabla_\theta J
=
\mathbb{E}_{s_0,\pi_\theta}
\sum_{t=0}^{T-1}
\nabla_\theta \log \pi_\theta(a_t \mid h_t)
\cdot
\mathbb{E}_{\mathcal{M}\sim p(\mathcal{M}\mid D,h_t)}
\left[
Q^{\pi_\theta}_{\mathcal{M}}(h_t,a_t)
\right].
\]

The posterior-weighted Q-value decomposes into three terms (Eq. 5.4):
\[
\mathbb{E}_{\mathcal{M}}
\left[
Q^{\pi_\theta}_{\mathcal{M}}(h_t,a_t)
\right]
=
\sum_{i=0}^{|\mathcal{M}|}
Q^{\pi_\theta}_{\mathcal{M}_i}(h_t,a_t)
\cdot
\underbrace{
\pi_\theta
\left(
y^{\mathcal{M}_i}_{s_0}
\mid
s_t + \texttt{</think>}
\right)
}_{\text{model plausibility}}
\cdot
\underbrace{
\prod_{t'=0}^{t-1}
\exp\left(
-\beta
\left|
r_{t'} - r_{\mathcal{M}_i}(s_{t'},a_{t'})
\right|
\right)
}_{\text{reward consistency}}.
\]

So concretely what BARL does is:

\begin{itemize}
    \item Sample $|\mathcal{M}|$ CoT rollouts from $\pi_\theta$ on prompt $s_0$.
    \item Each rollout defines a hypothesis MDP $\mathcal{M}_i$ --- the ``world'' consistent with that rollout's answer.
    \item For each step $t$ in each rollout, compute the posterior-weighted Q-value using all hypothesis MDPs.
    \item Weight each hypothesis by (a) how plausible the model finds its answer and (b) how consistent past rewards are with that hypothesis.
    \item Use this weighted Q-value as the advantage signal in policy gradient.
\end{itemize}

The key insight of BARL: instead of using a single reward signal (like GRPO uses the scalar reward per completion), \textbf{BARL uses a distribution over possible reward functions weighted by belief. This forces the policy to track uncertainty across hypotheses rather than committing to a single reward signal.}

\subsection{Where BARL Is Lacking --- Five Concrete Gaps}

\subsubsection{Gap 1 --- The Hypothesis Space Is Implicit and Unstructured}

BARL forms hypotheses $\{\mathcal{M}_i\}$ by sampling CoT rollouts and using each rollout's final answer as the hypothesis. This means the hypothesis space is determined by what the model happens to generate --- there is no principled way to ensure the hypotheses cover the relevant uncertainty. In their didactic example this works because the answer space is small (tokens 0, 1, 2). For harder problems the model may generate highly correlated hypotheses that don't span the true uncertainty.

In CSGG this gap is critical: the hypothesis space is the secret $x^* \in [N]$, which is explicit and enumerable. We know exactly which hypotheses are consistent with the history --- $\mathcal{C}_t$. We don't need to sample hypotheses from the model; we can compute them exactly. This is a structural advantage of CSGG over the open-domain reasoning settings BARL targets.

\subsubsection{Gap 2 --- The Reward Consistency Term Is a Soft Approximation}

Their reward consistency term
\[
\prod_{t'}
\exp\left(
-\beta
\left|
r_{t'} - r_{\mathcal{M}_i}(s_{t'},a_{t'})
\right|
\right)
\]
is a soft likelihood that assigns non-zero weight to hypotheses inconsistent with observed rewards. This is necessary in their setting because rewards are noisy and sparse. But it introduces approximation error --- the posterior is not exact. With $\beta \to \infty$ the term becomes an indicator function (their didactic example uses this), but for finite $\beta$ there is bias.

In CSGG we can compute the exact posterior --- $\mathcal{C}_t$ is the exact set of consistent hypotheses after filtering by oracle hints. The ``reward consistency'' is exact filtering, not soft weighting. This means our Bayesian update is exact while BARL's is approximate.

\subsubsection{Gap 3 --- No Per-Turn Credit Assignment}

BARL's posterior-weighted Q-value is computed at each step $t$, but the Q-value
\[
Q^{\pi_\theta}_{\mathcal{M}_i}(h_t,a_t)
\]
is still an episode-level quantity --- it's the expected sum of future rewards under hypothesis $\mathcal{M}_i$. This means BARL still conflates per-step contributions into an episode-level signal, just weighted by a belief distribution. So, if BARL reduces to GRPO in our empirical results, then we clearly show that our cross-generation method is learning credit assignment better than this reduced BARL. So, in a way, since current research in GRPO shows that it is not suitable for multiturn setups with stochastic transition dynamics, we can show that we resolve this per-turn credit assignment problem within episodes using the Bayesian formulation. So, we can make a claim that Bayesian GRPO can help solve the multi-turn credit assignment problem, but for this we would need experiments across some other domains like math-reasoning etc. 

This is exactly the same problem as episode return GRPO --- the advantage signal at turn $t$ doesn't isolate what turn $t$'s action contributed relative to turn $t+1$'s. BARL improves over standard RL by tracking belief evolution, but it doesn't solve the within-episode credit assignment problem.

\subsubsection{Gap 4 --- The Method Section (5) and the Theory Section (4) Are Disconnected}

Their theoretical results (Theorems 4.1--4.3) show that Bayesian RL admits an uncertainty-adaptive optimal policy and that the gap between Markovian and adaptive policies can be arbitrarily large. But their method section doesn't prove that BARL actually converges to the uncertainty-adaptive optimal policy --- it just approximates the Bayesian RL gradient. There is no convergence guarantee connecting their algorithm to their theory.

This is the most important theoretical gap. The theory says Bayesian RL is better than standard RL. The method approximates Bayesian RL. But there is no result saying the approximation is tight or that BARL converges to the same fixed point as exact Bayesian RL.

\subsubsection{Gap 5 --- No Treatment of Information Gain Per Step}

BARL's objective maximizes expected return under the posterior --- but it doesn't explicitly reward actions that reduce posterior uncertainty. The reward consistency term passively updates beliefs based on observed rewards, but there is no active incentive to ask informative questions. In exploration settings (including CSGG) the most valuable actions are those that maximally differentiate between hypotheses --- i.e., those with high mutual information with the hypothesis. BARL doesn't explicitly encourage this.

\subsection{Their Proofs --- What's Useful for Us}

\subsubsection{Theorem 4.1 Proof (Markovian Optimal Policy in Standard RL)}

The proof is elegant --- assume the optimal policy uses different distributions at the same state depending on history, take the better distribution and apply it always, strictly improve. The key step is that future trajectories are stochastically identical under the two policies once they reach the same state --- so history doesn't matter for future value. This is the argument we need to run in reverse: show that in the IA-MDP, future trajectories are NOT stochastically identical at the same turn number with different histories, because the belief state is different.

\subsubsection{Theorem 4.2 Proof (Bayesian RL Admits Adaptive Optimal Policy)}

The proof uses the epistemic POMDP construction --- embed the MDP uncertainty into the hidden state by defining
\[
z_t = (s_t,\mathcal{M}).
\]
The optimal policy for this POMDP is a function of the belief
\[
\beta_t(z)
=
\mathbb{P}(z_t=z\mid h_t,D),
\]
which in their setting factors as
\[
\beta_t(s,\mathcal{M})
=
\mathbf{1}[s=s_t]\cdot p(\mathcal{M}\mid h_t,D).
\]
This is exactly the belief factorization we need for the IA-MDP --- in our case the hidden state is $(t,x^*)$ and the belief factors as
\[
\beta_t(t',x^*)
=
\mathbf{1}[t'=t]\cdot b_t(x^*),
\]
where
\[
b_t = \mathrm{Uniform}(\mathcal{C}_t).
\]

\subsubsection{Theorem 4.3 Proof (Gap Between Markovian and Adaptive)}

The proof constructs a binary decision tree where each leaf has reward 1 under exactly one hypothesis MDP. The Markovian policy must distribute probability mass across leaves without adapting to observed rewards --- its expected return is
\[
\frac{1}{2^{T^*}}.
\]
The adaptive policy uses the history to eliminate inconsistent hypotheses and converges to the rewarding leaf --- its expected return is 1. We can adapt this directly: construct a version of CSGG where the Markovian gap is explicit --- a set of secrets that can only be distinguished by adaptive questioning, where any fixed question sequence has low expected resolve rate.




\begin{remark}
The history $h_t$ grows monotonically and is never revisited — 
the dynamics form a directed acyclic graph (DAG) over histories. 
This tree-structured property is shared with the token-level LLM 
MDP \citep{rafailov2024r} and the AOH-structured Dec-POMDP of 
\citet{obel2021}, and is essential for the bijection result in 
Lemma~\ref{lem:bijection}.
\end{remark}

The arm partition $\{\mathcal{A}_k\}$ reflects distinct 
\emph{information sources}, each with its own expertise domain. 
We impose a structural condition on how beliefs evolve:

\begin{assumption}[Arm Independence]
\label{assump:arm-independence}
The belief state factorizes across arms:
$b_t = \prod_{k=1}^{K} b_t^k$,
where $b_t^k \in \Delta(\mathcal{X}_k)$ is the marginal belief 
relevant to arm $k$'s information domain, and 
$\mathcal{X} = \prod_{k=1}^K \mathcal{X}_k$.
An action $a_t \in \mathcal{A}_k$ updates only $b_t^k$ and 
leaves $b_t^{k'}$ unchanged for all $k' \neq k$.
\end{assumption}

Assumption~\ref{assump:arm-independence} holds whenever arms 
provide non-overlapping information about the latent variable — 
a natural condition in many sequential question-answering, 
diagnostic reasoning, and information retrieval settings.

%─────────────────────────────────────────────────────────────────
\subsection{Reward Structure and Belief Collapse}
%─────────────────────────────────────────────────────────────────

We consider rewards that measure \emph{fractional belief 
reduction} — how much uncertainty is eliminated per step:

\begin{definition}[Fractional Elimination Reward]
\label{def:reward}
A reward function $r$ is a \emph{fractional elimination reward} 
if there exists a measure $\mu$ on $\mathcal{X}$ such that:
$$r(h_t, a_t) = 1 - \frac{\mu(b_{t+1})}{\mu(b_t)}$$
where $\mu(b_t) = \int_\mathcal{X} b_t(x) d\mu(x)$ measures 
the spread of the current belief.
\end{definition}

Under a uniform prior and finite $\mathcal{X}$, this reduces to
$r(h_t, a_t) = 1 - |\mathcal{C}_{t+1}|/|\mathcal{C}_t|$,
where $\mathcal{C}_t = \mathrm{supp}(b_t)$ is the support of 
the current belief. This reward is dense, non-negative, and 
bounded: $r_t \in [0, 1]$ for all $t$, with $r_t = 1$ 
if and only if the latent variable is uniquely identified.

\begin{remark}
Maximizing cumulative fractional elimination reward is equivalent 
to minimizing the entropy of the belief state trajectory 
$\{b_t\}_{t=0}^T$. This connects the RL objective directly to 
information-theoretic optimal sequential hypothesis testing 
\citep{wald1947sequential}.
\end{remark}

%─────────────────────────────────────────────────────────────────
\subsection{Bijection Between Rewards and Q-Functions}
%─────────────────────────────────────────────────────────────────

We first establish that the IA-MDP has a bijection between 
reward functions and optimal Q-functions, following the approach 
of \citet{rafailov2024r} extended to the history-based setting.

\begin{lemma}[Reward-Q Bijection]
\label{lem:bijection}
In any IA-MDP with tree-structured history dynamics, there is 
a bijection between reward functions $r: \mathcal{H} \times 
\mathcal{A} \rightarrow \mathbb{R}$ and optimal Q-functions 
$Q^*: \mathcal{H} \times \mathcal{A} \rightarrow \mathbb{R}$ 
under the KL-regularized objective:
$$\max_\pi \mathbb{E}_\pi\left[\sum_{t=0}^{T-1} \gamma^t r_t\right] 
- \beta \cdot \mathrm{KL}(\pi \| \pi_\mathrm{ref})$$
\end{lemma}

\begin{proof}[Proof Sketch]
The tree-structured history dynamics ensure that $h_t$ cannot 
be revisited after any sequence of actions. This allows a 
backward induction argument: starting from terminal histories, 
the optimal Q-function $Q^*(h_T, \cdot) = 0$ is uniquely 
determined; proceeding backwards, $Q^*(h_t, a_t)$ is uniquely 
determined by $r(h_t, a_t)$ and $Q^*$ at successor histories, 
which are unique by induction. The surjection follows by 
inverting the Bellman equation $r(h_t, a_t) = Q^*(h_t, a_t) - 
\beta \log \pi_\mathrm{ref}(a_t|h_t) - V^*(h_{t+1})$.
Full proof in Appendix~\ref{app:proof-bijection}.
\end{proof}

\begin{corollary}[Policy as Q-Function]
\label{cor:policy-as-q}
Any policy $\pi$ is the optimal soft Q-function for some reward 
function $r$ in the IA-MDP. Specifically, the fixed-point 
relationship holds:
$$\pi^*(a_t|h_t) \propto \pi_\mathrm{ref}(a_t|h_t) \cdot 
\exp\!\left(\frac{Q^*(h_t, a_t)}{\beta}\right)$$
and equivalently:
$$\beta \log \frac{\pi^*(a_t|h_t)}{\pi_\mathrm{ref}(a_t|h_t)} 
= A^*(h_t, a_t)$$
where $A^*(h_t, a_t) = Q^*(h_t, a_t) - V^*(h_t)$ is the 
optimal advantage function.
\end{corollary}

Corollary~\ref{cor:policy-as-q} establishes that the log ratio 
of the learned policy to the reference policy is exactly the 
optimal advantage at convergence — the central identity 
connecting online RL to the Gittins structure we develop next.

%─────────────────────────────────────────────────────────────────
\subsection{Gittins Index Decomposition of the Optimal Q-Function}
%─────────────────────────────────────────────────────────────────

We now show that under Assumption~\ref{assump:arm-independence}, 
the optimal Q-function decomposes into per-arm Gittins indices 
— establishing that reward-maximizing RL in the IA-MDP is 
equivalent to Gittins index learning.

Recall the Gittins index for arm $k$ with belief state $b^k$:

\begin{definition}[Gittins Index \citep{gittins1979bandit}]
\label{def:gittins}
The \emph{Gittins index} of arm $k$ at belief state $b^k$ is:
$$G(k, b^k) = \sup_{\tau > 0} 
\frac{\mathbb{E}\!\left[\sum_{s=0}^{\tau-1} \gamma^s 
r_s^k \,\Big|\, b^k\right]}
{\mathbb{E}\!\left[\sum_{s=0}^{\tau-1} \gamma^s 
\,\Big|\, b^k\right]}$$
where the supremum is over all stopping times $\tau$, and 
$r_s^k$ is the reward obtained from arm $k$ at step $s$.
\end{definition}

The Gittins index is the maximum reward rate achievable before 
retiring arm $k$, given its current belief state. Under 
independence (Assumption~\ref{assump:arm-independence}), the 
optimal policy is to always select the highest-index arm 
\citep{gittins1979bandit, weber1992gittins}.

\begin{proposition}[Gittins Decomposition of $Q^*$]
\label{prop:gittins-decomp}
Under Assumption~\ref{assump:arm-independence} and fractional 
elimination rewards (Definition~\ref{def:reward}), the optimal 
Q-function decomposes as:
$$Q^*(h_t, a_t \in \mathcal{A}_k) = 
G(k, b_t^k) \cdot \mathbb{E}\!\left[\sum_{s=t}^{\tau_k^*-1} 
\gamma^{s-t}\right] + V^*(h_t^{-k})$$
where $\tau_k^*$ is the Gittins-optimal stopping time for arm 
$k$, and $V^*(h_t^{-k})$ is the value contribution from all 
arms $k' \neq k$.
\end{proposition}

\begin{proof}[Proof Sketch]
By arm independence, the optimal value function decomposes as 
$V^*(h_t) = \sum_{k=1}^K G(k, b_t^k) \cdot w_t^k$, where 
$w_t^k$ is the expected discounted time allocated to arm $k$ 
under the optimal policy \citep{weber1992gittins}. The 
Q-function for selecting arm $k$ at turn $t$ is the marginal 
contribution of arm $k$ to $V^*(h_t)$, which by the Gittins 
decomposition equals $G(k, b_t^k)$ times the expected remaining 
time budget, plus the independent contribution from other arms.
Full proof in Appendix~\ref{app:proof-gittins-decomp}.
\end{proof}

\begin{theorem}[Optimal Advantage as Gittins Index]
\label{thm:advantage-gittins}
The optimal advantage function in the IA-MDP satisfies:
$$A^*(h_t, a_t \in \mathcal{A}_k) = 
G(k, b_t^k) \cdot c_t - \lambda_t$$
where $c_t = \mathbb{E}[\sum_{s=t}^{\tau_k^*-1} \gamma^{s-t}]$ 
is the expected discounted time budget for arm $k$, and 
$\lambda_t = V^*(h_t) - V^*(h_t^{-k})$ is the \emph{opportunity 
cost} of allocating a turn to arm $k$ rather than the best 
alternative.
\end{theorem}

\begin{proof}
Follows directly from Proposition~\ref{prop:gittins-decomp} 
and $A^*(h_t, a_t) = Q^*(h_t, a_t) - V^*(h_t)$:
\begin{align*}
A^*(h_t, a_t \in \mathcal{A}_k) 
&= Q^*(h_t, a_t) - V^*(h_t) \\
&= G(k, b_t^k) \cdot c_t + V^*(h_t^{-k}) - V^*(h_t) \\
&= G(k, b_t^k) \cdot c_t - \lambda_t. \qquad \square
\end{align*}
\end{proof}

Theorem~\ref{thm:advantage-gittins} is the central theoretical 
result: \textbf{the optimal advantage function in any IA-MDP 
is proportional to the Gittins index of the selected arm, 
minus a turn-dependent baseline}. Combined with 
Corollary~\ref{cor:policy-as-q}, this means the log ratio of 
the optimal policy to the reference policy encodes the Gittins 
index at each step. An RL algorithm that accurately estimates 
$A^*(h_t, a_t)$ per step implicitly approximates the Gittins 
index from experience — without requiring an explicit belief 
state tracker or index computation.

%─────────────────────────────────────────────────────────────────
\subsection{Policy Gradient Convergence and Suboptimality Bounds}
%─────────────────────────────────────────────────────────────────

We now analyze the online policy gradient in the IA-MDP and 
bound the suboptimality gap. The online KL-regularized objective 
is:
$$J(\pi) = \mathbb{E}_\pi\!\left[\sum_{t=0}^{T-1} \gamma^t r_t\right] 
- \beta \cdot \mathbb{E}_{h_t \sim d^\pi}\!\left[
\mathrm{KL}(\pi(\cdot|h_t) \,\|\, \pi_\mathrm{ref}(\cdot|h_t))
\right]$$
where $d^\pi \in \Delta(\mathcal{H} \times \mathcal{A})$ is the 
state-action occupancy measure of $\pi$.

\begin{lemma}[Performance Difference in IA-MDP]
\label{lem:pdl}
For any two policies $\pi, \pi'$ in the IA-MDP:
$$J(\pi^*) - J(\pi) = 
\frac{1}{1-\gamma} \mathbb{E}_{h_t \sim d^\pi}
\!\left[A^{\pi^*}(h_t, \pi(h_t))\right]$$
where $A^{\pi^*}$ is the advantage function under $\pi^*$ and 
the expectation is over the state visitation distribution 
$d^\pi$ induced by the current policy $\pi$.
\end{lemma}

\begin{proof}
The Performance Difference Lemma \citep{kakade2002approximately} 
applies to any MDP with a well-defined value function and 
occupancy measure. The IA-MDP satisfies these conditions: 
$V^\pi(h)$ is well-defined for all $h \in \mathcal{H}$, the 
occupancy measure $d^\pi$ is well-defined over the DAG of 
histories, and the Bellman equations hold by 
Lemma~\ref{lem:bijection}. The standard PDL derivation applies 
directly.
Full proof in Appendix~\ref{app:proof-pdl}.
\end{proof}

Substituting Theorem~\ref{thm:advantage-gittins} into 
Lemma~\ref{lem:pdl}:

\begin{corollary}[Suboptimality Gap via Gittins]
\label{cor:subopt-gittins}
The suboptimality of any policy $\pi$ relative to the 
Gittins-optimal policy $\pi^*_G$ satisfies:
$$J(\pi^*_G) - J(\pi) = 
\frac{1}{1-\gamma} \mathbb{E}_{h_t \sim d^\pi}
\!\left[G(k^*(h_t), b_t^{k^*}) \cdot c_t - \lambda_t - 
A^\pi(h_t, a_t)\right]$$
where $k^*(h_t) = \arg\max_k G(k, b_t^k)$ is the 
Gittins-optimal arm at history $h_t$.
\end{corollary}

This corollary formalizes the intuition that suboptimality 
arises from failing to select the highest-index arm — the gap 
is zero if and only if $\pi$ always selects $k^*$ and produces 
correct advantage estimates.

\begin{theorem}[Suboptimality Bound via Advantage Estimation Error]
\label{thm:subopt-bound}
Let $\hat{A}_t$ be any advantage estimator used in policy 
gradient updates, and let 
$\epsilon_A = \mathbb{E}_{h_t \sim d^\pi}\!\left[
(\hat{A}_t(h_t, a_t) - A^*(h_t, a_t))^2\right]^{1/2}$ 
be its root-mean-squared error. Then after $N$ gradient steps:
$$J(\pi^*_G) - J(\pi_N) \leq 
\frac{C_{\nu,\pi}}{1-\gamma} \cdot \epsilon_A$$
where $C_{\nu,\pi} = \max_{h,a} d^\pi(h,a)/\nu(h,a)$ is the 
density ratio coefficient measuring distribution shift between 
the current policy and the data distribution $\nu$.
\end{theorem}

\begin{proof}[Proof Sketch]
From Lemma~\ref{lem:pdl}, $J(\pi^*_G) - J(\pi) \leq 
\frac{1}{1-\gamma} \mathbb{E}_{d^\pi}[|A^*(h_t,a_t) - 
\hat{A}_t(h_t,a_t)|]$. By Cauchy-Schwarz and 
Assumption~\ref{assump:arm-independence}, this is bounded by 
$\frac{1}{1-\gamma} \sqrt{\mathbb{E}_{d^\pi}[(\hat{A}_t - 
A^*)^2]}$. The density ratio $C_{\nu,\pi}$ converts the 
expectation from $d^\pi$ to $\nu$ following 
\citet{zhan2022offline}.
Full proof in Appendix~\ref{app:proof-subopt}.
\end{proof}

Theorem~\ref{thm:subopt-bound} has a direct consequence: 
\textbf{the quality of convergence to the Gittins-optimal policy 
is entirely determined by the quality of the per-step advantage 
estimator}. An estimator with $\epsilon_A \to 0$ as $N \to 
\infty$ converges to $\pi^*_G$; an estimator with structural 
bias that does not vanish with $N$ is bounded away from 
$\pi^*_G$ regardless of training duration.

%─────────────────────────────────────────────────────────────────
\subsection{Structural Bias of Conflated Advantage Estimators}
%─────────────────────────────────────────────────────────────────

We now formalize the structural requirement on advantage 
estimators and show that a broad class of estimators — those 
that conflate per-step advantages into episode-level quantities 
— have irreducible bias in the IA-MDP.

\begin{definition}[Turn-Level Decomposability]
\label{def:decomposable}
An advantage estimator $\hat{A}_t$ is \emph{turn-level 
decomposable} if it produces statistically independent 
estimates at each step $t$:
$$\hat{A}_t(h_t, a_t) \perp\!\!\!\perp \hat{A}_{t'}(h_{t'}, a_{t'}) 
\quad \text{for } t \neq t', \text{ given } h_t$$
and satisfies $\mathbb{E}[\hat{A}_t | h_t, a_t] = 
A^*(h_t, a_t)$ (unbiasedness at each step).
\end{definition}

\begin{definition}[Episode-Conflated Estimator]
\label{def:conflated}
An advantage estimator $\hat{A}_t$ is \emph{episode-conflated} 
if it assigns the same value to all steps within an episode:
$$\hat{A}_t(h_t, a_t) = \phi\!\left(\sum_{t'=0}^{T-1} 
r_{t'}\right) \quad \forall t \in [T]$$
for some function $\phi: \mathbb{R} \rightarrow \mathbb{R}$.
\end{definition}

\begin{theorem}[Irreducible Bias of Episode-Conflated Estimators]
\label{thm:bias}
Any episode-conflated estimator has irreducible bias in the 
IA-MDP: there exists a constant $\delta > 0$ such that for 
all $N \geq 1$:
$$\mathbb{E}\!\left[\left(\hat{A}_t^{\mathrm{conflated}}(h_t, a_t) 
- A^*(h_t, a_t)\right)^2\right]^{1/2} \geq \delta$$
where $\delta$ depends on the variance of per-step Gittins 
indices $\{G(k, b_t^k)\}_{t=0}^{T-1}$ across turns.
\end{theorem}

\begin{proof}[Proof Sketch]
By Definition~\ref{def:conflated}, the conflated estimator 
assigns $\hat{A}_t = \phi(\sum_{t'} r_{t'})$ to all turns. 
The true advantage at turn $t$ is $A^*(h_t, a_t) = G(k, b_t^k) 
\cdot c_t - \lambda_t$ (Theorem~\ref{thm:advantage-gittins}), 
which varies across turns because the belief state $b_t^k$ 
evolves. The conflated estimator cannot recover this variation: 
$\mathbb{E}[\hat{A}_t | \sum_{t'} r_{t'}] = 
\mathbb{E}[A^*(h_t, a_t) | \sum_{t'} r_{t'}]$, which equals 
$A^*(h_t, a_t)$ only if $A^*(h_t, a_t)$ is constant across 
turns — a condition violated whenever Gittins indices vary, 
which holds generically in any non-trivial IA-MDP.
The lower bound $\delta$ is the standard deviation of 
$\{A^*(h_t, a_t)\}_{t=0}^{T-1}$, which is strictly positive.
Full proof in Appendix~\ref{app:proof-bias}.
\end{proof}

Combining Theorems~\ref{thm:subopt-bound} 
and~\ref{thm:bias}:

\begin{corollary}[Performance Gap Lower Bound]
\label{cor:gap-lower}
Any policy optimized with an episode-conflated advantage 
estimator satisfies:
$$J(\pi^*_G) - J(\pi_N) \geq 
\frac{C_{\nu,\pi}}{1-\gamma} \cdot \delta > 0$$
for all $N \geq 1$, where $\delta > 0$ is the 
irreducible bias from Theorem~\ref{thm:bias}.
\end{corollary}

This is the strongest theoretical result: \textbf{episode-level 
advantage estimation produces a policy that is bounded away 
from the Gittins-optimal policy by a positive constant, 
regardless of the amount of training data or gradient steps}. 
This is not a capacity or sample complexity problem — it is a 
structural incompatibility between the estimator and the 
per-step Gittins structure of the IA-MDP.

%─────────────────────────────────────────────────────────────────
\subsection{Policy Improvement as Gittins Index Approximation}
%─────────────────────────────────────────────────────────────────

We conclude with a result connecting iterative policy 
improvement to progressive refinement of Gittins index 
estimates, analogous to the OBL policy improvement operator 
of \citet{obel2021}.

\begin{theorem}[Online Policy Gradient as Gittins Approximation]
\label{thm:pg-gittins}
Let $\pi_0 = \pi_\mathrm{ref}$ be any reference policy, and 
let $\pi_{n+1}$ be the policy after one gradient step from 
$\pi_n$ using a turn-level decomposable advantage estimator 
$\hat{A}_t$ with estimation error $\epsilon_A$. Then:
\begin{enumerate}
    \item \textbf{Policy improvement:} 
          $J(\pi_{n+1}) \geq J(\pi_n) - \epsilon_A / e$
    \item \textbf{Convergence:} 
          $\lim_{n \to \infty} J(\pi_n) = J(\pi^*_G) - 
          O(\epsilon_A)$
    \item \textbf{Log-ratio recovery:} At convergence, 
          $\beta \log \pi_n(a|h_t) / \pi_\mathrm{ref}(a|h_t) 
          \approx G(k, b_t^k) \cdot c_t - \lambda_t$
\end{enumerate}
\end{theorem}

\begin{proof}[Proof Sketch]
Part (1) follows from the OBL policy improvement argument 
\citep{obel2021} adapted to the single-agent IA-MDP: the 
softmax policy improvement step guarantees 
$J(\pi_{n+1}) \geq \max_a Q^{\pi_n}(h_t, a) - \epsilon_A/e$, 
where the $\epsilon_A/e$ term follows from Lemma 2 of 
\citet{obel2021} applied to the advantage estimation error.
Part (2) follows from the contraction of the softmax policy 
improvement operator and Theorem~\ref{thm:subopt-bound}.
Part (3) follows from Corollary~\ref{cor:policy-as-q} and 
Theorem~\ref{thm:advantage-gittins} at convergence.
Full proof in Appendix~\ref{app:proof-pg-gittins}.
\end{proof}

Part (3) of Theorem~\ref{thm:pg-gittins} is the central 
empirical prediction of our theory: \textbf{the log ratio of 
the trained policy to the reference policy encodes the Gittins 
indices of the arms at each step}. An RL-trained language model 
that has implicitly learned to optimize fractional elimination 
rewards should exhibit, in its log-probabilities, a signal 
proportional to the Gittins index — the value of continuing 
to query each arm given the current belief state. We provide 
empirical evidence for this prediction in 
Section~\ref{sec:experiments}.

%─────────────────────────────────────────────────────────────────
\subsection{Instantiation: The Clue Selector Game}
%─────────────────────────────────────────────────────────────────

The Clue Selector Game (Section~\ref{sec:csgg}) is an instance 
of the IA-MDP with $\mathcal{X} = [N]$, $K = 5$ arms with 
disjoint question domains, and fractional elimination rewards 
over a uniform prior. Assumption~\ref{assump:arm-independence} 
holds by the disjoint arm property construction 
(Appendix~\ref{app:arm-independence}).

The theorems above therefore predict:
\begin{enumerate}
    \item A policy gradient algorithm with turn-level advantage 
          estimation converges to near-Gittins-optimal behavior 
          — characterized by progressive belief collapse toward 
          the theoretical optimum of halving candidates per turn.
    \item An episode-level advantage estimator produces a policy 
          that is bounded away from Gittins-optimal by 
          $\delta > 0$ regardless of training duration.
    \item At convergence, the trained model's log-probabilities 
          encode arm-level Gittins indices — manifesting 
          empirically as concentrated queries toward arms with 
          high current discrimination power.
\end{enumerate}
These predictions are verified empirically in 
Section~\ref{sec:experiments}, where we observe near-perfect 
binary search emergence, bounded suboptimality for episode-level 
baselines, and arm concentration patterns consistent with 
implicit Gittins index learning.



\section{Notes: Gittins-Structured Credit Assignment for Language Agents}

\subsection{Motivation}

\an{Standard policy gradient methods assign credit to actions or to an LLM's response through 
advantage estimates derived from learned value functions or group-normalized 
returns~\citep{sutton2018reinforcement, schulman2017proximalpolicyoptimizationalgorithms, schulman2015high, shao2024deepseekmath}. In interactive tasks where an LLM needs to communicate with a simulator\footnote{This can be a group of agents \textit{roleplayed}~\citep{ivison2023camels, nath2026collaboratedeliberateevaluatellm} by a central coordinator agent or simply a “user” simulator initialized with a strong LLM.} before acting in the environment or in multi-agent information-seeking tasks, this is theoretically 
suboptimal: the value of allocating scarce “attention” or turn-budget—proxies for the cost of communication in interactive settings—is similar to the classic multi-armed bandit (MAB) setup of querying an arm depends on its posterior belief 
state. While most current research in designing LLM agents for such tasks leverage standard temporal difference (TD)-based critics~\citep{schulman2015high}, it is a challenge for such RL-trained agents to recover the optimal structure of such resource allocation—without relying on more principled theories of arm selection from classic problem in optimal stopping theory. Notably, the Gittins index~\citep{weber1992gittins, whittle1980multi} provides an analytical solution to this 
allocation problem in theory where the distribution of the “cost” (e.g., the reward distribution) of communicating is known prior—an assumption that is rarely satisfied in complex real-world environments where agents operate under \textit{partial information} but are \textit{expected to act optimally under uncertainty} and where these agents are increasingly being deployed today. Importantly, the idea of using such indices  has never been applied in settings where the 
belief state is encoded in natural language rather than tracked numerically—more so in the context of “learning” in the environment as in standard Reinforcement learning.}

\textbf{We establish two theoretically principled methods for incorporating Gittins 
indices into policy gradient training, both of which preserve the optimal 
policy under standard conditions.}

\subsection{Gittins-Structured Reward Shaping (GSRS)}

\paragraph{Setup.} Let $s_t$ denote the agent's state at turn $t$, 
comprising the current candidate set $\mathcal{C}_t$ and the arm histories 
$\mathcal{H}_t = \{\mathcal{H}_t^k\}_{k=1}^K$. The agent selects arm $k_t$ 
and asks question $q_t$, receiving reward $r_t = |\mathcal{C}_t| - 
|\mathcal{C}_{t+1}| / |\mathcal{C}_t|$. Let $\lambda^*(k, \mathcal{H}_t^k)$ 
denote the Gittins index of arm $k$ given its history, approximated by the 
regression model $f_\theta$.

\paragraph{Shaped reward.} We define the Gittins-shaped reward as:
\begin{equation}
    \tilde{r}_t = r_t + \gamma \Phi(s_{t+1}) - \Phi(s_t)
\end{equation}
where the potential function is:
\begin{equation}
    \Phi(s_t) = \frac{1}{K} \sum_{k=1}^K f_\theta(\mathcal{H}_t^k)
    = \frac{1}{K} \sum_{k=1}^K \lambda^*(k, \mathcal{H}_t^k)
\end{equation}
That is, the potential at state $s_t$ is the mean Gittins index across all 
arms given their current histories.

\paragraph{Policy invariance.} By the potential-based reward shaping 
theorem of \citet{ng1999policy}, any reward function of the form 
$\tilde{r}(s, a, s') = r(s, a, s') + \gamma\Phi(s') - \Phi(s)$ yields 
the same optimal policy as the original reward $r$, for any bounded 
potential function $\Phi$. Since the Gittins index is bounded in $[0, 1]$ 
for Beta-Bernoulli arms, $\Phi(s_t) \in [0, 1]$ and the condition is 
satisfied. This guarantee is analogous to the PBRS framework used in 
\citet{cao2025scar, nath2026owen} for Shapley (or Owen value)-based credit assignment in RL training for language-based tasks. 

\paragraph{Interpretation.} The potential $\Phi(s_t)$ increases when 
arms have high remaining information value and decreases as arms become 
exhausted. \an{The shaping term $\gamma\Phi(s_{t+1}) - \Phi(s_t)$ therefore 
provides a dense signal that rewards transitions which preserve high-index 
arms for future use—a form of information conservation credit that 
standard reward functions cannot express.}

\subsection{Gittins Index Advantage Weighting (GIAW)}

\paragraph{Setup.} In GRPO, for each prompt $x$, $G$ completions 
$\{o_i\}_{i=1}^G$ are sampled and advantages are computed by 
group-normalizing the rewards:
\begin{equation}
    A_i^{\text{GRPO}} = \frac{r_i - \mu_G}{\sigma_G + \epsilon}
\end{equation}
where $\mu_G$ and $\sigma_G$ are the mean and standard deviation of rewards 
within the group.

\paragraph{Gittins-weighted advantage.} Let $k_i$ denote the arm selected 
in completion $o_i$ and let $\rho_i = f_\theta(\mathcal{H}^{k_i})$ be its 
Gittins index. We define the arm's rank-normalized index score:
\begin{equation}
    \hat{\rho}_i = \frac{\text{rank}(k_i \mid \{f_\theta(\mathcal{H}^k)\}_{k=1}^K)}{K}
    \in [0, 1]
\end{equation}
The Gittins-weighted advantage is then:
\begin{equation}
    \tilde{A}_i = A_i^{\text{GRPO}} \cdot (1 + \alpha \cdot \hat{\rho}_i)
\end{equation}
for a hyperparameter $\alpha \geq 0$ controlling the strength of the 
index weighting.

\paragraph{Relationship to OSPO.} This formulation is analogous to the 
Owen-Shapley Policy Optimization framework \citep{nath2026owen}, where 
sequence-level advantages are redistributed to tokens proportionally to 
their Owen-Shapley values. Here, the redistribution operates at the 
\emph{arm selection} level rather than the token level: completions that 
selected high-index arms receive amplified advantages, biasing the policy 
gradient toward Gittins-optimal arm allocation. Unlike token-level 
redistribution, GIAW preserves sequence-level advantages and applies a 
multiplicative scaling, making it compatible with standard GRPO 
implementations without modifying the loss function.

\paragraph{Practical note.} Rank normalization is used rather than raw 
index values to avoid scale sensitivity — the Gittins index magnitude 
varies with $(\alpha, \beta)$ in ways that may not align with the scale 
of GRPO advantages. Rank normalization maps all arms to $[0,1]$ 
regardless of their absolute index values.

\subsection{Summary of Theoretical Properties}

\begin{enumerate}

\item \textbf{GSRS preserves the optimal policy by construction.} 
The potential-based shaping guarantee holds for any bounded $\Phi$, 
and the Gittins index satisfies this bound. No additional assumptions 
about the task structure are required beyond those already needed for 
standard policy gradient convergence.

\item \textbf{GIAW introduces a Gittins-aware inductive bias without 
modifying the reward.} By weighting advantages rather than rewards, 
GIAW does not alter the target policy and does not require verifying 
PBRS conditions. The weighting amplifies the learning signal for 
Gittins-optimal actions while leaving suboptimal actions with reduced 
but non-zero gradient contribution.

\item \textbf{Both methods degrade gracefully to their baselines.} 
Setting $\Phi \equiv 0$ recovers standard PPO from GSRS. Setting 
$\alpha = 0$ recovers standard GRPO from GIAW. This makes both 
methods directly comparable to their baselines in ablation studies.

\item \textbf{The Gittins structure is a sufficient condition, not a 
necessary one.} Neither method requires the task to exactly satisfy 
the Beta-Bernoulli bandit assumptions. As long as the regression 
model $f_\theta$ produces monotonically meaningful index estimates — 
which the empirical results confirm at $r = 0.936$ Spearman correlation 
— the shaped rewards and weighted advantages will preferentially 
reinforce informative arm selections over exhausted ones.

\item \textbf{The connection to natural language belief states is the 
key generalization.} Classical Gittins index theory assumes numerically 
tracked posteriors. The contribution here is showing that the index 
can be approximated from implicit belief states encoded in conversational 
histories, extending the theory to language agent settings where 
explicit posterior tracking is infeasible.

\end{enumerate}

\section{Notes: Gittins Index Regression — Design Choices and Insights}

\subsection{Goal}

The central question motivating this section is whether the Gittins index—a theoretically optimal solution to the multi-armed bandit problem under 
uncertainty—can be approximated from natural language descriptions of arm 
states. If so, it becomes possible to deploy index-based arm selection in 
settings where the belief state is not numerically tracked but is instead 
encoded implicitly in a conversational history. \an{We train a regression model 
that maps a sentence-transformer embedding of an arm's interaction history 
to its Gittins index value, and evaluate five distinct label construction 
strategies that differ in what ground truth they approximate and what 
information they require. \textbf{A core hypothesis is that LLM's under RL-training can implicitly learn to encode Gittin index properties (e.g., of knowing when to stop querying an information reservoir or “user” simulator or who to seek more information from) in its CoT reasoning traces—purely from empirical estimate of such values using a strong function approximator like a neural network. During our initial experiments, we found strong reasoning agents like Claude Sonnet 4.6 to show emergence of such Gittins-encoded optimal reasoning in the clue selection game; which is a game where an optimal strategy of gameplay naturally requires one to make such index-optimal choices to solve the task most efficiently.
}}

\subsection{Label Construction Strategies}

Each strategy defines a different notion of ``arm value'' that serves as 
the regression target. All strategies share the same input representation: 
the arm state is converted to a natural language string encoding the arm 
index, number of prior queries, the full question-answer history, and the 
number of remaining candidates. This string is embedded using a 
sentence transformer (\texttt{all-MiniLM-L6-v2}), producing a 384-dimensional 
vector augmented with two scalar features (normalized query count and 
normalized candidate count), yielding a 386-dimensional input.

\paragraph{Discounted Future Return (DFR).}
The label for an arm at turn $t$ is the discounted sum of future rewards 
received from that arm across the remainder of the episode:
\begin{equation}
    y_{\text{DFR}} = \frac{\sum_{j \geq t,\, a_j = k} \gamma^{s} \cdot r_j}
    {\sum_{s=0}^{n_k - 1} \gamma^s}
\end{equation}
where $s$ indexes pulls of arm $k$ after turn $t$, $r_j$ is the normalized 
elimination reward, $\gamma = 0.9$, and $n_k$ is the number of subsequent 
pulls. This label is derived entirely from offline rollout data and requires 
no external computation. However, it is confounded by the quality of the 
policy that generated the data: if a suboptimal policy asked poor questions 
to arm $k$, the label underestimates the true value of that arm.

\paragraph{Immediate Next Reward (INR).}
The label is simply the reward received the next time arm $k$ is pulled 
after turn $t$:
\begin{equation}
    y_{\text{INR}} = r_{j^*}, \quad j^* = \min\{j > t : a_j = k\}
\end{equation}
This is the most myopic label — it captures only the immediate value of 
the next pull without accounting for the arm's remaining information budget. 
It is fast to compute but discards all long-horizon structure.

\paragraph{Oracle Discrimination Budget (ODB).}
For each arm, the label is the total discrimination power of its remaining 
unasked questions, computed against the actual secret:
\begin{equation}
    y_{\text{ODB}} = \frac{1}{|\mathcal{C}|} \sum_{p \in \mathcal{P}_k \setminus \mathcal{A}_k} 
    \left| \mathcal{C} \right| - \left| \text{filter}(\mathcal{C}, \text{hint}(q_p, s)) \right|
\end{equation}
where $\mathcal{P}_k$ is the set of properties known to arm $k$, 
$\mathcal{A}_k$ is the set of already-asked questions, $\mathcal{C}$ is 
the current candidate set, and $s$ is the secret. This label requires 
access to the secret during training and thus cannot be used at inference 
time without oracle information. It measures how much information the arm 
could still provide if queried optimally.

\paragraph{Beta-Bernoulli Gittins Index (BBGI).}
\an{The label is the exact Gittins index for a Beta-Bernoulli arm whose 
posterior belief is tracked from the arm's interaction history.} Starting 
from a uniform prior $\text{Beta}(1, 1)$, each informative response 
(one that is neither redundant nor uninformative) increments $\alpha$ by 
one, and each uninformative response increments $\beta$ by one:
\begin{align}
    \alpha &= 1 + |\{(q, h) \in \mathcal{H}_k : h \text{ is informative}\}| \\
    \beta  &= 1 + |\{(q, h) \in \mathcal{H}_k : h \text{ is uninformative}\}|
\end{align}
The Gittins index $\lambda^*(\alpha, \beta)$ is then looked up from a 
precomputed dynamic programming table — the largest $\lambda$ for which 
a Beta$(\alpha, \beta)$ arm is preferred over a safe arm paying $\lambda$ 
per step:
\begin{equation}
    \lambda^*(\alpha, \beta) = \sup \left\{ \lambda : 
    V_{\text{risky}}(\alpha, \beta, \lambda) > \frac{\lambda}{1 - \gamma} \right\}
\end{equation}
The table is precomputed over a grid of $(\alpha, \beta) \in [1, 60]^2$ 
using binary search and backward induction with discount factor $\gamma = 0.9$ 
and horizon $N = 20$. In practice, since each arm can be pulled at most 
$T_{\max} = 10$ times, the reachable belief states satisfy 
$\alpha + \beta \leq T_{\max} + 2 = 12$, meaning only a small corner of 
the table is ever accessed during inference. The larger table is precomputed 
for safety but a $15 \times 15$ grid would suffice without loss.

\paragraph{Binary Exhaustion Indicator (BEI).}
The label is a binary indicator of whether the arm has any remaining 
unasked properties:
\begin{equation}
    y_{\text{BEI}} = \mathbf{1}\left[\mathcal{P}_k \setminus \mathcal{A}_k \neq \emptyset\right]
\end{equation}
This label is the simplest to compute and requires no external resources. 
It captures only whether an arm is worth pulling at all, not how valuable 
it is relative to other arms.

\subsection{Experimental Setup}

All five models share the same architecture: a three-layer MLP with hidden 
dimensions 128 and 64, ReLU activations, dropout of 0.1, and a sigmoid 
output. Models are trained for 100 epochs with Adam ($\text{lr} = 10^{-3}$) 
and cosine annealing. The train/validation split is by secret value: 
secrets 1--70 for training, 71--100 for validation, testing generalization 
to unseen game instances. Episodes are drawn from all available offline 
rollouts spanning Qwen models (0.5B--14B), frontier models 
(GPT-4o, Claude Sonnet, GPT-4.1-mini), and rule-based baselines 
(greedy oracle, random, round-robin).


\fi
\subsection{ RL Results}


% \begin{table*}[t]
% \centering
% \small
% \setlength{\tabcolsep}{5pt}
% \begin{tabular}{lcccccccc}
% \toprule
% \textbf{Model} & \textbf{Resolve} & \textbf{Eff.} & \textbf{Reward} & \textbf{Elim.} & \textbf{Final} & \textbf{Arm} & \textbf{Red.} & \textbf{Turns} \\
% \midrule
% [T] cross\_generation & \textbf{45.3$\pm$1.26} & \textbf{0.515} & \textbf{4.303} & \textbf{14.064} & 1.698 & 19.0\% & 0.5\% & \textbf{8.880} \\
% [F] gpt-4.1-mini & 31.6$\pm$1.92 & 0.472 & 4.025 & 13.830 & \textbf{1.604} & 17.2\% & 0.5\% & 9.116 \\
% [F] gpt-4o       & 25.8$\pm$2.21 & 0.441 & 4.017 & 13.136 & 1.733 & 17.5\% & 0.5\% & 9.449 \\
% [F] gpt-4o-mini  & 22.7$\pm$1.26 & 0.411 & 3.759 & 12.910 & 2.742 & 16.6\% & \textbf{0.2\%} & 9.538 \\

% [B] Qwen-14B     & 22.7$\pm$1.66 & 0.410 & 3.622 & 13.328 & 2.329 & 11.9\% & 1.6\% & 9.378 \\
% [B] Qwen-7B      & 20.9$\pm$0.73 & 0.402 & 3.504 & 13.229 & 3.249 & \textbf{20.1\%} & 1.7\% & 9.369 \\
% \bottomrule
% \end{tabular}
% \caption{Headline comparison across representative frontier, open-weight, and trained models. The cross-generation MC method achieves the highest resolve rate and efficiency, outperforming both frontier and open-weight baselines.}
% \label{tab:headline_results}
% \end{table*}



% \begin{figure*}[t]
%     \centering
%     \includegraphics[width=\textwidth]{rough_plots/grpo_multiturn_cross_generation_3seed.png}
%     \caption{
%     Cross-generation GRPO training dynamics for Qwen2.5-1.5B-Instruct with LoRA rank $r=16$, averaged over three seeds. Metrics show steady gains in reward, episode resolve rate, turn resolve rate, efficiency bonus, and turns saved, while redundancy falls to near zero and diversity remains high. Entropy decreases gradually, indicating specialization without immediate collapse.
%     }
%     \label{fig:grpo_multiturn_cross_generation}
% \end{figure*}


% \paragraph{Training Dynamics and Mechanistic Analysis.}
% Fig.~\ref{fig:grpo_multiturn_cross_generation} shows training dynamics across outcome, mechanism, and optimization-health metrics to understand both performance and the underlying learning behavior  

% At the outcome level, \textit{episode resolve rate} provides the primary signal of task success, improving substantially over training (reaching $\sim$74.8\%). This increase reflects the emergence of effective querying strategies once the environment provides a consistent reward signal. Supporting this, the \textit{redundancy rate} drops sharply (from $\sim$14\% to $\sim$0.2\%), confirming that the filtering mechanism successfully removes spurious or non-informative actions, thereby enabling meaningful learning. In parallel, \textit{mean eliminated} rises to $\sim$13.5 candidates per turn, indicating that the \textbf{policy converges toward near-optimal information-gathering behavior, approximating a binary partitioning strategy over the candidate set.}

% At the mechanism level, the policy demonstrates strong within-step reasoning but incomplete higher-level strategy. While candidate sets consistently shrink across turns (e.g., candidates-before decreasing from $\sim$28.7 to $\sim$25.6), the \textit{arm match rate} remains relatively low ($\sim$25.9\%), only slightly above random selection. This suggests that the policy has learned to ask informative questions conditioned on a chosen arm, but has not yet fully learned how to select the most informative arm itself. Consequently, \textit{mean episode length} decreases gradually (to $\sim$7.4), and \textit{turns saved} increases ($\sim$0.39), reflecting improved efficiency on successful episodes but leaving room for further gains through better arm selection. A complementary pattern emerges in \textit{turn resolve rate} and \textit{efficiency bonus}, both of which increase steadily, indicating that improvements in local decision quality translate into incremental gains in overall efficiency.

% At the optimization level, training remains stable but exhibits signs of progressive specialization. \textit{Entropy} decreases steadily (to $\sim$0.166), indicating that the policy is concentrating probability mass on a smaller set of actions. At the same time, \textit{diversity rate} remains high ($\sim$99.6\%), suggesting that this concentration reflects improved decision confidence rather than premature mode collapse. \textit{Reward standard deviation} remains relatively high ($\sim$0.47), indicating continued exposure to heterogeneous trajectories and the absence of degenerate policy collapse. Finally, \textit{gradient norms} remain below the clipping threshold, and the near-zero clip ratio indicates that optimization proceeds within a stable trust region.


\iffalse
\section{Lens: The pipeline is the asset; the task framing on top of it is the variable.}
Taken together, these results reveal a clear decomposition of phase-wise learning: the policy first acquires strong within-step information-gathering behavior, followed by gradual improvements in efficiency, while higher-level strategic components such as arm selection remain underdeveloped yet showing tendency of learning as multi-turn RL progresses. This is expected since learning the “joint”-policy over both the arm selection and the question is a challenging problem for three reasons: first, most “priors” that base models inherit from finetuning or post-training are about the skill of generating the final answer conditioned on the initial part (e.g., the CoT) of the response itself. This means that when the learning (or rather the way the task is structured or the action space itself) necessitates first generating the arm itself (unlike a standard CoT trace) and the shaped-reward for the policy—to learn the optimal joint policy of connecting the arm domain to the question asked-is marginal (small part of the total terminal reward). In these cases, we can expect the task itself to be challenging, so the lower arm-matching rate is expected on the first glance. 

\textbf{However, on second glance it could be that the limitation here is “definitional” or in the state construction shortcut that standard language RL applies}; can we think of tasks where this packing of trajectories could be problematic or difficult to achieve? This could help motivate a new dimension of learning in multiturn RL; Under what conditions does applying multi-turn RL with standard state construction assumptions leads to suboptimal or unsafe situations? The usual excuse here is that multi-turn interactions with an external simulator is costly which is true but it is also true that many works identify the challenge of finding safer ways for this interaction. For example, if the interaction is costly in a different way—say the user simulator (or our clue model) is constrained on their end (not our end) to not reveal certain types of information or if they find the policy is interacting in a certain way, provide false information;[\textcolor{purple}{at this point, my thoughts are strongly hinting at defining policy gradient in "natural language itself"—the most expressive form of reward~\citep{su2026end} that humans know of but not in a way as in the policy gets feedback in its context window and learns from it; something more fundamental }] The slight hint here is that many folks are using RL with a GPT simulator, so if we can show that for certain domains or tasks this can be risky not just because the feedback itself is difficult to get, but how to \textit{validate} (e.g., , how does the model learn to check for good sources of feedback while learning in the real world interactions that the field is pushing for these days) the feedback (or which feedback should the model pay attention to) and which ones to ignore; if you cannot define this with a scalar reward, you can show that with a standard reward, it tends to disturb the user model more in a certain behavioral pattern; Still not sure how much we have covered here but you get the hint of where I am going. 


Additionally, it is also true that since this run with the same base model which is on-policy from the start (i.e., no other policy generated all the trajectories that it saw during training vs say, starting with offline trajectory logs of a stronger model like Claude where they already contain “better” states for the model to optimize for during RL training where we earlier saw much stronger arm-match rates towards later steps due to this advantage)—this explain how off-policy learning from a stored trajectory buffer can help learn to optimize this joint policy. At a theoretical level, the true transition function in this multi-turn RL environment is non-deterministic since the environment feedback comes directly from the clue giver model that is a GPT4o-mini instance that can behave non-deterministically~\citep{gpt4}. So, although language RL assumptions allow us to pack all parts of the trajectory (feedback and remaining candidates etc, except the reward) as one single concatenated block and call it the “state”, in reality this does not capture the true dynamics of this interaction in the stated definition of the transition function~\citep{samineni2026rlonlyanalyzingstructural, kazemnejad2025vinepporefiningcreditassignment}. 

This degeneracy in the RL MDP formulation itself for language tasks is the second level of challenge in this task. Finally, this also makes a direct application of the value function difficult since generating the arm itself should ideally not be given any hint of how valuable the question is going to be (since the arm is generated first from the left)—however, some of it still makes sense since given a context, the policy should still get a value signal in selecting a good arm since arm selection itself directly (e.g., autoregressively) affects the probability of it generating the relevant question; Looking at the third angle, it seems like employing a value model could boost performance—but the core question that this leads to is that if standard MDP for language is “degenerate” due to the nature of this multiturn interaction, what could be a better alternative to formulate this? 



This separation highlights a key avenue for future work—improving long-horizon credit assignment to enable coordinated decision-making across steps, rather than only within them.

\fi


\paragraph{CSG Reward Composition} The reward  $r_t$ is computed as follows. See \Cref{sec:experiments} for motivation on why this reward design is appropriate for our theoretical and experimental setup.  Specifically, if the policy has already asked the same question to the selected arm earlier in the episode, the turn returns immediately with the redundancy penalty $r_t = -0.1$ and the oracle is not consulted. Otherwise the oracle returns a hint, the candidate set is filtered, and $r_t$ is the sum of up to five components: the fractional information gain $r_{\text{info}} = (\text{candidates\_before} - \text{candidates\_after}) / \text{candidates\_before}$, which lies in $[0,1]$ and dominates the dense per-turn signal; an arm-match bonus of $+0.05$ when the question matches a property the selected arm covers (keyword match against \texttt{ARM\_PROPERTIES}); a diversity bonus of $+0.05$ when the question has not been asked to any arm earlier in the episode; a sparse resolution bonus of $+1.0$ when the post-oracle candidate set equals $\{\text{secret}\}$; and an efficiency bonus of $0.1 \cdot (T_{\max} - t)$ on resolution, with $T_{\max} = 10$ turns. 


\paragraph{Training Hyperparameters} The KL coefficient is set to $\beta_{\text{KL}} = 0$ to remove the reference-model anchor and permit policy drift during the continuation phase; gradients are clipped at $\|\nabla\| \le 0.5$; the PPO clip range uses the TRL default. The information-gain term dominates over most of training (range $[0,1]$ per turn vs.\ $\le 0.1$ for the shaping bonuses), making per-turn elimination the primary gradient signal. The resolution bonus is sparse but large at $+1.0$ once per resolved episode, supplying the long-horizon reward needed to escape the local optimum of asking individually informative questions that fail to fully resolve. The redundancy penalty acts as a hard early-return: in the pre-fix regime the policy frequently selected duplicate (arm, question) pairs, triggering the penalty on a large fraction of turns and short-circuiting all other reward components, so the gradient saw $-0.1$ rather than the information-gain landscape underneath. Setting $\beta_{\text{KL}} = 0$ removes the reference-model anchor, trading trajectory-drift risk for faster specialization during the continuation phase. \\



\begin{table*}[t]
\centering
\scriptsize
\setlength{\tabcolsep}{5pt}
\renewcommand{\arraystretch}{1.15}
\begin{tabularx}{\linewidth}{p{3.0cm}p{4.6cm}X}
\toprule
\textbf{Metric} & \textbf{Definition} & \textbf{Diagnostic Role in Clue Selector Training} \\
\midrule
\multicolumn{3}{l}{\textit{Outcome metrics --- did the policy solve the task?}} \\


\midrule
Episode Resolve Rate & Fraction of episodes (capped at 10 turns) where the policy identifies the secret number. & Primary outcome. The only metric that directly answers whether the policy solves the task; bounded above by 1, so the trajectory bends as performance approaches the ceiling. \\
Turn Resolve Rate & Fraction of individual turns that reduce the candidate set to a single element. & Per-turn analogue of episode resolve. Distinguishes ``the policy resolved this episode'' from ``the policy resolved on this specific turn,'' surfacing whether resolution comes from a single decisive action or from gradual narrowing. \\
Reward & Mean scalar reward across all turns in the batch; primary GRPO signal. & Coarse summary of all reward components. Useful as a sanity check but conflates information gain, resolution bonus, and efficiency bonus; should always be read alongside the component metrics. \\
Normalized Reward & Reward divided by the maximum achievable per episode. & Scale-free measure of proximity to optimal play. Lets cross-configuration comparisons abstract away changes in reward weighting. \\
\midrule
\multicolumn{3}{l}{\textit{Mechanism metrics --- how is the policy solving (or failing to solve) the task?}} \\
\midrule
Eliminated Mean & Average number of candidates removed per turn after applying the oracle's response. & Direct measure of information gain per question. With candidate sets near 25--30, an information-maximizing question eliminates roughly half; values approaching this threshold indicate the policy has discovered near-optimal binary partitioning at the question level. \\
Candidates Before / After & Average candidate set size at the start of a turn / after applying the oracle's response. & Pre- and post-action uncertainty. The gap is per-turn information gain; the ratio is the fractional reduction. A falling \emph{Candidates Before} across the episode indicates earlier turns are doing useful work that compresses later turns' starting state. \\
Arm Match Rate & Fraction of turns where the selected arm matches the oracle's preferred arm. & Strategic alignment metric, distinct from question quality. With $K$ arms, the random baseline is $1/K$. Values near baseline while \emph{Episode Resolve Rate} climbs indicate the policy is winning through within-arm question quality rather than informed arm selection --- a decomposition that flags untapped headroom in higher-order strategy. \\
Mean Turn & Average turn index at which resolution occurs, conditioned on successful episodes. & Efficiency on solved episodes. Low values indicate the policy resolves early in the budget; the gap between this and \emph{Mean Episode Length} reflects the cost of unresolved episodes running to the cap. \\
Mean Episode Length & Average total turns per episode, including unresolved runs. & Aggregate efficiency. Drops toward \emph{Mean Turn} as \emph{Episode Resolve Rate} climbs, because fewer episodes are exhausting the turn cap without resolving. \\
Turns Saved & Average turns saved relative to a baseline episode length on resolved episodes. & Efficiency expressed in turn units. Couples resolution with promptness; nonzero values indicate the policy is not just resolving but resolving \emph{faster} than the baseline expects. \\
Efficiency Bonus & Additional reward for resolving episodes earlier than baseline. & Reward-side counterpart to \emph{Turns Saved}. Grows non-linearly with both arm match rate and resolve rate, since both must improve for early resolution to become routine. \\
\midrule
\multicolumn{3}{l}{\textit{Environment-correctness metrics --- is the policy seeing the reward signal we intended?}} \\
\midrule
Redundancy Rate & Fraction of turns where the oracle's response eliminates zero candidates. & Diagnostic of the filter rather than the policy. High values indicate that questions intended to be informative are being scored as wasted, gating the reward signal upstream of any learning. Near-zero values are a precondition for the mechanism metrics above to behave as designed. \\
Diversity Rate & Fraction of turns where the question was not asked previously in the same episode. & Guards against trivial repetition. Once near 1, it functions as a passive monitor: a drop is the first sign of mode collapse onto a small set of high-reward questions. \\
\midrule
% \multicolumn{3}{l}{\textit{Optimization-health metrics --- is training stable, and is the gain real or about to break?}} \\
% \midrule
% Entropy & Entropy of the policy's action distribution over questions and arms. & Specialization-vs.-collapse indicator. Read jointly with \emph{Diversity Rate} and \emph{Arm Match Rate}: falling entropy alongside stable diversity and rising arm match indicates healthy specialization; falling entropy alongside falling diversity or stalling arm match is the mode-collapse signature. \\
% Reward Std & Standard deviation of rewards across episodes in the batch. & Variance over outcomes. Sustained high values during training reflect ongoing exploration; collapse alongside falling entropy is a second mode-collapse signal. \\
% Policy Loss & GRPO objective at the current step. & Stable negative values indicate ongoing improvement; sudden zeroing or sign flips signal optimizer issues. Less informative on its own than as a co-monitor with gradient norm. \\
% Gradient Norm & $L_2$ norm of the policy gradient prior to clipping. & Update-magnitude diagnostic. Sustained values near the clip threshold indicate the optimizer is being throttled; sudden drops to near zero suggest the policy has entered a flat region of the loss landscape. \\
% Clip Ratio & Fraction of tokens whose probability ratio exceeds the PPO-style clip bound. & Trust-region activity. Near-zero values confirm updates lie comfortably within the clip range; rising values flag aggressive updates that would destabilize without clipping. \\
\bottomrule
\end{tabularx}
\caption{Metrics used to analyze policy performance during training, training dynamics over turns, and environment correctness in the Clue selector task. Metrics are grouped by diagnostic role: outcome metrics report task success; mechanism metrics decompose \emph{how} success is achieved (information gain per question vs.\ strategic arm selection vs.\ efficiency); environment-correctness metrics confirm the reward signal is reaching the policy as intended}

% ; and optimization-health metrics distinguish genuine specialization from impending mode collapse.
\label{tab:metric_definitions}
\end{table*}


\section{CSG Episode Dataset and Statistics}

\input{csg_dataformat}

\end{document}
